% Vectors — Pure Mathematics with Mechanics Upper Sixth — Cours signé OnBuch+
% Official MINESEC syllabus. Compiler avec : tectonic PMM03-vectors.tex
\newcommand{\DOCMATIERE}{Pure Mathematics with Mechanics}
\newcommand{\DOCNIVEAU}{Upper Sixth}
\newcommand{\DOCMODULE}{Upper Sixth — Pure mathematics with mechanics}
\newcommand{\DOCLECON}{3}
\newcommand{\DOCTITRE}{Vectors}
% OnBuch+ — préambule « cours signés » ENRICHI (compilé avec tectonic / XeLaTeX)
% Système visuel « Pop » d'OnBuch+ : fond crème (jamais blanc), bordures encre
% épaisses, ombres « dures » décalées, titres Archivo Black en capitales.
% Version MATHÉMATIQUES : graphes (pgfplots/tikz), boîtes pédagogiques
% (définition, propriété, méthode, attention, le savais-tu, exercice résolu,
% exercice, TP, bilan), page de garde et signature OnBuch+ ancrée en bas.
%
% Macros attendues dans main.tex AVANT \input{preamble} :
%   \DOCMATIERE  \DOCNIVEAU  \DOCMODULE  \DOCLECON (numéro)  \DOCTITRE
\documentclass[11pt]{article}

\usepackage{fontspec}
\usepackage[english]{babel}
\usepackage[a4paper,margin=2cm,top=3.4cm,bottom=2.8cm,headheight=1.4cm,footskip=1.3cm]{geometry}
\usepackage{amsmath,amssymb,amsfonts}
\usepackage{mathtools} % \xmapsto, \coloneqq... souvent utilisés par les modèles
\usepackage{cancel} % simplifications barrées (bilans de forces)
% \abs{z} (module, valeur absolue) et \norm{u} (norme), utilisés par les modèles
\providecommand{\abs}{}\let\abs\relax\DeclarePairedDelimiter{\abs}{\lvert}{\rvert}
\providecommand{\norm}{}\let\norm\relax\DeclarePairedDelimiter{\norm}{\lVert}{\rVert}
\usepackage[table]{xcolor} % [table] : fournit aussi \rowcolors (alternance de lignes)
\usepackage{pagecolor}
\usepackage{tikz}
\usetikzlibrary{arrows.meta,positioning,calc,shapes.geometric,shapes.misc,shapes.arrows,decorations.pathmorphing,decorations.pathreplacing,decorations.markings,patterns,shadows,mindmap,trees,fit,backgrounds,matrix,intersections,angles,quotes}
\usepackage{pgfplots}
\usepackage[version=4]{mhchem} % équations nucléaires \ce{^{235}_{92}U}
\usepackage[european,siunitx]{circuitikz} % schémas électriques
\pgfplotsset{compat=1.18}
\usepgfplotslibrary{fillbetween,groupplots} % aires entre courbes (calcul intégral)
\usepackage{enumitem}
\usepackage{fancyhdr}
% [most] charge toutes les bibliothèques tcolorbox (listings, minted, raster,
% documentation...) et gonfle inutilement la mémoire TeX sur un document long
% (~300 boîtes) : on ne charge que ce qui est réellement utilisé.
\usepackage[skins,breakable]{tcolorbox}
\usepackage{graphicx}
\usepackage{colortbl}
\usepackage{booktabs}
\usepackage{tabularx}
\usepackage{diagbox} % case d'en-tête diagonale des tableaux à double entrée
% Colonnes extensibles centrée / gauche / droite (C, L, R), souvent utilisées par les modèles
\newcolumntype{C}{>{\centering\arraybackslash}X}
\newcolumntype{L}{>{\raggedright\arraybackslash}X}
\newcolumntype{R}{>{\raggedleft\arraybackslash}X}
\usepackage{array}
\usepackage{multicol}
\usepackage{siunitx}
% parse-numbers=false : accepte aussi \SI{2,5 \times 10^{-3}}{...}
\sisetup{locale=UK,output-decimal-marker={.},group-minimum-digits=4,parse-numbers=false}
\DeclareSIUnit\ohm{\ensuremath{\symup{\Omega}}} % U+2126 absent de la police
\DeclareSIUnit\division{div} % réglages d'oscilloscope (ms/div, V/div)
\DeclareSIUnit\tr{tr} % tours (vitesse de rotation, ex. \SI{1500}{\tr\per\minute})
\usepackage{qrcode}
% \degree : commande fréquemment utilisée par les modèles pour un angle ou une
% température en dehors de \SI{}{} (non fournie par les paquets chargés).
\providecommand{\degree}{^{\circ}}
% \qed (fin de démonstration), utilisé par les modèles sans amsthm
\providecommand{\qed}{\hfill$\square$}
% \barre : double barre des valeurs interdites dans un tableau de variations (tkz-tab)
\providecommand{\barre}{\ensuremath{\|}}
% environnement proof (amsthm non chargé) utilisé par les modèles
\ifdefined\proof\else\newenvironment{proof}[1][Proof]{\par\noindent\textit{#1.}\ }{\qed\par}\fi
\usepackage{lastpage}
\usepackage{hyperref}
\hypersetup{hidelinks}

% ============================================================
% Palette « Pop » OnBuch+ — mêmes hex que la classe P (plus_ui.dart)
% ============================================================
\definecolor{poporange}{HTML}{F59321}
\definecolor{popdark}{HTML}{E07A0C}
\definecolor{popcream}{HTML}{FFF6E8}
\definecolor{popink}{HTML}{1C1714}
\definecolor{poppurple}{HTML}{7A5AE0}
\definecolor{popgreen}{HTML}{2E9E6A}
\definecolor{poppink}{HTML}{E0457B}
\definecolor{popblue}{HTML}{2F7DE1}
\definecolor{popgold}{HTML}{C98A12}
\definecolor{popmuted}{HTML}{8A7F76}
\definecolor{popline}{HTML}{EADFD2}
\definecolor{popgray}{HTML}{8A7F76}
\colorlet{popgrey}{popgray}
% Alias tolérants (le modèle nomme parfois les couleurs différemment)
\colorlet{popwhite}{white}
\colorlet{popblack}{popink}
\colorlet{popred}{poppink}
\colorlet{popyellow}{popgold}
% Teintes claires pour fonds de tableaux / zones
\colorlet{popblueL}{popblue!10!white}
\colorlet{poporangeL}{poporange!14!white}
\colorlet{popgreenL}{popgreen!12!white}
\colorlet{poppurpleL}{poppurple!12!white}
\colorlet{poppinkL}{poppink!12!white}
\colorlet{popgoldL}{popgold!14!white}

\pagecolor{popcream}

% ============================================================
% Typographie
% ============================================================
\defaultfontfeatures{Ligatures=TeX}
\setmainfont{PlusJakartaSans-Medium.ttf}[Path=../../../fonts/,BoldFont=PlusJakartaSans-Bold.ttf]
% Maths en sans-serif (Fira Math) avec chiffres et lettres droites de Plus
% Jakarta, pour que les formules se fondent dans le texte.
\usepackage[math-style=ISO,bold-style=ISO]{unicode-math}
\setmathfont{FiraMath-Regular.otf}[Path=../../../fonts/]
\setmathfont{PlusJakartaSans-Medium.ttf}[Path=../../../fonts/,range={up/num,up/{Latin,latin},\mathup}]
% (icomma retiré : décimales avec point en anglais)
% Caractères mathématiques Unicode absents des polices de texte (Plus Jakarta,
% Archivo Black) : rendus par leur équivalent mathématique, en texte comme en
% formule — sinon ils s'affichent en carré vide (ex. « Arithmétique dans ℤ »).
\usepackage{newunicodechar}
\newunicodechar{ℝ}{\ensuremath{\mathbb{R}}}
\newunicodechar{ℤ}{\ensuremath{\mathbb{Z}}}
\newunicodechar{ℂ}{\ensuremath{\mathbb{C}}}
\newunicodechar{ℕ}{\ensuremath{\mathbb{N}}}
\newunicodechar{ℚ}{\ensuremath{\mathbb{Q}}}
\newunicodechar{𝔻}{\ensuremath{\mathbb{D}}}
\newunicodechar{α}{\ensuremath{\alpha}}
\newunicodechar{β}{\ensuremath{\beta}}
\newunicodechar{γ}{\ensuremath{\gamma}}
\newunicodechar{δ}{\ensuremath{\delta}}
\newunicodechar{ε}{\ensuremath{\varepsilon}}
\newunicodechar{θ}{\ensuremath{\theta}}
\newunicodechar{λ}{\ensuremath{\lambda}}
\newunicodechar{μ}{\ensuremath{\mu}}
\newunicodechar{σ}{\ensuremath{\sigma}}
\newunicodechar{τ}{\ensuremath{\tau}}
\newunicodechar{φ}{\ensuremath{\varphi}}
\newunicodechar{ψ}{\ensuremath{\psi}}
\newunicodechar{ω}{\ensuremath{\omega}}
\newunicodechar{Σ}{\ensuremath{\Sigma}}
% majuscules produites par \MakeUppercase dans les titres (α-aminés -> Α)
\newunicodechar{Α}{\ensuremath{\alpha}}
\newunicodechar{Β}{\ensuremath{\beta}}
\newunicodechar{Γ}{\ensuremath{\Gamma}}
\newunicodechar{Δ}{\ensuremath{\Delta}}
\newunicodechar{Ε}{\ensuremath{\varepsilon}}
\newunicodechar{Θ}{\ensuremath{\Theta}}
\newunicodechar{Λ}{\ensuremath{\Lambda}}
\newunicodechar{Μ}{\ensuremath{\mu}}
\newunicodechar{Τ}{\ensuremath{\tau}}
\newunicodechar{Φ}{\ensuremath{\Phi}}
\newunicodechar{Ψ}{\ensuremath{\Psi}}
\newunicodechar{Ω}{\ensuremath{\Omega}}
\newunicodechar{↦}{\ensuremath{\mapsto}}
\newunicodechar{∈}{\ensuremath{\in}}
\newunicodechar{∉}{\ensuremath{\notin}}
\newunicodechar{∧}{\ensuremath{\wedge}}
\newunicodechar{⇔}{\ensuremath{\Leftrightarrow}}
\newunicodechar{⟺}{\ensuremath{\Longleftrightarrow}}
\newunicodechar{⇒}{\ensuremath{\Rightarrow}}
\newunicodechar{⟹}{\ensuremath{\Longrightarrow}}
\newunicodechar{∘}{\ensuremath{\circ}}
\newunicodechar{∪}{\ensuremath{\cup}}
\newunicodechar{∩}{\ensuremath{\cap}}
\newunicodechar{≡}{\ensuremath{\equiv}}
\newunicodechar{⊂}{\ensuremath{\subset}}
\newunicodechar{⊥}{\ensuremath{\perp}}
\newunicodechar{∃}{\ensuremath{\exists}}
\newunicodechar{∀}{\ensuremath{\forall}}
\newunicodechar{∛}{\ensuremath{\sqrt[3]{\,}}}
\newunicodechar{∜}{\ensuremath{\sqrt[4]{\,}}}
\newunicodechar{⃗}{\ensuremath{{}^{\to}}}
\newunicodechar{ⁿ}{\ifmmode^{n}\else\textsuperscript{\ensuremath{n}}\fi}
\newunicodechar{ˣ}{\ifmmode^{x}\else\textsuperscript{\ensuremath{x}}\fi}
\newunicodechar{ᵇ}{\ifmmode^{b}\else\textsuperscript{\ensuremath{b}}\fi}
\newunicodechar{ʳ}{\ifmmode^{r}\else\textsuperscript{\ensuremath{r}}\fi}
\newunicodechar{ᵘ}{\ifmmode^{u}\else\textsuperscript{\ensuremath{u}}\fi}
\newunicodechar{ʸ}{\ifmmode^{y}\else\textsuperscript{\ensuremath{y}}\fi}
\newunicodechar{ᵃ}{\ifmmode^{a}\else\textsuperscript{\ensuremath{a}}\fi}
\newunicodechar{ᵉ}{\ifmmode^{e}\else\textsuperscript{\ensuremath{e}}\fi}
\newunicodechar{ᵏ}{\ifmmode^{k}\else\textsuperscript{\ensuremath{k}}\fi}
\newunicodechar{ᵖ}{\ifmmode^{p}\else\textsuperscript{\ensuremath{p}}\fi}
\newunicodechar{ᴺ}{\ifmmode^{N}\else\textsuperscript{\ensuremath{N}}\fi}
\newunicodechar{ᶜ}{\ifmmode^{c}\else\textsuperscript{\ensuremath{c}}\fi}
\newunicodechar{ʰ}{\ifmmode^{h}\else\textsuperscript{\ensuremath{h}}\fi}
\newunicodechar{ᵠ}{\ifmmode^{q}\else\textsuperscript{\ensuremath{q}}\fi}
\newunicodechar{ₙ}{\ifmmode_{n}\else\textsubscript{\ensuremath{n}}\fi}
\newunicodechar{ₐ}{\ifmmode_{a}\else\textsubscript{\ensuremath{a}}\fi}
\newunicodechar{ᵢ}{\ifmmode_{i}\else\textsubscript{\ensuremath{i}}\fi}
\newunicodechar{ᵣ}{\ifmmode_{r}\else\textsubscript{\ensuremath{r}}\fi}
\newunicodechar{ₖ}{\ifmmode_{k}\else\textsubscript{\ensuremath{k}}\fi}
\newunicodechar{ᵦ}{\ifmmode_{\beta}\else\textsubscript{\ensuremath{\beta}}\fi}
\newunicodechar{ₓ}{\ifmmode_{x}\else\textsubscript{\ensuremath{x}}\fi}
\newfontfamily\popdisplay{ArchivoBlack-Regular.ttf}[Path=../../../fonts/]
\newfontfamily\poplogo{SpaceGrotesk-Bold.ttf}[Path=../../../fonts/]
\newfontfamily\popsemi{PlusJakartaSans-SemiBold.ttf}[Path=../../../fonts/]
\newfontfamily\popxbold{PlusJakartaSans-ExtraBold.ttf}[Path=../../../fonts/]
\newfontfamily\popxboldls{PlusJakartaSans-ExtraBold.ttf}[Path=../../../fonts/,LetterSpace=3.0]

\setlist[enumerate]{leftmargin=1.7em,itemsep=5pt,topsep=4pt}
\setlist[itemize]{leftmargin=1.7em,itemsep=4pt,topsep=4pt,label={\color{poporange}\textbullet}}
\setlength{\parindent}{0pt}
\setlength{\parskip}{5pt}
\renewcommand{\arraystretch}{1.25}

% pgfplots : style Pop par défaut
\pgfplotsset{
  every axis/.append style={axis line style={popink,line width=1pt},tick style={popink},
    grid style={popline,line width=0.5pt},label style={font=\small},tick label style={font=\footnotesize}},
  popcurve/.style={poporange,line width=1.8pt,no markers,smooth},
}
% Même style disponible dans un \draw TikZ simple (hors pgfplots)
\tikzset{popcurve/.style={poporange,line width=1.8pt}}
\tikzset{popgrid/.style={popline,very thin}}
\colorlet{popcurve}{poporange} % le nom du style est parfois utilisé comme couleur

% ============================================================
% Pill / squiggle
% ============================================================
\newcommand{\poppill}[3]{%
  \tikz[baseline=-0.55ex]\node[draw=popink,line width=1.2pt,rounded corners=2.6pt,%
    fill=#2,inner xsep=6.5pt,inner ysep=3pt]{%
    \popxboldls\fontsize{7}{7}\selectfont\color{#3}\MakeUppercase{#1}};%
}
\newcommand{\popsquiggle}[1]{%
  \tikz[baseline=-0.4ex]{
    \draw[#1,line width=2.6pt,line cap=round]
      plot [smooth] coordinates {(0,0.09) (0.34,0.01) (0.68,0.21) (1.02,0.01) (1.36,0.10)};
  }%
}
% Mot-clé surligné (marqueur orange sous le texte)
\newcommand{\cle}[1]{{\popxbold\color{popdark}#1}}

% ============================================================
% En-tête / pied
% ============================================================
\pagestyle{fancy}
\fancyhf{}
\renewcommand{\headrulewidth}{0pt}
\renewcommand{\footrulewidth}{0pt}
\lhead{\raisebox{-0.15ex}{\poplogo\fontsize{13}{13}\selectfont\color{popink}OnBuch\color{poporange}+}}
\rhead{\poppill{\DOCMATIERE\ \textbullet\ \DOCNIVEAU\ \textbullet\ Chapter \DOCLECON}{white}{popink}}
\lfoot{\popxbold\fontsize{7.6}{7.6}\selectfont\color{popmuted}\MakeUppercase{Signed course OnBuch+}\ \textbullet\ {\popsemi\DOCTITRE}}
\rfoot{\tikz[baseline=-0.6ex]\node[rounded corners=8.5pt,fill=popink,minimum height=17pt,minimum width=17pt,inner xsep=5pt,inner sep=0]{\popxbold\fontsize{8}{8}\selectfont\color{popcream}\thepage\,/\,\pageref*{LastPage}};}

% ============================================================
% Titres
% ============================================================
\newcommand{\chaptitre}[2]{%
  \begin{tcolorbox}[enhanced,colback=poporange,colframe=popink,boxrule=2.6pt,arc=11pt,
      drop shadow=popink,left=15pt,right=15pt,top=12pt,bottom=11pt,before skip=3pt,after skip=18pt]
    \poppill{#2}{popink}{popcream}\par\vspace{9pt}
    {\raggedright\hyphenpenalty=10000\exhyphenpenalty=10000\popdisplay\fontsize{22}{24}\selectfont\color{popcream}\MakeUppercase{#1}\par}
  \end{tcolorbox}
}
% Section numérotée : pastille orange avec le numéro + titre Archivo
\newcounter{popsec}
\newcommand{\coursec}[1]{%
  \stepcounter{popsec}\setcounter{popsub}{0}%
  \par\vspace{16pt}\noindent
  \begin{minipage}{\linewidth}
  \tikz[baseline=-0.7ex]\node[draw=popink,line width=1.6pt,fill=poporange,rounded corners=4pt,minimum size=22pt,inner sep=0]{\popdisplay\fontsize{12}{12}\selectfont\color{popcream}\thepopsec};%
  \hspace{8pt}{\popdisplay\fontsize{15}{17}\selectfont\color{popink}\MakeUppercase{#1}}\par
  \vspace{4pt}\noindent{\color{poporange}\rule{36pt}{3.6pt}}
  \end{minipage}\par\nopagebreak\vspace{8pt}
}
\newcounter{popsub}
\newcommand{\courssub}[1]{%
  \stepcounter{popsub}%
  \par\vspace{9pt}\noindent{\popxbold\fontsize{11.5}{13}\selectfont\color{popink}{\color{poporange}\thepopsec.\thepopsub}\hspace{6pt}#1}\par\nopagebreak\vspace{4pt}
}

% ============================================================
% Boîtes pédagogiques — même recette : fond blanc, bordure épaisse colorée,
% ombre dure encre, badge plein. Titre optionnel [#1].
% ============================================================
\tcbset{popbase/.style={breakable,enhanced,colback=white,boxrule=2.3pt,arc=9pt,
  shadow={3pt}{-3pt}{0pt}{popink},left=14pt,right=14pt,top=11pt,bottom=11pt,before skip=11pt,after skip=15pt,
  fontupper=\popsemi\fontsize{10.6}{14.5}\selectfont\color{popink},
  fontlower=\popsemi\fontsize{10.6}{14.5}\selectfont\color{popink}}}
% Coche dessinée (le glyphe ✓ n'existe pas dans les polices du document)
\newcommand{\popcheck}[1]{\tikz[baseline=-0.5ex]\draw[#1,line width=1.6pt,line cap=round,line join=round](-0.13,0.0)--(-0.04,-0.09)--(0.13,0.1);}
\providecommand{\coche}{\popcheck{popgreen}}
\providecommand{\resultat}[1]{\par\smallskip\noindent\fcolorbox{popgreen}{popgreenL}{\parbox{\dimexpr\linewidth-2\fboxsep-2\fboxrule\relax}{\textbf{Result.} #1}}\par\smallskip}
\newcommand{\popboxhead}[3]{\poppill{#1}{#2}{white}\ifx\relax#3\relax\else\hspace{6pt}{\popxbold\fontsize{10.6}{12}\selectfont\color{popink}#3}\fi\par\vspace{7pt}}

\newtcolorbox{aretenir}[1][]{popbase,colframe=popgreen,before upper={\popboxhead{Key points}{popgreen}{#1}}}
\newtcolorbox{exemplebox}[1][]{popbase,colframe=poporange,before upper={\popboxhead{Exemple}{poporange}{#1}}}
\newtcolorbox{definition}[1][]{popbase,colframe=popblue,before upper={\popboxhead{Definition}{popblue}{#1}}}
\newtcolorbox{propriete}[1][]{popbase,colframe=poppurple,before upper={\popboxhead{Property}{poppurple}{#1}}}
\newtcolorbox{methode}[1][]{popbase,colframe=popgold,before upper={\popboxhead{Method}{popgold}{#1}}}
\newtcolorbox{attention}[1][]{popbase,colframe=poppink,before upper={\popboxhead{Warning}{poppink}{#1}}}
\newtcolorbox{experience}[1][]{popbase,colframe=popblue,colback=popblueL,before upper={\popboxhead{Experiment}{popblue}{#1}}}
% « Le savais-tu ? » : carte violette pleine (contexte, culture, Cameroun)
\newtcolorbox{savaistu}[1][]{popbase,colback=poppurple,colframe=popink,
  fontupper=\popsemi\fontsize{10.4}{14.2}\selectfont\color{white},
  before upper={\poppill{Did you know?}{white}{poppurple}\ifx\relax#1\relax\else\hspace{6pt}{\popxbold\color{white}#1}\fi\par\vspace{7pt}}}
% Exercice résolu : énoncé en haut, \tcblower puis la solution sur fond crème
\newcounter{popexo}\newcounter{popexr}
\newtcolorbox{exoresolu}[1][]{popbase,colframe=poporange,colbacklower=popcream,
  segmentation style={popink,line width=1.2pt,dashed},
  before upper={\stepcounter{popexr}\popboxhead{Worked example \thepopexr}{poporange}{#1}},
  before lower={\poppill{Solution}{popink}{popcream}\par\vspace{6pt}}}
% Exercice (non résolu en place) — la correction est dans la partie Corrigés
\newtcolorbox{exercice}[1][]{popbase,colframe=popink,
  before upper={\stepcounter{popexo}\popboxhead{Exercise \thepopexo}{popink}{#1}}}
% Corrigé
\newtcolorbox{corrige}[1][]{popbase,colframe=popgreen,colback=popgreenL,
  before upper={\popboxhead{Solution}{popgreen}{#1}}}
% Objectifs / prérequis / bilan
\newtcolorbox{objectifs}{popbase,colframe=popink,colback=poporangeL,
  before upper={\popboxhead{Chapter objectives}{popink}{}}}
\newtcolorbox{prerequis}{popbase,colframe=popink,colback=popblueL,
  before upper={\popboxhead{Prerequisites}{popblue}{}}}
\newtcolorbox{bilan}{popbase,colframe=popink,colback=white,boxrule=2.8pt,
  before upper={\popboxhead{Fiche bilan}{poporange}{L'essentiel à connaître pour l'examen}}}
% Figure encadrée avec légende
\newtcolorbox{popfigure}[1][]{enhanced,breakable=false,colback=white,colframe=popline,boxrule=1.4pt,arc=8pt,
  left=10pt,right=10pt,top=10pt,bottom=8pt,before skip=10pt,after skip=12pt,halign=center,
  lower separated=false,halign lower=center,
  fontlower=\popsemi\fontsize{9}{11.5}\selectfont\color{popmuted},
  #1}
\newcommand{\legende}[1]{\par\vspace{6pt}{\popsemi\fontsize{9}{11.5}\selectfont\color{popmuted}#1}}

% ============================================================
% Signature OnBuch+
% ============================================================
\newcommand{\poplogoinline}[1]{{\poplogo\fontsize{#1}{#1}\selectfont\color{popink}OnBuch\color{poporange}+}}
\newcommand{\signatureonbuch}{%
  \begin{tcolorbox}[enhanced,colback=popink,colframe=popink,boxrule=2.2pt,arc=10pt,
      drop shadow=poporange,left=14pt,right=14pt,top=10pt,bottom=10pt,before skip=0pt,after skip=0pt]
    \tikz[baseline=-0.7ex]\node[circle,fill=popgreen,draw=popcream,line width=1.3pt,minimum size=22pt,inner sep=0]{\popcheck{white}};%
    \hspace{9pt}\begin{minipage}[c]{0.62\linewidth}
      {\popxbold\fontsize{12}{13}\selectfont\color{popcream}Signed\ \textbullet\ {\poplogo OnBuch\color{poporange}+}}\par\vspace{2pt}
      {\raggedright\popsemi\fontsize{8}{10.5}\selectfont\color{popline}\MakeUppercase{OnBuch+ teaching team}\\\MakeUppercase{Follows the official MINESEC syllabus}\par}
    \end{minipage}\hfill
    {\poplogo\fontsize{22}{22}\selectfont\color{popcream}OnBuch\color{poporange}+}
  \end{tcolorbox}%
}

% ============================================================
% Page de garde
% #1 titre · #2 matière · #3 niveau · #4 module · #5 n° leçon · #6 durée ·
% #7 description / objectifs (texte ou itemize)
% ============================================================
\newcommand{\pagegarde}[7]{%
  \thispagestyle{empty}
  \newgeometry{a4paper,margin=1.7cm,top=1.6cm,bottom=1.4cm}%
  \begin{tikzpicture}[remember picture,overlay]
    \fill[poporange!18] ([xshift=-3.2cm,yshift=-3.6cm]current page.north east) circle (4.6cm);
    \fill[poppurple!14] ([xshift=2.4cm,yshift=7.6cm]current page.south west) circle (3.8cm);
  \end{tikzpicture}%
  {\poplogo\fontsize{30}{30}\selectfont\color{popink}OnBuch\color{poporange}+}\hfill
  \raisebox{0.6ex}{\poppill{Signed course \textbullet\ Premium}{popink}{popcream}}\par
  \vspace{1pt}\popsquiggle{poppurple}\par
  \vspace{8pt}
  \begin{tcolorbox}[enhanced,colback=poporange,colframe=popink,boxrule=2.8pt,arc=13pt,
      drop shadow=popink,left=18pt,right=18pt,top=12pt,bottom=13pt,after skip=10pt]
    \poppill{#2 \textbullet\ #3}{popink}{popcream}\hspace{5pt}\poppill{#4}{white}{popink}\par\vspace{8pt}
    {\popdisplay\fontsize{14}{14}\selectfont\color{popink}CHAPTER #5}\par\vspace{4pt}
    {\raggedright\hyphenpenalty=10000\exhyphenpenalty=10000\popdisplay\fontsize{25}{27}\selectfont\color{popcream}\MakeUppercase{#1}\par}
    \vspace{8pt}
    {\popxbold\fontsize{9.5}{9.5}\selectfont\color{popink}\MakeUppercase{Suggested time: #6}}
  \end{tcolorbox}
  \begin{tcolorbox}[enhanced,colback=white,colframe=popink,boxrule=2.2pt,arc=10pt,
      drop shadow=popink,left=16pt,right=16pt,top=10pt,bottom=10pt,after skip=8pt]
    \poppill{In this chapter}{poppurple}{white}\par\vspace{5pt}
    \popsemi\fontsize{9.3}{12.6}\selectfont\color{popink}#7
  \end{tcolorbox}
  \noindent\begin{minipage}[t]{0.487\linewidth}
    \begin{tcolorbox}[enhanced,colback=popgreen,colframe=popink,boxrule=2.2pt,arc=10pt,
        drop shadow=popink,left=11pt,right=11pt,top=10pt,bottom=10pt,height=3.1cm,valign=center]
      \tcbox[colback=white,colframe=popink,boxrule=1.3pt,arc=5pt,boxsep=0pt,nobeforeafter,
        left=3pt,right=3pt,top=3pt,bottom=3pt,valign=center]{\qrcode[height=1.9cm]{https://play.google.com/store/apps/details?id=com.luvvix.onbuch&hl=en}}%
      \hspace{8pt}\begin{minipage}[c]{0.5\linewidth}
        {\popdisplay\fontsize{11}{12.5}\selectfont\color{white}\MakeUppercase{Download OnBuch}}\par\vspace{4pt}
        {\raggedright\popsemi\fontsize{8.4}{11}\selectfont\color{white}Free \textbullet\ Google Play\\AI tutor, past papers, results\par}
      \end{minipage}
    \end{tcolorbox}
  \end{minipage}\hfill
  \begin{minipage}[t]{0.487\linewidth}
    \begin{tcolorbox}[enhanced,colback=poppurple,colframe=popink,boxrule=2.2pt,arc=10pt,
        drop shadow=popink,left=11pt,right=11pt,top=10pt,bottom=10pt,height=3.1cm,valign=center]
      \tikz[baseline=-0.7ex]\node[circle,fill=white,draw=popink,line width=1.3pt,minimum size=1.6cm,inner sep=0]{\popdisplay\fontsize{24}{24}\selectfont\color{poppurple}+};%
      \hspace{8pt}\begin{minipage}[c]{0.6\linewidth}
        {\popdisplay\fontsize{11}{12.5}\selectfont\color{white}\MakeUppercase{Go OnBuch+}}\par\vspace{4pt}
        {\raggedright\popsemi\fontsize{8.4}{11}\selectfont\color{white}All signed courses, unlimited AI tutor, PDF sheets — OnBuch+ tab in the app\par}
      \end{minipage}
    \end{tcolorbox}
  \end{minipage}
  \vspace{8pt}
  \popcontenu
  \vfill
  \signatureonbuch
  \restoregeometry
  \newpage
}

% Rangée « Ce cours contient » (page de garde)
\newcommand{\poptile}[3]{% #1 couleur #2 gros texte #3 légende
  \begin{tcolorbox}[enhanced,colback=#1,colframe=popink,boxrule=2pt,arc=9pt,shadow={2.5pt}{-2.5pt}{0pt}{popink},
      left=6pt,right=6pt,top=9pt,bottom=9pt,halign=center,valign=center,height=1.75cm,width=\linewidth,nobeforeafter]
    {\popdisplay\fontsize{12.5}{13}\selectfont\color{white}\MakeUppercase{#2}}\par\vspace{3pt}
    {\centering\popsemi\fontsize{7.8}{9.5}\selectfont\color{white}#3\par}
  \end{tcolorbox}}
\newcommand{\popcontenu}{%
  \par\noindent{\popxboldls\fontsize{7.5}{7.5}\selectfont\color{popmuted}THIS SIGNED COURSE CONTAINS}\par\vspace{7pt}
  \noindent\begin{minipage}[t]{0.235\linewidth}\poptile{popblue}{Course}{complete, illustrated, step by step}\end{minipage}\hfill
  \begin{minipage}[t]{0.235\linewidth}\poptile{poporange}{Methods}{and worked examples}\end{minipage}\hfill
  \begin{minipage}[t]{0.235\linewidth}\poptile{poppink}{Exercises}{exam style + solutions}\end{minipage}\hfill
  \begin{minipage}[t]{0.235\linewidth}\poptile{popgreen}{Summary sheet}{the essentials to remember}\end{minipage}\par}

% Sommaire
\newtcolorbox{sommaire}{enhanced,colback=white,colframe=popink,boxrule=2.2pt,arc=10pt,
  drop shadow=popink,left=18pt,right=18pt,top=13pt,bottom=13pt,before skip=6pt,after skip=20pt,
  before upper={\poppill{Contents}{poppurple}{white}\par\vspace{9pt}\popsemi\fontsize{10.6}{14.5}\selectfont\color{popink}}}

% Signature de fin de cours — poussée tout en bas de la dernière page
\newcommand{\findecours}{%
  \par\vfill\nopagebreak
  \begin{center}\popsquiggle{poporange}\end{center}\vspace{4pt}
  \signatureonbuch
}
\AtBeginDocument{\def\llbracket{\mathopen{[\mkern-3.5mu[}}\def\rrbracket{\mathclose{]\mkern-3.5mu]}}} % intervalles d'entiers (glyphe ⟦ absent de la police)
% Glyphes absents de Fira Math : redéfinis à partir de glyphes disponibles
\AtBeginDocument{%
  \def\setminus{\mathbin{\backslash}}\let\smallsetminus\setminus
  \def\vdots{\mathord{\vbox{\baselineskip3.2pt\lineskiplimit0pt\kern4pt\hbox{.}\hbox{.}\hbox{.}}}}%
  \def\ddots{\mathinner{\mkern1mu\raise6.5pt\hbox{.}\mkern2mu\raise3.6pt\hbox{.}\mkern2mu\raise0.7pt\hbox{.}\mkern1mu}}%
  \def\triangle{\mathord{\scalebox{0.9}{$\symup{\Delta}$}}}%
  \def\nabla{\mathord{\raisebox{\depth}{\scalebox{1}[-1]{$\symup{\Delta}$}}}}%
  \def\checkmark{\popcheck{popdark}}%
  \def\star{\mathbin{\ast}}\def\bigstar{\mathord{\ast}}%
  \def\diamond{\mathbin{\scalebox{0.6}{$\Diamond$}}}%
  \def\bigcup{\mathop{\mathchoice{\raisebox{-0.3em}{\scalebox{1.7}{$\cup$}}}{\scalebox{1.25}{$\cup$}}{\cup}{\cup}}\displaylimits}%
  \def\bigcap{\mathop{\mathchoice{\raisebox{-0.3em}{\scalebox{1.7}{$\cap$}}}{\scalebox{1.25}{$\cap$}}{\cap}{\cap}}\displaylimits}%
  \def\lesseqqgtr{\mathrel{\lessgtr}}\def\bigoplus{\mathop{\oplus}\displaylimits}%
}
\providecommand{\ssi}{if and only if} % abréviation fréquente
\AtBeginDocument{\providecommand{\wideparen}{\overparen}} % arc de cercle

\begin{document}

\pagegarde{Vectors}{Pure Mathematics with Mechanics}{Upper Sixth}{Upper Sixth — Pure mathematics with mechanics}{3}{11 periods}{This chapter extends vector methods to the geometry of planes in three-dimensional space. Learners will derive vector, parametric and Cartesian equations of planes, compute distances and angles, and study the intersection of lines and planes — essential tools for the GCE Advanced Level examination.
\vspace{4pt}
\begin{multicols}{2}\small
\begin{itemize}
\item By the end of this chapter I can write the vector equation of a plane given a point and a normal vector, or a point and two direction vectors.
\item By the end of this chapter I can convert between the vector, parametric and Cartesian forms of the equation of a plane.
\item By the end of this chapter I can find the Cartesian equation of a plane through three non-collinear points.
\item By the end of this chapter I can calculate the perpendicular distance from a point to a plane.
\item By the end of this chapter I can find the angle between two planes and the angle between a line and a plane.
\item By the end of this chapter I can determine the point of intersection of a line and a plane and discuss existence and uniqueness.
\end{itemize}\end{multicols}}

\begin{sommaire}
\begin{multicols}{2}
\begin{enumerate}
\item Equations of a Plane: Vector, Parametric and Cartesian Forms
\item Plane Through Three Points and Special Planes
\item Perpendicular Distance from a Point to a Plane
\item Angles: Between Two Planes and Between a Line and a Plane
\item Intersection of a Line with a Plane and of Two Planes
\item Pencils of Planes and Synthesis Problems
\item Integration activity
\item Methods and worked examples
\item Exercises
\item Solutions to the exercises
\item Summary sheet
\end{enumerate}
\end{multicols}
\end{sommaire}

\begin{objectifs}
\begin{itemize}
\item By the end of this chapter I can write the vector equation of a plane given a point and a normal vector, or a point and two direction vectors.
\item By the end of this chapter I can convert between the vector, parametric and Cartesian forms of the equation of a plane.
\item By the end of this chapter I can find the Cartesian equation of a plane through three non-collinear points.
\item By the end of this chapter I can calculate the perpendicular distance from a point to a plane.
\item By the end of this chapter I can find the angle between two planes and the angle between a line and a plane.
\item By the end of this chapter I can determine the point of intersection of a line and a plane and discuss existence and uniqueness.
\item By the end of this chapter I can find the line of intersection of two planes and recognise parallel or coincident planes.
\item By the end of this chapter I can write the equation of a plane through the line of intersection of two given planes (pencil of planes).
\item By the end of this chapter I can solve GCE-style problems combining several of these skills.
\end{itemize}
\end{objectifs}

\begin{prerequis}
\begin{itemize}
\item Scalar (dot) product of two vectors and its link with the angle between them (Lower Sixth).
\item Vector and parametric equations of a straight line in 3D (Lower Sixth).
\item Vector (cross) product and its geometric meaning as area and perpendicular direction.
\item Position vectors, magnitude of a vector, and basic vector algebra.
\item Solving simultaneous linear equations in two and three unknowns.
\end{itemize}
\end{prerequis}

\newpage

\coursec{Equations of a Plane: Vector, Parametric and Cartesian Forms}

When engineers designed the roof of the Ahmadou Ahidjo Stadium in Yaoundé, every steel panel had to be positioned in space with millimetre precision: each panel is a piece of a plane, and the angle between two panels determines how rainwater flows off. Similarly, when an aircraft lands at Douala International Airport, the pilot must control the angle between the aircraft's path (a line) and the runway (a plane). How can we describe a flat surface in space with an equation, measure distances from it, and find where lines and planes meet? Vectors give us a powerful and elegant language to answer all these questions.

\courssub{The Concept of a Plane in Space}

A plane in three-dimensional space is a flat, two-dimensional surface that extends infinitely in all directions. Unlike a line, which is one-dimensional and requires only one direction vector, a plane is \textbf{two-dimensional}. To fix a plane in space we need either:
\begin{itemize}
    \item A \textbf{point} on the plane and a \textbf{normal vector} (a vector perpendicular to the plane), \emph{or}
    \item A \textbf{point} on the plane and \textbf{two non-parallel direction vectors} lying in the plane.
\end{itemize}

\begin{popfigure}
\begin{tikzpicture}[scale=1.2, >=Stealth]
% Axes
\draw[->, thick] (0,0,0) -- (3,0,0) node[anchor=north west] {$x$};
\draw[->, thick] (0,0,0) -- (0,3,0) node[anchor=south east] {$y$};
\draw[->, thick] (0,0,0) -- (0,0,3) node[anchor=south] {$z$};

% Plane: x + y + z = 2 (intercepts 2,2,2)
\fill[popblueL, opacity=0.3] (2,0,0) -- (0,2,0) -- (0,0,2) -- cycle;
\draw[popblue, thick] (2,0,0) -- (0,2,0) -- (0,0,2) -- cycle;

% Point A on plane
\coordinate (A) at (1,0.5,0.5);
\fill[popink] (A) circle (2pt) node[anchor=south west] {$A(\vec{a})$};

% Normal vector n
\draw[poporange, very thick, ->] (A) -- ++(0.6,0.6,0.6) node[anchor=south west] {$\vec{n}$};

% Direction vectors u, v in plane
\draw[popgreen, very thick, ->] (A) -- ++(-0.5,0.5,0) node[anchor=north] {$\vec{u}$};
\draw[poppurple, very thick, ->] (A) -- ++(0,-0.5,0.5) node[anchor=east] {$\vec{v}$};

% Dashed lines to axes for intercepts
\draw[dashed, thin, gray] (2,0,0) -- (0,0,0);
\draw[dashed, thin, gray] (0,2,0) -- (0,0,0);
\draw[dashed, thin, gray] (0,0,2) -- (0,0,0);
\node[anchor=north east] at (2,0,0) {$\alpha$};
\node[anchor=north west] at (0,2,0) {$\beta$};
\node[anchor=south] at (0,0,2) {$\gamma$};
\end{tikzpicture}
\legende{A plane defined by point $A$, normal vector $\vec{n}$, and two direction vectors $\vec{u}, \vec{v}$ lying in the plane. The intercepts on the axes are $\alpha, \beta, \gamma$.}
\end{popfigure}

\begin{definition}[Normal Vector to a Plane]
A vector $\vec{n}$ is a \textbf{normal vector} to a plane $\Pi$ if $\vec{n}$ is perpendicular to \emph{every} vector lying in $\Pi$. If $\vec{n}$ is a normal vector, then any non-zero scalar multiple $k\vec{n}$ ($k \neq 0$) is also a normal vector.
\end{definition}

\begin{definition}[Vector Equation of a Plane (Normal Form)]
Let $\Pi$ be a plane with normal vector $\vec{n}$. Let $A$ be a fixed point on $\Pi$ with position vector $\vec{a}$, and let $P$ be any point on $\Pi$ with position vector $\vec{r}$. Then the vector $\overrightarrow{AP} = \vec{r} - \vec{a}$ lies in the plane, so it is perpendicular to $\vec{n}$.
\[
(\vec{r} - \vec{a}) \cdot \vec{n} = 0 \quad \Longrightarrow \quad \vec{r} \cdot \vec{n} = \vec{a} \cdot \vec{n}.
\]
The quantity $\vec{a} \cdot \vec{n}$ is a constant, denoted by $p$ (or $d$). The equation is
\[
\[\boxed{\vec{r} \cdot \vec{n} = p}\]
\]
where $\vec{n}$ is a normal vector and $p$ is the perpendicular distance from the origin to the plane multiplied by $|\vec{n}|$.
\end{definition}

\begin{aretenir}[Key Points: Vector Equation $\vec{r} \cdot \vec{n} = p$]
\begin{itemize}
\item $\vec{n} = \begin{pmatrix} a \\ b \\ c \end{pmatrix}$ is a normal vector to the plane.
\item $p = \vec{a} \cdot \vec{n}$ is a scalar constant.
\item If $\vec{r} = \begin{pmatrix} x \\ y \\ z \end{pmatrix}$, the equation expands to the \textbf{Cartesian form}: $ax + by + cz = p$.
\item The coefficients of $x, y, z$ in the Cartesian equation give the components of a normal vector directly.
\end{itemize}
\end{aretenir}

\begin{exoresolu}[Finding the Vector Equation Given a Point and a Normal]
Find the vector equation of the plane $\Pi$ that passes through the point $A(2, -1, 3)$ and has normal vector $\vec{n} = 2\vec{i} - 3\vec{j} + 4\vec{k}$.
\tcblower
\textbf{Solution:}
The position vector of $A$ is $\vec{a} = \begin{pmatrix} 2 \\ -1 \\ 3 \end{pmatrix}$.
The normal vector is $\vec{n} = \begin{pmatrix} 2 \\ -3 \\ 4 \end{pmatrix}$.
The constant $p = \vec{a} \cdot \vec{n} = (2)(2) + (-1)(-3) + (3)(4) = 4 + 3 + 12 = 19$.
The vector equation is:
\[
\vec{r} \cdot \begin{pmatrix} 2 \\ -3 \\ 4 \end{pmatrix} = 19
\]
In Cartesian form: $2x - 3y + 4z = 19$.
\end{exoresolu}

\courssub{Parametric Vector Form of a Plane}

If we know a point $A(\vec{a})$ on the plane and two non-parallel vectors $\vec{u}$ and $\vec{v}$ parallel to the plane (direction vectors), then any point $P(\vec{r})$ on the plane can be reached by starting at $A$ and moving some distance along $\vec{u}$ and some distance along $\vec{v}$.

\begin{definition}[Parametric (Two-Parameter) Vector Equation]
Let $\vec{a}$ be the position vector of a point $A$ on the plane $\Pi$, and let $\vec{u}, \vec{v}$ be two \textbf{non-parallel} vectors parallel to $\Pi$. Then the position vector $\vec{r}$ of any point $P$ on $\Pi$ is given by:
\[
\[\boxed{\vec{r} = \vec{a} + \lambda \vec{u} + \mu \vec{v} \qquad (\lambda, \mu \in \mathbb{R})}\]
\]
The parameters $\lambda$ and $\mu$ are real numbers. Varying $\lambda$ and $\mu$ independently generates the whole plane.
\end{definition}

\begin{attention}[Vectors Must Not Be Parallel]
In the equation $\vec{r} = \vec{a} + \lambda \vec{u} + \mu \vec{v}$, the vectors $\vec{u}$ and $\vec{v}$ \textbf{must not be parallel} (i.e., $\vec{u} \neq k\vec{v}$ for any scalar $k$). If they are parallel, the equation describes only a \textbf{line}, not a plane. Always check that $\vec{u} \times \vec{v} \neq \vec{0}$.
\end{attention}

\courssub{Converting Parametric Form to Cartesian Form}

Given $\vec{r} = \vec{a} + \lambda \vec{u} + \mu \vec{v}$, we write the components:
\[
\begin{cases}
x = a_1 + \lambda u_1 + \mu v_1 \\
y = a_2 + \lambda u_2 + \mu v_2 \\
z = a_3 + \lambda u_3 + \mu v_3
\end{cases}
\]
We eliminate the parameters $\lambda$ and $\mu$ to obtain a single equation in $x, y, z$ of the form $ax + by + cz = d$.

\begin{methode}[Converting Parametric $\to$ Cartesian]
\begin{enumerate}
    \item Write the parametric equations for $x, y, z$ in terms of $\lambda, \mu$.
    \item Treat two of the equations as a linear system in $\lambda, \mu$ (e.g., solve for $\lambda, \mu$ in terms of $x, y$).
    \item Substitute the expressions for $\lambda, \mu$ into the third equation.
    \item Simplify to get the Cartesian form $ax + by + cz = d$.
    \item \textbf{Alternative (faster for exams):} Compute the normal vector $\vec{n} = \vec{u} \times \vec{v}$. Then use $\vec{r} \cdot \vec{n} = \vec{a} \cdot \vec{n}$.
\end{enumerate}
\end{methode}

\begin{exoresolu}[Parametric to Cartesian via Cross Product]
A plane $\Pi$ has parametric equation:
\[
\vec{r} = \begin{pmatrix} 1 \\ 2 \\ -1 \end{pmatrix} + \lambda \begin{pmatrix} 2 \\ 1 \\ 0 \end{pmatrix} + \mu \begin{pmatrix} -1 \\ 0 \\ 3 \end{pmatrix}.
\]
Find its Cartesian equation.
\tcblower
\textbf{Method 1: Cross Product (Recommended)}
Direction vectors: $\vec{u} = \begin{pmatrix} 2 \\ 1 \\ 0 \end{pmatrix}$, $\vec{v} = \begin{pmatrix} -1 \\ 0 \\ 3 \end{pmatrix}$.
Normal vector $\vec{n} = \vec{u} \times \vec{v} = \begin{vmatrix} \vec{i} & \vec{j} & \vec{k} \\ 2 & 1 & 0 \\ -1 & 0 & 3 \end{vmatrix}
= \vec{i}(3 - 0) - \vec{j}(6 - 0) + \vec{k}(0 - (-1)) = \begin{pmatrix} 3 \\ -6 \\ 1 \end{pmatrix}$.
Point $A(1, 2, -1)$, so $\vec{a} = \begin{pmatrix} 1 \\ 2 \\ -1 \end{pmatrix}$.
Constant $p = \vec{a} \cdot \vec{n} = 3(1) - 6(2) + 1(-1) = 3 - 12 - 1 = -10$.
Cartesian equation: $3x - 6y + z = -10$ or $\boxed{3x - 6y + z + 10 = 0}$.

\textbf{Method 2: Elimination of Parameters}
$x = 1 + 2\lambda - \mu$ \quad (1) \\
$y = 2 + \lambda$ \quad $\Rightarrow \lambda = y - 2$ \quad (2) \\
$z = -1 + 3\mu$ \quad $\Rightarrow \mu = \frac{z+1}{3}$ \quad (3)
Substitute (2) and (3) into (1):
$x = 1 + 2(y-2) - \frac{z+1}{3}$
$3x = 3 + 6y - 12 - z - 1$
$3x - 6y + z = -10$. \quad \checkmark
\end{exoresolu}

\courssub{Cartesian Equation and the Normal Vector}

The general Cartesian equation of a plane is
\[
ax + by + cz = d \quad \text{or} \quad ax + by + cz + d = 0 \quad (d \text{ sign convention varies}).
\]

\begin{propriete}[Normal Vector from Cartesian Equation]
For the plane $ax + by + cz = d$, a normal vector is
\[
\vec{n} = \begin{pmatrix} a \\ b \\ c \end{pmatrix} = a\vec{i} + b\vec{j} + c\vec{k}.
\]
\end{propriete}

\begin{attention}[Coefficients Give the Normal, Not a Direction]
A very common mistake is to think that $\begin{pmatrix} a \\ b \\ c \end{pmatrix}$ is a vector \emph{in} the plane. It is \textbf{perpendicular} to the plane. Direction vectors in the plane satisfy $a u_1 + b u_2 + c u_3 = 0$ (they are orthogonal to $\vec{n}$).
\end{attention}

\begin{exoresolu}[Reading the Normal and Finding a Point]
For the plane $\Pi: 2x - y + 3z = 6$:
\begin{enumerate}
    \item State a normal vector.
    \item Find the coordinates of three points on $\Pi$.
\end{enumerate}
\tcblower
\textbf{Solution:}
\begin{enumerate}
    \item Normal vector $\vec{n} = \begin{pmatrix} 2 \\ -1 \\ 3 \end{pmatrix}$.
    \item Choose arbitrary values for two variables and solve for the third.
    \begin{itemize}
        \item Let $x=0, y=0 \Rightarrow 3z=6 \Rightarrow z=2$. Point $A(0,0,2)$.
        \item Let $x=0, z=0 \Rightarrow -y=6 \Rightarrow y=-6$. Point $B(0,-6,0)$.
        \item Let $y=0, z=0 \Rightarrow 2x=6 \Rightarrow x=3$. Point $C(3,0,0)$.
    \end{itemize}
    These are exactly the intercepts on the axes.
\end{enumerate}
\end{exoresolu}

\courssub{Intercepts and the Intercept Form}

If a plane cuts the $x$-, $y$-, and $z$-axes at $(\alpha,0,0)$, $(0,\beta,0)$, and $(0,0,\gamma)$ respectively, then $\alpha, \beta, \gamma$ are the \textbf{intercepts}. The plane does not pass through the origin, so $\alpha, \beta, \gamma \neq 0$.

\begin{popfigure}
\begin{tikzpicture}[scale=1.3, >=Stealth]
% Axes
\draw[->, thick] (0,0,0) -- (3.5,0,0) node[anchor=north west] {$x$};
\draw[->, thick] (0,0,0) -- (0,3.5,0) node[anchor=south east] {$y$};
\draw[->, thick] (0,0,0) -- (0,0,3.5) node[anchor=south] {$z$};

% Plane intercepts: alpha=3, beta=2, gamma=2.5
\coordinate (X) at (3,0,0);
\coordinate (Y) at (0,2,0);
\coordinate (Z) at (0,0,2.5);

% Plane surface
\fill[popgreenL, opacity=0.4] (X) -- (Y) -- (Z) -- cycle;
\draw[popgreen, thick] (X) -- (Y) -- (Z) -- cycle;

% Labels
\node[anchor=north east] at (X) {$\alpha$};
\node[anchor=north west] at (Y) {$\beta$};
\node[anchor=south] at (Z) {$\gamma$};
\draw[dashed, thin] (X) -- (0,0,0) -- (Y) (0,0,0) -- (Z);
\end{tikzpicture}
\legende{A plane cutting the coordinate axes at intercepts $\alpha, \beta, \gamma$.}
\end{popfigure}

Substituting these points into $ax+by+cz=d$ gives $a\alpha=d$, $b\beta=d$, $c\gamma=d$. Thus $a=d/\alpha$, $b=d/\beta$, $c=d/\gamma$. The equation becomes:
\[
\frac{d}{\alpha}x + \frac{d}{\beta}y + \frac{d}{\gamma}z = d \quad \Longrightarrow \quad \boxed{\frac{x}{\alpha} + \frac{y}{\beta} + \frac{z}{\gamma} = 1}
\]
This is the \textbf{intercept form}.

\begin{aretenir}[GCE Exam Tip: Intercept Form]
The intercept form $\frac{x}{\alpha} + \frac{y}{\beta} + \frac{z}{\gamma} = 1$ is extremely useful in \textbf{multiple-choice questions} for quickly identifying the intercepts or checking if a plane passes through a given point on an axis. If a plane is parallel to an axis, the corresponding intercept is infinite (the term is absent).
\end{aretenir}

\begin{savaistu}[Syllabus Note: Three Points Define a Plane]
The official syllabus mentions "Lesson 43: Cartesian equation of a plane through three collinear points". This is a known typo: three \textbf{collinear} points lie on infinitely many planes (they only define a line). A unique plane is defined by three \textbf{non-collinear} points. In the exam, you will be given three non-collinear points (or asked to verify they are non-collinear first). The method is: find two direction vectors $\vec{AB}$ and $\vec{AC}$, then $\vec{n} = \vec{AB} \times \vec{AC}$, and use $\vec{r} \cdot \vec{n} = \vec{a} \cdot \vec{n}$.
\end{savaistu}

\begin{exoresolu}[Using Intercept Form]
Find the equation of the plane that cuts the $x$-axis at $4$, the $y$-axis at $-2$, and is parallel to the $z$-axis.
\tcblower
\textbf{Solution:}
Intercepts: $\alpha = 4$, $\beta = -2$. Parallel to $z$-axis $\Rightarrow$ no $z$-intercept ($\gamma \to \infty$, term $\frac{z}{\gamma} \to 0$).
Equation: $\frac{x}{4} + \frac{y}{-2} = 1$.
Multiply by 4: $x - 2y = 4$ or $x - 2y - 4 = 0$.
(Check: Normal vector $\vec{n} = (1, -2, 0)$ has zero $z$-component, so plane is vertical/parallel to $z$-axis.)
\end{exoresolu}

\courssub{Checking Whether a Point Lies on a Plane}

Given a plane in any form (vector, parametric, or Cartesian) and a point $P(x_0, y_0, z_0)$, we substitute the coordinates into the equation.

\begin{itemize}
    \item \textbf{Cartesian $ax+by+cz=d$:} Check if $ax_0 + by_0 + cz_0 = d$.
    \item \textbf{Vector $\vec{r}\cdot\vec{n}=p$:} Check if $\vec{OP}\cdot\vec{n} = p$.
    \item \textbf{Parametric $\vec{r}=\vec{a}+\lambda\vec{u}+\mu\vec{v}$:} Solve the system $x_0 = a_1+\lambda u_1+\mu v_1$, etc., for $\lambda, \mu$. If a consistent solution exists, the point lies on the plane.
\end{itemize}

\begin{exoresolu}[Point on Plane?]
Determine whether the point $P(1, 4, -2)$ lies on the plane $\Pi: \vec{r} = \begin{pmatrix} 2 \\ 1 \\ 0 \end{pmatrix} + \lambda \begin{pmatrix} 1 \\ -1 \\ 2 \end{pmatrix} + \mu \begin{pmatrix} 0 \\ 2 \\ -1 \end{pmatrix}$.
\tcblower
\textbf{Solution:}
We need $\lambda, \mu$ such that:
\[
\begin{pmatrix} 1 \\ 4 \\ -2 \end{pmatrix} = \begin{pmatrix} 2 \\ 1 \\ 0 \end{pmatrix} + \lambda \begin{pmatrix} 1 \\ -1 \\ 2 \end{pmatrix} + \mu \begin{pmatrix} 0 \\ 2 \\ -1 \end{pmatrix}
\]
Component-wise:
\begin{align*}
1 &= 2 + \lambda \quad \Rightarrow \quad \lambda = -1 \quad \text{(1)} \\
4 &= 1 - \lambda + 2\mu \quad \text{(2)} \\
-2 &= 2\lambda - \mu \quad \text{(3)}
\end{align*}
Substitute $\lambda = -1$ into (2): $4 = 1 - (-1) + 2\mu \Rightarrow 4 = 2 + 2\mu \Rightarrow \mu = 1$.
Check (3): $2(-1) - 1 = -3 \neq -2$.
The system is inconsistent. \textbf{Point $P$ does NOT lie on the plane.}
\end{exoresolu}

\begin{exercice}[Practice: Equations of a Plane]
\begin{enumerate}
    \item Find the vector equation $\vec{r}\cdot\vec{n}=p$ of the plane through $A(3, -2, 1)$ with normal $\vec{n} = \vec{i} + 2\vec{j} - 2\vec{k}$. Write the Cartesian equation.
    \item A plane has parametric equation $\vec{r} = \begin{pmatrix} 0 \\ 1 \\ 2 \end{pmatrix} + \lambda \begin{pmatrix} 1 \\ 0 \\ -1 \end{pmatrix} + \mu \begin{pmatrix} 2 \\ -1 \\ 0 \end{pmatrix}$. Find its Cartesian equation.
    \item Find the intercepts of the plane $2x - 3y + 4z = 12$ and write its equation in intercept form.
    \item Show that the point $P(2, -1, 3)$ lies on the plane $\Pi: \vec{r}\cdot(2\vec{i} - \vec{j} + \vec{k}) = 8$. Find the parametric form of $\Pi$ using $P$ as the base point and two direction vectors of your choice.
    \item (GCE Style) The plane $\Pi$ has equation $x + 2y - 2z = 4$. Which of the following vectors is \textbf{not} a normal vector to $\Pi$?
    \begin{itemize}
        \item[A] $\begin{pmatrix} 1 \\ 2 \\ -2 \end{pmatrix}$
        \item[B] $\begin{pmatrix} 2 \\ -1 \\ 1 \end{pmatrix}$
        \item[C] $\begin{pmatrix} -1 \\ -2 \\ 2 \end{pmatrix}$
        \item[D] $\begin{pmatrix} 2 \\ 4 \\ -4 \end{pmatrix}$
    \end{itemize}
\end{enumerate}
\end{exercice}

\begin{corrige}[Exercise 1]
\begin{enumerate}
    \item $\vec{a} = (3, -2, 1)$, $\vec{n} = (1, 2, -2)$. $p = \vec{a}\cdot\vec{n} = 3(1) + (-2)(2) + 1(-2) = 3 - 4 - 2 = -3$.
    Equation: $\vec{r}\cdot\begin{pmatrix} 1 \\ 2 \\ -2 \end{pmatrix} = -3$. Cartesian: $x + 2y - 2z = -3$.
    \item $\vec{u} = (1, 0, -1)$, $\vec{v} = (2, -1, 0)$. $\vec{n} = \vec{u}\times\vec{v} = \begin{vmatrix} \vec{i} & \vec{j} & \vec{k} \\ 1 & 0 & -1 \\ 2 & -1 & 0 \end{vmatrix} = \vec{i}(0 - 1) - \vec{j}(0 - (-2)) + \vec{k}(-1 - 0) = (-1, -2, -1)$.
    Use $\vec{n} = (1, 2, 1)$ (simpler, multiply by $-1$). Point $\vec{a} = (0, 1, 2)$. $p = 0 + 2 + 2 = 4$.
    Cartesian: $x + 2y + z = 4$.
    \item $x$-intercept: $y=z=0 \Rightarrow 2x=12 \Rightarrow x=6 \Rightarrow \alpha=6$.
    $y$-intercept: $x=z=0 \Rightarrow -3y=12 \Rightarrow y=-4 \Rightarrow \beta=-4$.
    $z$-intercept: $x=y=0 \Rightarrow 4z=12 \Rightarrow z=3 \Rightarrow \gamma=3$.
    Intercept form: $\frac{x}{6} + \frac{y}{-4} + \frac{z}{3} = 1$.
    \item Check: $\vec{OP}\cdot\vec{n} = (2)(2) + (-1)(-1) + (3)(1) = 4 + 1 + 3 = 8 = p$. Yes, $P$ lies on $\Pi$.
    Need two vectors $\perp \vec{n}=(2,-1,1)$. Choose $\vec{u}=(1,2,0)$ (check: $2-2+0=0$). Choose $\vec{v}=(0,1,1)$ (check: $0-1+1=0$).
    Parametric: $\vec{r} = \begin{pmatrix} 2 \\ -1 \\ 3 \end{pmatrix} + \lambda \begin{pmatrix} 1 \\ 2 \\ 0 \end{pmatrix} + \mu \begin{pmatrix} 0 \\ 1 \\ 1 \end{pmatrix}$.
    \item A normal vector is $\vec{n} = (1, 2, -2)$. Any non-zero scalar multiple is also a normal vector.
    Option A: $(1, 2, -2)$ is $\vec{n}$ itself. \textbf{Normal}.
    Option B: $(2, -1, 1)$. Dot product with $\vec{n}$: $2(1) + (-1)(2) + 1(-2) = 2 - 2 - 2 = -2 \neq 0$. Not parallel to $\vec{n}$. \textbf{Not a normal vector}.
    Option C: $(-1, -2, 2) = -1 \times \vec{n}$. \textbf{Normal}.
    Option D: $(2, 4, -4) = 2 \times \vec{n}$. \textbf{Normal}.
    \textbf{Answer: B}.
\end{enumerate}
\end{corrige}

\coursec{Plane Through Three Points and Special Planes}

In the previous section we learned that a plane is completely determined as soon as we know \textbf{one point of the plane} and \textbf{one normal vector}: the vector equation $\vec{r}\cdot\vec{n}=\vec{a}\cdot\vec{n}$ and the Cartesian equation $ax+by+cz=d$ both follow immediately. We also saw the parametric form $\vec{r}=\vec{a}+\lambda\vec{u}+\mu\vec{v}$, which needs one point and \textbf{two non-parallel direction vectors} lying in the plane. In this section we put these tools to work on the most frequent GCE situation: \emph{a plane specified by three points}. We then treat a family of quick but very examinable cases: planes parallel to the coordinate planes, planes containing an axis, planes parallel to a given plane, and planes containing a given line and a given point.

\courssub{Plane Through Three Non-Collinear Points}

\textbf{Intuition.} Place three fingertips of one hand on your desk so that they are not in a straight line, and rest a sheet of cardboard on them: the cardboard settles in exactly one position. Three non-collinear points behave like a tripod — they fix a \cle{unique plane}. But if the three fingertips are in a straight line, the cardboard can still rotate around that line: uniqueness is lost. This simple observation is the whole content of the section.

\begin{propriete}[Uniqueness of the plane through three points]
Three \textbf{non-collinear} points $A$, $B$, $C$ of space determine one and only one plane. (The statement is assumed; its use is examinable, not its proof.)
\end{propriete}

\textbf{How do we build the equation?} If $A$, $B$, $C$ lie in the plane, then the vectors $\vec{AB}$ and $\vec{AC}$ lie in the plane. Since $A$, $B$, $C$ are not collinear, $\vec{AB}$ and $\vec{AC}$ are not parallel, so their cross product
\[ \vec{n}=\vec{AB}\times\vec{AC} \]
is a non-zero vector perpendicular to \emph{both} of them, hence perpendicular to the whole plane: it is a \cle{normal vector}. We then apply the point-normal form of Section~1 with the point $A$:
\[ \vec{r}\cdot\vec{n}=\vec{OA}\cdot\vec{n}. \]

\begin{popfigure}
\begin{center}
\begin{tikzpicture}[scale=1.15,>=stealth]
\coordinate (A) at (0,0);
\coordinate (B) at (4.2,0.6);
\coordinate (C) at (1.3,2.6);
\coordinate (D) at (5.5,3.2);
\coordinate (E) at (-0.9,2.2);
% plane parallelogram
\fill[popblueL,opacity=0.55] (E) -- (C) -- (D) -- (B) -- cycle;
\draw[popblue] (E) -- (C) -- (D) -- (B) -- cycle;
% triangle
\fill[popink] (A) circle (1.6pt) node[below left]{$A$};
\fill[popink] (3.1,0.45) circle (1.6pt) node[below]{$B$};
\fill[popink] (1.05,2.05) circle (1.6pt) node[above left]{$C$};
\draw[popdark,thick] (A) -- (3.1,0.45) -- (1.05,2.05) -- cycle;
\draw[->,poporange,very thick] (A) -- (3.1,0.45) node[midway,below]{$\vec{AB}$};
\draw[->,popgreen,very thick] (A) -- (1.05,2.05) node[midway,left]{$\vec{AC}$};
\draw[->,poppurple,very thick] (A) -- (0.35,1.05) node[right]{$\vec{n}=\vec{AB}\times\vec{AC}$};
\end{tikzpicture}
\end{center}
\legende{The plane through $A$, $B$, $C$: $\vec{AB}$ and $\vec{AC}$ lie in the plane, and $\vec{n}=\vec{AB}\times\vec{AC}$ is normal to it.}
\end{popfigure}

\begin{methode}[Plane through three points $A$, $B$, $C$ — cross product method]
\begin{enumerate}
\item Compute $\vec{AB}=\vec{OB}-\vec{OA}$ and $\vec{AC}=\vec{OC}-\vec{OA}$.
\item Compute the normal $\vec{n}=\vec{AB}\times\vec{AC}$. \textbf{Check that $\vec{n}\neq\vec{0}$} (otherwise the points are collinear — see the warning below).
\item Write $\vec{r}\cdot\vec{n}=\vec{OA}\cdot\vec{n}$, or directly the Cartesian form $n_1x+n_2y+n_3z=d$ with $d=\vec{OA}\cdot\vec{n}$.
\item \textbf{Verify}: the coordinates of $A$, $B$ and $C$ must all satisfy the final equation.
\end{enumerate}
\end{methode}

\begin{exoresolu}[Plane through three points]
Find the Cartesian equation of the plane through $A(1,0,2)$, $B(2,1,1)$ and $C(-1,2,1)$.
\tcblower
\textbf{Solution.} Direction vectors in the plane:
\[ \vec{AB}=\begin{pmatrix}1\\1\\-1\end{pmatrix},\qquad \vec{AC}=\begin{pmatrix}-2\\2\\-1\end{pmatrix}. \]
Normal vector:
\[ \vec{n}=\vec{AB}\times\vec{AC}=\begin{pmatrix}(1)(-1)-(-1)(2)\\ (-1)(-2)-(1)(-1)\\ (1)(2)-(1)(-2)\end{pmatrix}=\begin{pmatrix}1\\3\\4\end{pmatrix}. \]
Since $\vec{n}\neq\vec{0}$, the points are non-collinear and determine a unique plane. Then
\[ d=\vec{OA}\cdot\vec{n}=(1)(1)+(0)(3)+(2)(4)=9. \]
Hence the plane has equation $\boxed{x+3y+4z=9}$.
\textbf{Check:} $B$: $2+3+4=9$ \checkmark; $C$: $-1+6+4=9$ \checkmark.
\end{exoresolu}

\courssub{The Direct Cartesian Method (Simultaneous Equations)}

There is an alternative that some candidates prefer because it needs no cross product. Every plane has an equation $ax+by+cz=d$; if the three points lie on it, their coordinates satisfy it, giving three equations in the four unknowns $a,b,c,d$. Since the equation is only defined up to a non-zero multiple, we may fix one coefficient (for instance set $d=1$, or keep $d$ as a parameter) and solve for the ratios $a:b:c:d$.

\begin{methode}[Plane through three points — substitution method]
\begin{enumerate}
\item Substitute the coordinates of $A$, $B$, $C$ into $ax+by+cz=d$ to obtain three simultaneous equations.
\item Solve for $a$, $b$, $c$ in terms of $d$ (or set a convenient value of one coefficient).
\item Substitute back and simplify; clear fractions by multiplying through.
\item Verify that all three points satisfy the equation.
\end{enumerate}
\end{methode}

\begin{exemplebox}[Same plane, second method]
For $A(1,0,2)$, $B(2,1,1)$, $C(-1,2,1)$, substitution into $ax+by+cz=d$ gives
\[ a+2c=d,\qquad 2a+b+c=d,\qquad -a+2b+c=d. \]
From the first equation $a=d-2c$. Substituting into the second: $b=d-2a-c=d-2(d-2c)-c=-d+3c$. The third equation then gives $-(d-2c)+2(-d+3c)+c=d$, i.e.\ $-3d+9c=d$, so $c=\dfrac{d}{2.5}\cdot\ldots$ — more cleanly, $9c=4d$, so take $d=9$, $c=4$, whence $a=9-8=1$ and $b=-9+12=3$. Equation: $x+3y+4z=9$, exactly as before.
\end{exemplebox}

\begin{attention}[Choose your method wisely]
The cross product method is usually faster and less error-prone. The substitution method is useful when the cross product is awkward or when the question explicitly asks you to ``show'' the equation by substitution. Whichever you use, \textbf{always verify the three points} in your final answer — this catches arithmetic slips and earns method marks even after an error.
\end{attention}

\courssub{The GCE Trap: Three Collinear Points}

The syllabus explicitly mentions ``the equation of a plane through three collinear points'' — this is a deliberate trap. If $A$, $B$, $C$ are collinear, then $\vec{AC}=k\,\vec{AB}$ for some scalar $k$, so
\[ \vec{AB}\times\vec{AC}=k\,(\vec{AB}\times\vec{AB})=\vec{0}, \]
and the method collapses: there is \textbf{no normal vector}. Geometrically, the three points all lie on a single line $\ell$, and \emph{infinitely many planes contain a given line} — imagine all the pages of an open book hinged along its spine.

\begin{popfigure}
\begin{center}
\begin{tikzpicture}[scale=1.0,>=stealth]
% common line
\draw[popink,very thick] (-3.4,0) -- (3.4,0) node[right]{$\ell$};
\fill[popink] (-1.8,0) circle (1.6pt) node[below]{$A$};
\fill[popink] (0,0) circle (1.6pt) node[below]{$B$};
\fill[popink] (1.8,0) circle (1.6pt) node[below]{$C$};
% plane 1
\fill[popblueL,opacity=0.6] (-2.6,-1.2) -- (2.6,-1.2) -- (3.0,0) -- (-3.0,0) -- cycle;
\draw[popblue] (-2.6,-1.2) -- (2.6,-1.2);
% plane 2
\fill[popgreenL,opacity=0.55] (-2.2,1.7) -- (2.2,1.7) -- (2.9,0) -- (-2.9,0) -- cycle;
\draw[popgreen] (-2.2,1.7) -- (2.2,1.7);
% plane 3 (slanted)
\fill[poporange,opacity=0.35] (-2.7,-0.9) -- (1.5,1.6) -- (2.6,1.6) -- (-1.6,-0.9) -- cycle;
\draw[poporange] (-2.7,-0.9) -- (1.5,1.6);
\node[popblue] at (2.9,-1.1){$\Pi_1$};
\node[popgreen] at (2.5,1.75){$\Pi_2$};
\node[popdark] at (2.9,1.3){$\Pi_3$};
\end{tikzpicture}
\end{center}
\legende{Three collinear points lie on a line $\ell$; infinitely many planes $\Pi_1,\Pi_2,\Pi_3,\ldots$ are ``hinged'' about $\ell$. No unique plane is determined.}
\end{popfigure}

\begin{attention}[Three collinear points do not fix a plane]
If $A$, $B$, $C$ are collinear, then $\vec{AB}\times\vec{AC}=\vec{0}$ and there are \textbf{infinitely many planes} through them. Before computing anything, check that $\vec{AB}\times\vec{AC}\neq\vec{0}$. If an exam question gives you three collinear points, the correct response is to state that the plane is not unique and that an extra condition (a fourth point off the line, or a prescribed normal) is required.
\end{attention}

\courssub{Special Planes: Coordinate Planes, Axes, Parallel Planes}

\textbf{Planes parallel to the coordinate planes.} The coordinate plane $Oxy$ has equation $z=0$ and normal $\vec{k}=(0,0,1)$. Consequently any plane \textbf{parallel to $Oxy$} has equation $z=k$; similarly planes parallel to $Oxz$ have equation $y=k$, and planes parallel to $Oyz$ have equation $x=k$.

\textbf{Planes containing a coordinate axis.} A plane containing the $z$-axis passes through the origin and has normal perpendicular to $\vec{k}$, i.e.\ $c=0$ and $d=0$: its equation is $ax+by=0$. Similarly, planes containing the $x$-axis: $by+cz=0$; the $y$-axis: $ax+cz=0$.

\begin{propriete}[Parallel planes]
Two planes are parallel if and only if their normal vectors are proportional: $\vec{n}_1=\lambda\vec{n}_2$. Hence the plane \textbf{through a point $P$ and parallel to the plane $ax+by+cz=d$} keeps the same left-hand side and only the constant changes:
\[ ax+by+cz=ax_P+by_P+cz_P. \]
\end{propriete}

\begin{exemplebox}[Plane through a point, parallel to a given plane]
Find the plane through $P(2,-1,3)$ parallel to $\Pi: 2x-y+4z=7$. Same normal $(2,-1,4)$, so the equation is $2x-y+4z=d'$ with $d'=2(2)-(-1)+4(3)=4+1+12=17$. Answer: $2x-y+4z=17$.
\end{exemplebox}

\courssub{Plane Containing a Given Line and a Given Point}

\begin{attention}[Do not confuse the two situations]
``Plane through three points'' and ``plane through a line $\ell$ and a point $P\notin\ell$'' are \textbf{different strategies}. In the second case you do not have three points given directly — you must \emph{manufacture} them: pick two points $A$, $B$ on the line (e.g.\ by taking two values of the parameter), then apply the three-point method to $A$, $B$, $P$. Note that a direction vector $\vec{u}$ of the line and the vector $\vec{AP}$ both lie in the plane, so a shortcut is $\vec{n}=\vec{u}\times\vec{AP}$.
\end{attention}

\begin{exoresolu}[Plane containing a line and a point]
Find the equation of the plane containing the line $\vec{r}=\begin{pmatrix}1\\0\\1\end{pmatrix}+t\begin{pmatrix}2\\-1\\1\end{pmatrix}$ and the point $P(3,2,0)$.
\tcblower
\textbf{Solution.} The point $A(1,0,1)$ lies on the line, and $\vec{u}=(2,-1,1)$ is a direction vector of the line, hence parallel to the plane. Also
\[ \vec{AP}=\vec{OP}-\vec{OA}=\begin{pmatrix}2\\2\\-1\end{pmatrix} \]
lies in the plane. A normal vector is
\[ \vec{n}=\vec{u}\times\vec{AP}=\begin{pmatrix}(-1)(-1)-(1)(2)\\ (1)(2)-(2)(-1)\\ (2)(2)-(-1)(2)\end{pmatrix}=\begin{pmatrix}-1\\4\\6\end{pmatrix}. \]
Then $d=\vec{OA}\cdot\vec{n}=-1+0+6=5$, so the plane is $-x+4y+6z=5$, equivalently
\[ \boxed{x-4y-6z=-5}. \]
\textbf{Check:} $A$: $1-0-6=-5$ \checkmark; $P$: $3-8-0=-5$ \checkmark; a second point of the line, $t=1$: $(3,-1,2)$ gives $3+4-12=-5$ \checkmark.
\end{exoresolu}

\begin{aretenir}[Key points of this section]
\begin{itemize}
\item Three \textbf{non-collinear} points determine a unique plane: $\vec{n}=\vec{AB}\times\vec{AC}$, then $\vec{r}\cdot\vec{n}=\vec{OA}\cdot\vec{n}$.
\item Alternative: substitute the three points into $ax+by+cz=d$ and solve the simultaneous equations.
\item Three \textbf{collinear} points give $\vec{AB}\times\vec{AC}=\vec{0}$ and infinitely many planes — always check.
\item Planes parallel to $Oxy$, $Oxz$, $Oyz$: $z=k$, $y=k$, $x=k$. Planes through an axis have $d=0$ and one zero coefficient.
\item Parallel planes share the same normal: only the constant $d$ changes.
\item Line + external point: use $\vec{n}=\vec{u}\times\vec{AP}$ where $\vec{u}$ directs the line and $A$ is on the line.
\end{itemize}
\end{aretenir}

\begin{exercice}
\begin{enumerate}
\item Find the Cartesian equation of the plane through $A(2,1,-1)$, $B(0,3,1)$ and $C(1,-1,2)$.
\item Show that the points $A(1,2,3)$, $B(2,3,4)$ and $C(3,4,5)$ do not determine a unique plane, and explain why geometrically.
\item Find the equation of the plane through $(1,-2,4)$ parallel to the plane $3x+2y-z=6$.
\item Find the equation of the plane containing the line $\vec{r}=(1,1,0)+t(1,2,-1)$ and the point $Q(0,1,3)$.
\end{enumerate}
\end{exercice}

\begin{corrige}[Exercise 1]
$\vec{AB}=(-2,2,2)$, $\vec{AC}=(-1,-2,3)$. $\vec{n}=\vec{AB}\times\vec{AC}=(2\cdot3-2\cdot(-2),\ 2\cdot(-1)-(-2)\cdot3,\ (-2)(-2)-2(-1))=(10,4,6)$, simplify to $(5,2,3)$. $d=5(2)+2(1)+3(-1)=9$. Plane: $5x+2y+3z=9$. Check $B$: $0+6+3=9$ \checkmark; $C$: $5-2+6=9$ \checkmark.
\end{corrige}

\begin{corrige}[Exercise 2]
$\vec{AB}=(1,1,1)$, $\vec{AC}=(2,2,2)=2\vec{AB}$, so $\vec{AB}\times\vec{AC}=\vec{0}$: the points are collinear (they lie on the line $\vec{r}=(1,2,3)+t(1,1,1)$). Infinitely many planes contain this line, so no unique plane is determined; an extra condition is needed.
\end{corrige}

\begin{corrige}[Exercise 3]
Same normal $(3,2,-1)$: $3x+2y-z=d'$ with $d'=3(1)+2(-2)-4=-5$. Plane: $3x+2y-z=-5$.
\end{corrige}

\begin{corrige}[Exercise 4]
$A(1,1,0)$ on the line, $\vec{u}=(1,2,-1)$, $\vec{AQ}=(-1,0,3)$. $\vec{n}=\vec{u}\times\vec{AQ}=(2\cdot3-(-1)\cdot0,\ (-1)(-1)-1\cdot3,\ 1\cdot0-2(-1))=(6,-2,2)$, simplify to $(3,-1,1)$. $d=3(1)-1+0=2$. Plane: $3x-y+z=2$. Check $Q$: $0-1+3=2$ \checkmark.
\end{corrige}

\begin{savaistu}[Why the tripod never wobbles]
Carpenters, surveyors and photographers all exploit this section's first result: a tripod stands firm on any ground because three (non-collinear) feet determine exactly one plane, whereas a four-legged table can rock if the four feet are not coplanar. The theodolites used by surveyors on road projects across Cameroon are mounted on tripods for precisely this reason — geometry you can literally stand on!
\end{savaistu}

\coursec{Perpendicular Distance from a Point to a Plane}

In the previous section we learned how to pin down a plane completely: through three non-collinear points, through a point with a given normal, or through a line with an extra condition. Once a plane is fixed in space, a natural question arises: \emph{how far is a given point from that plane?} This question is not a luxury of pure geometry. When an engineer checks the clearance between a component and a wall, when a pilot must know the altitude above a tilted runway surface, or when an architect verifies that a beam stays clear of a sloping roof, the computation is exactly the \cle{perpendicular distance from a point to a plane}. This section gives you the complete toolbox: the distance formula and its derivation, the signed distance, the distance between parallel planes, the foot of the perpendicular, the mirror image of a point, and locus problems leading to bisector planes.

\courssub{Geometric motivation: what do we mean by ``distance''?}

Stand in a classroom and look at the floor. You are a point $P$ above a plane $\Pi$. You could walk from your position to \emph{any} point of the floor: the distances to those points are all different. But there is exactly one path that is shortest: drop straight down, along the vertical. In other words, the shortest route from $P$ to $\Pi$ is along the \cle{normal} to the plane.

\begin{definition}[Perpendicular distance from a point to a plane]
Let $\Pi$ be a plane and $P$ a point in space. The \cle{perpendicular distance} (or \emph{shortest distance}) from $P$ to $\Pi$ is the length $PN$, where $N$ is the unique point of $\Pi$ such that $\vec{PN}$ is perpendicular to $\Pi$. The point $N$ is called the \cle{foot of the perpendicular} from $P$ to $\Pi$.
\end{definition}

Why is $PN$ the shortest? Take any other point $M$ of the plane. The triangle $PNM$ is right-angled at $N$ (because $PN \perp \Pi$ and $NM \subset \Pi$), so by Pythagoras $PM^2 = PN^2 + NM^2 \geq PN^2$, with equality only when $M = N$. Hence $PN \leq PM$ for every $M \in \Pi$: the perpendicular truly gives the minimum distance.

\courssub{Derivation of the distance formula}

Let the plane $\Pi$ have Cartesian equation $ax + by + cz = d$, so that $\vec{n} = \begin{pmatrix} a \\ b \\ c \end{pmatrix}$ is a normal vector. Let $P(x_1, y_1, z_1)$ be any point in space. Choose any point $A(x_0, y_0, z_0)$ lying \emph{on} the plane, so $ax_0 + by_0 + cz_0 = d$.

The key observation is that the perpendicular distance $PN$ is the \emph{absolute value of the scalar projection} of $\vec{AP}$ onto the normal direction. Indeed, projecting $\vec{AP}$ onto $\vec{n}$ collapses the component of $\vec{AP}$ lying inside the plane and keeps exactly the perpendicular part, which has length $PN$.

\begin{popfigure}
\begin{center}
\begin{tikzpicture}[scale=1.05, >=stealth]
  % plane
  \draw[fill=popblueL, draw=popblue, thick] (-3.4,-0.7) -- (3.4,-0.7) -- (4.4,0.9) -- (-2.4,0.9) -- cycle;
  \node[popblue] at (3.6,0.55) {$\Pi$};
  % points
  \coordinate (P) at (0.6,3.1);
  \coordinate (N) at (0.6,0.1);
  \coordinate (A) at (-1.7,0.35);
  \fill[popdark] (P) circle (2pt) node[above right] {$P(x_1,y_1,z_1)$};
  \fill[poporange] (N) circle (2pt) node[below right] {$N$};
  \fill[popgreen] (A) circle (2pt) node[below left] {$A$};
  % perpendicular
  \draw[very thick, poporange] (P) -- (N) node[midway, right] {$d = PN$};
  % right angle mark
  \draw[popdark] (0.6,0.32) -- (0.82,0.32) -- (0.82,0.1);
  % AP vector
  \draw[->, very thick, popgreen] (A) -- (P) node[midway, above left] {$\vec{AP}$};
  % unit normal
  \draw[->, very thick, poppurple] (N) -- (0.6,1.5) node[midway, right] {$\hat{n}$};
  % dashed projection hint
  \draw[dashed, popmuted] (A) -- (N);
\end{tikzpicture}
\end{center}
\legende{The distance from $P$ to $\Pi$ is the absolute value of the scalar projection of $\vec{AP}$ onto the unit normal $\hat{n}$.}
\end{popfigure}

The scalar projection of $\vec{AP}$ onto $\vec{n}$ is $\dfrac{\vec{AP} \cdot \vec{n}}{|\vec{n}|}$. Now compute:
\[
\vec{AP} \cdot \vec{n} = a(x_1 - x_0) + b(y_1 - y_0) + c(z_1 - z_0) = (ax_1 + by_1 + cz_1) - (ax_0 + by_0 + cz_0).
\]
Since $A \in \Pi$, the second bracket equals $d$, so
\[
\vec{AP} \cdot \vec{n} = ax_1 + by_1 + cz_1 - d, \qquad |\vec{n}| = \sqrt{a^2 + b^2 + c^2}.
\]
Taking the absolute value of the projection gives the distance. We have proved the central result of this section.

\begin{propriete}[Perpendicular distance from a point to a plane — examinable derivation]
The perpendicular distance from the point $P(x_1, y_1, z_1)$ to the plane $ax + by + cz = d$ is
\[
\[\boxed{\; \delta = \frac{|ax_1 + by_1 + cz_1 - d|}{\sqrt{a^2 + b^2 + c^2}} \;}\]
\]
The derivation above (scalar projection of $\vec{AP}$ onto $\hat{n}$) is examinable at GCE Advanced Level: you may be asked to \emph{establish} the formula, not merely to use it.
\end{propriete}

\begin{attention}[The two most common errors]
\begin{itemize}
\item \textbf{Forgetting the absolute value.} A distance is never negative. Writing $\delta = \dfrac{ax_1+by_1+cz_1-d}{\sqrt{a^2+b^2+c^2}}$ without the modulus can produce a negative ``distance'' --- an instant loss of marks.
\item \textbf{Forgetting the square root.} The denominator is $\sqrt{a^2+b^2+c^2}$, the \emph{length} of the normal vector, not $a^2+b^2+c^2$. Check: the numerator is in length units times coefficient units, so the denominator must carry the same coefficient units --- only the square root achieves this.
\end{itemize}
\end{attention}

\begin{exoresolu}[Direct application of the formula]
Find the perpendicular distance from the point $P(2, -1, 3)$ to the plane $\Pi: 2x - y + 2z = 5$.
\tcblower
Here $a = 2$, $b = -1$, $c = 2$, $d = 5$, and $(x_1, y_1, z_1) = (2, -1, 3)$.
\[
ax_1 + by_1 + cz_1 - d = 2(2) - (-1) + 2(3) - 5 = 4 + 1 + 6 - 5 = 6,
\]
\[
\sqrt{a^2+b^2+c^2} = \sqrt{4+1+4} = \sqrt{9} = 3.
\]
Therefore
\[
\delta = \frac{|6|}{3} = 2 \text{ units.}
\]
\end{exoresolu}

\courssub{Signed distance: which side of the plane?}

If we simply \emph{remove} the absolute value, the quantity
\[
\sigma(P) = \frac{ax_1 + by_1 + cz_1 - d}{\sqrt{a^2+b^2+c^2}}
\]
is called the \cle{signed distance} from $P$ to the plane. Its sign tells us on which \emph{side} of $\Pi$ the point lies, relative to the direction of the chosen normal $\vec{n}$:

\begin{itemize}
\item $\sigma(P) > 0$: $P$ lies on the side towards which $\vec{n}$ points;
\item $\sigma(P) < 0$: $P$ lies on the opposite side;
\item $\sigma(P) = 0$: $P$ lies on the plane itself.
\end{itemize}

\begin{exemplebox}[Testing two points]
For the plane $2x - y + 2z = 5$, consider $P(2,-1,3)$ and $Q(0,0,0)$. We found $\sigma(P) = +2$. For $Q$: $\sigma(Q) = \dfrac{0 - 0 + 0 - 5}{3} = -\dfrac{5}{3}$. The signs differ, so $P$ and $Q$ lie on \emph{opposite sides} of the plane. This trick is extremely useful for checking whether a segment crosses a plane without computing the intersection point.
\end{exemplebox}

\courssub{Distance between two parallel planes}

Two planes are parallel exactly when they share the same normal direction. Suppose
\[
\Pi_1: ax + by + cz = d_1 \qquad \text{and} \qquad \Pi_2: ax + by + cz = d_2.
\]
Take any point $A(x_0,y_0,z_0)$ on $\Pi_1$, so $ax_0+by_0+cz_0 = d_1$. The distance between the planes equals the distance from $A$ to $\Pi_2$:
\[
\delta = \frac{|ax_0 + by_0 + cz_0 - d_2|}{\sqrt{a^2+b^2+c^2}} = \frac{|d_1 - d_2|}{\sqrt{a^2+b^2+c^2}}.
\]

\begin{propriete}[Distance between parallel planes]
The distance between the parallel planes $ax + by + cz = d_1$ and $ax + by + cz = d_2$ is
\[
\delta = \frac{|d_1 - d_2|}{\sqrt{a^2+b^2+c^2}}.
\]
\end{propriete}

\begin{popfigure}
\begin{center}
\begin{tikzpicture}[scale=0.95, >=stealth]
  % lower plane
  \draw[fill=popblueL, draw=popblue, thick] (-3.2,-0.6) -- (3.2,-0.6) -- (4.2,1.0) -- (-2.2,1.0) -- cycle;
  \node[popblue] at (3.7,0.7) {$\Pi_1$};
  % upper plane
  \draw[fill=popgreenL, draw=popgreen, thick] (-3.2,2.2) -- (3.2,2.2) -- (4.2,3.8) -- (-2.2,3.8) -- cycle;
  \node[popgreen!60!black] at (3.7,3.5) {$\Pi_2$};
  % common normal
  \draw[<->, very thick, poporange] (0.3,0.2) -- (0.3,3.0) node[midway, right] {$\delta = \dfrac{|d_1-d_2|}{\sqrt{a^2+b^2+c^2}}$};
  \fill[popdark] (0.3,0.2) circle (2pt) node[below left] {$A$};
  \draw[->, thick, poppurple] (0.3,3.0) -- (0.3,3.9) node[right] {$\vec{n}$};
\end{tikzpicture}
\end{center}
\legende{Two parallel planes share the same normal; the distance between them is measured along any common perpendicular.}
\end{popfigure}

\begin{attention}[Same normal coefficients before subtracting]
The formula $|d_1 - d_2|/\sqrt{a^2+b^2+c^2}$ is valid \textbf{only} if the two equations are written with the \emph{same} coefficients $a, b, c$. For example, the planes $2x - y + 2z = 5$ and $4x - 2y + 4z = 6$ are parallel, but you must first rewrite the second as $2x - y + 2z = 3$ (dividing by 2) \emph{before} subtracting the constants. Subtracting $5 - 6 = -1$ directly would give the wrong answer.
\end{attention}

\begin{exoresolu}[Distance between parallel planes]
Find the distance between the planes $\Pi_1: x - 2y + 2z = 3$ and $\Pi_2: 2x - 4y + 4z = 21$.
\tcblower
First rewrite $\Pi_2$ with the same normal as $\Pi_1$: dividing by 2,
\[
\Pi_2: x - 2y + 2z = \tfrac{21}{2}.
\]
Now apply the formula with $d_1 = 3$, $d_2 = \tfrac{21}{2}$ and $\sqrt{a^2+b^2+c^2} = \sqrt{1+4+4} = 3$:
\[
\delta = \frac{\left|3 - \frac{21}{2}\right|}{3} = \frac{\frac{15}{2}}{3} = \frac{5}{2} = 2.5 \text{ units.}
\]
\end{exoresolu}

\courssub{Foot of the perpendicular and mirror image}

Beyond the distance itself, examiners sometimes ask for the \emph{foot of the perpendicular} $N$ or the \emph{mirror image} $P'$ of $P$ in the plane. Both are obtained by the same idea: parametrise the line through $P$ in the direction of the normal.

\begin{methode}[Finding the foot of the perpendicular]
To find the foot $N$ of the perpendicular from $P(x_1,y_1,z_1)$ to the plane $ax+by+cz=d$:
\begin{enumerate}
\item Write the line through $P$ along $\vec{n} = (a,b,c)$: $\vec{r} = \begin{pmatrix} x_1 \\ y_1 \\ z_1 \end{pmatrix} + t \begin{pmatrix} a \\ b \\ c \end{pmatrix}$.
\item Substitute $x = x_1 + at$, $y = y_1 + bt$, $z = z_1 + ct$ into the plane equation and solve for $t$.
\item Substitute this value of $t$ back into the line: the point obtained is $N$.
\item For the \textbf{mirror image} $P'$ of $P$, note that $N$ is the midpoint of $PP'$, so $P' = 2N - P$ (equivalently, use the parameter value $2t$ on the line).
\end{enumerate}
\end{methode}

\begin{exoresolu}[Foot of the perpendicular and mirror image]
Find the foot of the perpendicular from $P(1, 2, 3)$ to the plane $x + y + z = 3$, and deduce the mirror image of $P$ in the plane.
\tcblower
The normal is $\vec{n} = (1,1,1)$. The perpendicular line through $P$ is
\[
x = 1 + t, \qquad y = 2 + t, \qquad z = 3 + t.
\]
Substituting into the plane equation:
\[
(1+t) + (2+t) + (3+t) = 3 \;\Longrightarrow\; 6 + 3t = 3 \;\Longrightarrow\; t = -1.
\]
Hence the foot is $N = (1-1,\ 2-1,\ 3-1) = (0, 1, 2)$. Check: $0+1+2 = 3$ \checkmark.

The mirror image is $P' = 2N - P = (2(0)-1,\ 2(1)-2,\ 2(2)-3) = (-1, 0, 1)$. As a check, the distances: $PN = \sqrt{1+1+1} = \sqrt{3}$ and $NP' = \sqrt{1+1+1} = \sqrt{3}$, equal as expected.
\end{exoresolu}

\courssub{Locus problems: points equidistant from two planes}

A classic GCE locus question asks for the set of points \emph{equidistant from two given planes}. If $\Pi_1: a_1x + b_1y + c_1z = d_1$ and $\Pi_2: a_2x + b_2y + c_2z = d_2$, then $M(x,y,z)$ is equidistant from both planes when
\[
\frac{|a_1x + b_1y + c_1z - d_1|}{\sqrt{a_1^2+b_1^2+c_1^2}} = \frac{|a_2x + b_2y + c_2z - d_2|}{\sqrt{a_2^2+b_2^2+c_2^2}}.
\]
Removing the absolute values gives \emph{two} equations (one with $+$, one with $-$), each of which is linear in $x, y, z$: the locus is a \textbf{pair of planes}. When $\Pi_1$ and $\Pi_2$ intersect, these are the two \cle{bisector planes} of the dihedral angles formed by $\Pi_1$ and $\Pi_2$ --- exactly the planes that bisect the angles between the two planes, a notion we shall meet again in Section 6 when studying pencils of planes.

\begin{exoresolu}[Bisector planes]
Find the locus of points equidistant from the planes $\Pi_1: 3x - 4y = 5$ and $\Pi_2: 4x + 3y = 10$.
\tcblower
A point $M(x,y,z)$ is equidistant from both planes when
\[
\frac{|3x - 4y - 5|}{\sqrt{9+16}} = \frac{|4x + 3y - 10|}{\sqrt{16+9}}.
\]
Both denominators equal $5$, so $|3x - 4y - 5| = |4x + 3y - 10|$, which splits into two cases.

\textbf{Case 1:} $3x - 4y - 5 = 4x + 3y - 10 \;\Longrightarrow\; x + 7y = 5$.

\textbf{Case 2:} $3x - 4y - 5 = -(4x + 3y - 10) \;\Longrightarrow\; 7x - y = 15$.

The locus is the pair of planes $x + 7y = 5$ and $7x - y = 15$. Note that their normals $(1,7,0)$ and $(7,-1,0)$ are perpendicular ($7 - 7 + 0 = 0$): the two bisector planes are always perpendicular to each other.
\end{exoresolu}

\begin{attention}[Reading the plane equation correctly]
Before applying the distance formula, make sure the equation is in the form $ax + by + cz = d$ with \emph{everything on one side}. For the plane $3x - 4y = 5$, the missing $z$ means $c = 0$ --- do not invent a coefficient. Also, if the equation is given as $ax+by+cz+d=0$, identify the constant carefully and move it to the right-hand side with the correct sign.
\end{attention}

\begin{aretenir}[Key points of this section]
\begin{itemize}
\item The distance from $P(x_1,y_1,z_1)$ to $ax+by+cz=d$ is $\delta = \dfrac{|ax_1+by_1+cz_1-d|}{\sqrt{a^2+b^2+c^2}}$, derived by projecting $\vec{AP}$ onto the unit normal.
\item The \textbf{signed distance} (without the modulus) tells you on which side of the plane the point lies.
\item Distance between parallel planes $ax+by+cz=d_1$ and $ax+by+cz=d_2$: $\delta = \dfrac{|d_1-d_2|}{\sqrt{a^2+b^2+c^2}}$, after writing both equations with the \emph{same} normal coefficients.
\item Foot of the perpendicular: parametrise the line through $P$ along $\vec{n}$, meet the plane. Mirror image: $P' = 2N - P$.
\item Points equidistant from two planes form a pair of bisector planes, obtained by equating the two distance expressions with $\pm$.
\end{itemize}
\end{aretenir}

\begin{exercice}[Practice]
\begin{enumerate}
\item Find the distance from $A(3, -2, 1)$ to the plane $2x + 6y - 3z = 11$.
\item Show that the origin and the point $(4, 2, 6)$ lie on opposite sides of the plane $x + 2y + 2z = 15$.
\item Find the distance between the planes $2x - 3y + 6z = 7$ and $4x - 6y + 12z = 35$.
\item Find the foot of the perpendicular from $P(2, 0, -1)$ to the plane $x - y + z = 2$, and the image of $P$ by reflection in this plane.
\item Find the pair of bisector planes of $x + 2y - 2z = 3$ and $2x - y + 2z = 4$.
\end{enumerate}
\end{exercice}

\begin{corrige}[Exercise 1]
$\delta = \dfrac{|2(3) + 6(-2) - 3(1) - 11|}{\sqrt{4+36+9}} = \dfrac{|6 - 12 - 3 - 11|}{7} = \dfrac{20}{7}$ units.
\end{corrige}

\begin{corrige}[Exercise 2]
Signed distances: for $O$: $\dfrac{0 - 15}{3} = -5 < 0$; for $(4,2,6)$: $\dfrac{4 + 4 + 12 - 15}{3} = \dfrac{5}{3} > 0$. Opposite signs, so the points lie on opposite sides.
\end{corrige}

\begin{corrige}[Exercise 3]
Rewrite the second plane: $2x - 3y + 6z = \tfrac{35}{2}$. Then $\delta = \dfrac{\left|\frac{35}{2} - 7\right|}{\sqrt{4+9+36}} = \dfrac{\frac{21}{2}}{7} = \dfrac{3}{2}$ units.
\end{corrige}

\begin{corrige}[Exercise 4]
Line: $x = 2 + t$, $y = -t$, $z = -1 + t$. Substituting: $(2+t) - (-t) + (-1+t) = 2 \Rightarrow 1 + 3t = 2 \Rightarrow t = \tfrac{1}{3}$. Foot: $N = \left(\tfrac{7}{3}, -\tfrac{1}{3}, -\tfrac{2}{3}\right)$. Image: $P' = 2N - P = \left(\tfrac{8}{3}, -\tfrac{2}{3}, -\tfrac{1}{3}\right)$.
\end{corrige}

\begin{corrige}[Exercise 5]
Both normals have length $3$, so $|x + 2y - 2z - 3| = |2x - y + 2z - 4|$. Case 1: $x + 2y - 2z - 3 = 2x - y + 2z - 4 \Rightarrow x - 3y + 4z = 1$. Case 2: $x + 2y - 2z - 3 = -(2x - y + 2z - 4) \Rightarrow 3x + y = 7$. The bisector planes are $x - 3y + 4z = 1$ and $3x + y = 7$.
\end{corrige}

\begin{savaistu}[Why surveyors love this formula]
In road construction across the hills of the North West Region, engineers must constantly check clearances between design points and inclined surfaces such as roadbeds and embankments. Modern surveying software computes exactly the quantity $\dfrac{|ax_1+by_1+cz_1-d|}{\sqrt{a^2+b^2+c^2}}$ thousands of times per second. The mathematics you are learning here is the engine inside those instruments.
\end{savaistu}

\coursec{Angles: Between Two Planes and Between a Line and a Plane}

In the previous section we measured \emph{distances} involving planes: the perpendicular distance from a point to a plane. We now turn to \emph{angles}. Two natural questions arise. When two planes meet along a line (think of the two wings of an open book), how ``open'' are they? And when a line pierces a plane (think of a javelin stuck in the football pitch at the Amadou Ahidjo stadium), at what angle does it strike the ground? Both questions are answered beautifully using nothing more than the \cle{scalar product} and the \cle{normal vector} of a plane.

\courssub{The Angle Between Two Planes}

\textbf{Intuition.} Stand two exercise books against each other so that they meet along a common edge. The ``angle'' between them is measured in a plane \emph{perpendicular to that edge} — this is the \cle{dihedral angle}. Now, the normal vector of a plane is perpendicular to every line in the plane, so the angle between the two normals captures exactly the same opening. This observation is the key to everything that follows.

\begin{definition}[Angle between two planes]
Let $\mathcal{P}_1$ and $\mathcal{P}_2$ be two planes with normal vectors $\vec{n}_1$ and $\vec{n}_2$. The \cle{angle between the two planes} is defined as the angle between their normal vectors, taken as the \textbf{acute angle} (or right angle), i.e. the angle $\theta$ with $0 \leq \theta \leq \dfrac{\pi}{2}$.
\end{definition}

\begin{attention}[Why the absolute value?]
A plane has \emph{two} opposite normal directions ($\vec{n}$ and $-\vec{n}$ both work). Depending on the choice, the angle between the normals could be $\theta$ or $\pi - \theta$. Since the angle between two \emph{planes} must be unambiguous, we always take the \textbf{acute} value. In practice this means taking the absolute value of the scalar product in the formula below.
\end{attention}

\begin{propriete}[Formula for the angle between two planes — proof examinable]
Let $\mathcal{P}_1: a_1x+b_1y+c_1z+d_1=0$ and $\mathcal{P}_2: a_2x+b_2y+c_2z+d_2=0$ have normals $\vec{n}_1=\begin{pmatrix}a_1\\b_1\\c_1\end{pmatrix}$ and $\vec{n}_2=\begin{pmatrix}a_2\\b_2\\c_2\end{pmatrix}$. Then the angle $\theta$ between the planes satisfies
\[
\cos\theta=\frac{|\vec{n}_1\cdot\vec{n}_2|}{|\vec{n}_1|\,|\vec{n}_2|}.
\]
\textbf{Proof.} By definition, $\theta$ is the acute angle between $\vec{n}_1$ and $\vec{n}_2$. The scalar product gives $\vec{n}_1\cdot\vec{n}_2=|\vec{n}_1||\vec{n}_2|\cos\alpha$, where $\alpha$ is the angle between the chosen directions of the normals. If $\alpha$ is obtuse, $\cos\alpha<0$, but the acute angle between the planes is $\pi-\alpha$ with $\cos(\pi-\alpha)=-\cos\alpha=|\cos\alpha|$. Hence in all cases $\cos\theta=\dfrac{|\vec{n}_1\cdot\vec{n}_2|}{|\vec{n}_1||\vec{n}_2|}$. $\blacksquare$
\end{propriete}

\begin{popfigure}
\begin{center}
\begin{tikzpicture}[scale=1.1]
% Plane 1 (horizontal-ish parallelogram)
\draw[fill=popblueL, draw=popblue, thick] (-3,-0.8) -- (0.5,-0.8) -- (2.5,0.9) -- (-1,0.9) -- cycle;
\node[popblue] at (-2.3,0.45) {$\mathcal{P}_1$};
% Plane 2 (tilted parallelogram)
\draw[fill=popgreenL, draw=popgreen, thick] (-0.6,-1.6) -- (2.2,-0.2) -- (3.4,2.0) -- (0.6,0.6) -- cycle;
\node[popgreen] at (2.9,1.5) {$\mathcal{P}_2$};
% Line of intersection
\draw[popdark, very thick] (-1.4,0.05) -- (2.9,0.05);
\node[popdark, below] at (2.3,-0.05) {intersection line};
% Normals
\draw[-latex, very thick, poporange] (0.6,0.05) -- (0.6,1.9) node[right] {$\vec{n}_1$};
\draw[-latex, very thick, poppurple] (0.6,0.05) -- (-0.55,1.35) node[above] {$\vec{n}_2$};
% Angle arc
\draw[popink, thick] (0.6,1.0) arc[start angle=90, end angle=123, radius=0.95];
\node[popink] at (0.28,1.35) {$\theta$};
\end{tikzpicture}
\end{center}
\legende{Two intersecting planes with their normal vectors; the dihedral angle $\theta$ equals the (acute) angle between $\vec{n}_1$ and $\vec{n}_2$.}
\end{popfigure}

\begin{propriete}[Perpendicular and parallel planes]
Let $\mathcal{P}_1$ and $\mathcal{P}_2$ have normals $\vec{n}_1,\vec{n}_2$. Then:
\begin{itemize}
\item $\mathcal{P}_1 \perp \mathcal{P}_2 \iff \vec{n}_1\cdot\vec{n}_2 = 0$ \quad (the normals are perpendicular);
\item $\mathcal{P}_1 \parallel \mathcal{P}_2 \iff \vec{n}_1 \parallel \vec{n}_2 \iff \vec{n}_1 = k\,\vec{n}_2$ for some $k\in\mathbb{R}$, $k\neq 0$.
\end{itemize}
\end{propriete}

\begin{exemplebox}[Angle between two planes]
Find the acute angle between the planes $\mathcal{P}_1: 2x+y-2z=5$ and $\mathcal{P}_2: 3x-6y-2z=7$.

\textbf{Solution.} The normals are $\vec{n}_1=\begin{pmatrix}2\\1\\-2\end{pmatrix}$ and $\vec{n}_2=\begin{pmatrix}3\\-6\\-2\end{pmatrix}$.
\[
\vec{n}_1\cdot\vec{n}_2 = (2)(3)+(1)(-6)+(-2)(-2)=6-6+4=4,
\]
\[
|\vec{n}_1|=\sqrt{4+1+4}=3,\qquad |\vec{n}_2|=\sqrt{9+36+4}=7.
\]
Hence $\cos\theta=\dfrac{|4|}{3\times 7}=\dfrac{4}{21}$, so
\[
\theta=\cos^{-1}\!\left(\tfrac{4}{21}\right)\approx 79.0^{\circ}.
\]
\end{exemplebox}

\courssub{The Angle Between a Line and a Plane}

\textbf{Intuition.} When a javelin sticks into the ground, the natural angle to quote is the angle between the javelin and the \emph{ground itself} — not the angle with the vertical. The vertical direction is the \emph{normal} to the ground. So the angle $\varphi$ between a line and a plane and the angle $\psi$ between the line and the normal are \textbf{complementary}: $\varphi+\psi=90^{\circ}$. This single fact drives the whole theory — and is also the source of the most common examination error.

\begin{definition}[Angle between a line and a plane]
Let $\mathcal{L}$ be a line piercing a plane $\mathcal{P}$ at a point $I$. The \cle{angle between the line and the plane} is the acute angle $\varphi$ between $\mathcal{L}$ and its \textbf{orthogonal projection} onto $\mathcal{P}$. Equivalently, if $\psi$ is the acute angle between $\mathcal{L}$ and the normal $\vec{n}$ to $\mathcal{P}$, then
\[
\varphi = 90^{\circ}-\psi.
\]
\end{definition}

\begin{popfigure}
\begin{center}
\begin{tikzpicture}[scale=1.1]
% Plane
\draw[fill=popblueL, draw=popblue, thick] (-3.2,-0.9) -- (0.8,-0.9) -- (2.8,0.9) -- (-1.2,0.9) -- cycle;
\node[popblue] at (-2.5,0.5) {$\mathcal{P}$};
% Projection of the line on the plane (dashed)
\draw[dashed, thick, popgreen] (0.5,0) -- (2.3,0.35);
\node[popgreen, below right] at (2.0,0.15) {projection};
% The line piercing the plane
\draw[very thick, popdark] (-1.1,-1.5) -- (2.1,1.7);
\draw[fill=popdark] (0.5,0.1) circle (0.05) node[below left=1pt, popdark] {$I$};
\node[popdark] at (2.25,1.75) {$\mathcal{L}$};
% Normal at I
\draw[-latex, very thick, poporange] (0.5,0.1) -- (0.5,2.1) node[right] {$\vec{n}$};
% Angle phi between line and projection
\draw[popgreen, thick] (1.5,0.14) arc[start angle=8, end angle=42, radius=1.0];
\node[popgreen] at (1.75,0.5) {$\varphi$};
% Angle psi between line and normal
\draw[poppurple, thick] (0.5,1.1) arc[start angle=90, end angle=48, radius=1.0];
\node[poppurple] at (0.82,1.28) {$\psi$};
\node[poppurple, align=center] at (-1.9,1.9) {$\varphi+\psi=90^{\circ}$};
\end{tikzpicture}
\end{center}
\legende{A line $\mathcal{L}$ piercing a plane $\mathcal{P}$ at $I$. The angle $\varphi$ between line and plane is the complement of the angle $\psi$ between the line and the normal $\vec{n}$.}
\end{popfigure}

\begin{propriete}[Formula for the angle between a line and a plane — proof examinable]
Let $\mathcal{L}$ have direction vector $\vec{b}$ and let $\mathcal{P}$ have normal vector $\vec{n}$. Then the angle $\varphi$ between $\mathcal{L}$ and $\mathcal{P}$ satisfies
\[
\sin\varphi=\frac{|\vec{b}\cdot\vec{n}|}{|\vec{b}|\,|\vec{n}|}.
\]
\textbf{Proof.} Let $\psi$ be the acute angle between $\vec{b}$ and $\vec{n}$. Then $\cos\psi=\dfrac{|\vec{b}\cdot\vec{n}|}{|\vec{b}||\vec{n}|}$. Since $\varphi=90^{\circ}-\psi$, we have $\sin\varphi=\sin(90^{\circ}-\psi)=\cos\psi$, which gives the result. $\blacksquare$
\end{propriete}

\begin{attention}[SINE, not cosine!]
For the angle between \textbf{two planes} we use $\cos\theta$ with the two \emph{normals}. For the angle between a \textbf{line and a plane} we use $\sin\varphi$ with the \emph{direction vector of the line} and the \emph{normal of the plane}. Mixing these up — writing cosine where sine is required — is the single most frequent mistake in this topic at GCE Advanced Level. Remember: \emph{line--plane $\Rightarrow$ complement $\Rightarrow$ sine}.
\end{attention}

\begin{propriete}[Line perpendicular or parallel to a plane]
Let $\mathcal{L}$ have direction $\vec{b}$ and $\mathcal{P}$ have normal $\vec{n}$. Then:
\begin{itemize}
\item $\mathcal{L}\perp\mathcal{P} \iff \vec{b}\parallel\vec{n}$ \quad (equivalently $\varphi=90^{\circ}$, i.e. $\sin\varphi=1$);
\item $\mathcal{L}\parallel\mathcal{P}$ (or $\mathcal{L}\subset\mathcal{P}$) $\iff \vec{b}\cdot\vec{n}=0$ \quad (equivalently $\varphi=0$).
\end{itemize}
\end{propriete}

\begin{attention}[Parallel to a plane is not the same as lying in the plane]
If $\vec{b}\cdot\vec{n}=0$, the line is \emph{parallel to} the plane — but it may float above it without ever touching it. To show that the line \textbf{lies entirely in} the plane you must check \textbf{two} things:
\begin{enumerate}
\item $\vec{b}\cdot\vec{n}=0$ (the direction is parallel to the plane), \textbf{and}
\item one point of the line satisfies the equation of the plane.
\end{enumerate}
Omitting the point check is a classic trap in GCE questions.
\end{attention}

\begin{methode}[Angle between a line and a plane]
\begin{enumerate}
\item Read off the direction vector $\vec{b}$ of the line (from its vector or parametric equation).
\item Read off the normal $\vec{n}$ of the plane (coefficients of $x,y,z$).
\item Compute $\vec{b}\cdot\vec{n}$, $|\vec{b}|$ and $|\vec{n}|$.
\item Apply $\sin\varphi=\dfrac{|\vec{b}\cdot\vec{n}|}{|\vec{b}||\vec{n}|}$ and take $\varphi=\sin^{-1}(\cdots)$, giving an acute angle.
\end{enumerate}
\end{methode}

\begin{exoresolu}[GCE-style: angle between a line and a plane]
Find the angle between the line $\vec{r}=\begin{pmatrix}1\\0\\2\end{pmatrix}+t\begin{pmatrix}2\\1\\-2\end{pmatrix}$ and the plane $6x-2y+3z=4$.
\tcblower
\textbf{Solution.} Direction of line: $\vec{b}=\begin{pmatrix}2\\1\\-2\end{pmatrix}$. Normal of plane: $\vec{n}=\begin{pmatrix}6\\-2\\3\end{pmatrix}$.
\[
\vec{b}\cdot\vec{n}=(2)(6)+(1)(-2)+(-2)(3)=12-2-6=4,
\]
\[
|\vec{b}|=\sqrt{4+1+4}=3,\qquad |\vec{n}|=\sqrt{36+4+9}=7.
\]
Therefore
\[
\sin\varphi=\frac{|4|}{3\times 7}=\frac{4}{21}\quad\Longrightarrow\quad \varphi=\sin^{-1}\!\left(\tfrac{4}{21}\right)\approx 11.0^{\circ}.
\]
The line meets the plane at an angle of about $11.0^{\circ}$. (Notice how the same numbers $4$, $3$, $7$ as in the earlier example now appear under a \emph{sine} — a useful self-check that the two formulas are genuinely different.)
\end{exoresolu}

\begin{exoresolu}[GCE-style: mixed angles and a containment check]
The planes $\mathcal{P}_1: x+2y+2z=0$ and $\mathcal{P}_2: y-z=3$ are given, together with the line $\mathcal{L}: \vec{r}=\begin{pmatrix}1\\1\\-1\end{pmatrix}+\lambda\begin{pmatrix}2\\-2\\1\end{pmatrix}$.
\begin{enumerate}
\item Determine whether $\mathcal{P}_1$ and $\mathcal{P}_2$ are perpendicular. If not, find the angle between them.
\item Determine whether $\mathcal{L}$ is parallel to $\mathcal{P}_1$, and whether it lies in $\mathcal{P}_1$.
\item Find the angle between $\mathcal{L}$ and $\mathcal{P}_2$.
\end{enumerate}
\tcblower
\textbf{Solution.} Normals: $\vec{n}_1=\begin{pmatrix}1\\2\\2\end{pmatrix}$, $\vec{n}_2=\begin{pmatrix}0\\1\\-1\end{pmatrix}$; direction of line: $\vec{b}=\begin{pmatrix}2\\-2\\1\end{pmatrix}$.

\textbf{1.} $\vec{n}_1\cdot\vec{n}_2=(1)(0)+(2)(1)+(2)(-1)=0+2-2=0$. Hence $\mathcal{P}_1 \perp \mathcal{P}_2$ and the angle between them is $90^{\circ}$.

\textbf{2.} $\vec{b}\cdot\vec{n}_1=(2)(1)+(-2)(2)+(1)(2)=2-4+2=0$, so $\mathcal{L}$ is parallel to $\mathcal{P}_1$. Now test the point $(1,1,-1)$ of $\mathcal{L}$ in $\mathcal{P}_1$: $1+2(1)+2(-1)=1\neq 0$. The point is \emph{not} on the plane, so $\mathcal{L}$ is parallel to $\mathcal{P}_1$ but does \textbf{not} lie in it.

\textbf{3.} $\vec{b}\cdot\vec{n}_2=(2)(0)+(-2)(1)+(1)(-1)=0-2-1=-3$; $|\vec{b}|=\sqrt{4+4+1}=3$; $|\vec{n}_2|=\sqrt{0+1+1}=\sqrt{2}$. Hence
\[
\sin\varphi=\frac{|-3|}{3\times \sqrt{2}}=\frac{1}{\sqrt{2}}\quad\Longrightarrow\quad \varphi=\sin^{-1}\!\left(\tfrac{1}{\sqrt{2}}\right)=45^{\circ}.
\]
\end{exoresolu}

\begin{aretenir}[Key points]
\begin{itemize}
\item Angle between two planes = acute angle between the normals: $\cos\theta=\dfrac{|\vec{n}_1\cdot\vec{n}_2|}{|\vec{n}_1||\vec{n}_2|}$.
\item Planes perpendicular $\iff \vec{n}_1\cdot\vec{n}_2=0$; planes parallel $\iff \vec{n}_1\parallel\vec{n}_2$.
\item Angle between a line and a plane uses the \textbf{sine} of the complementary angle: $\sin\varphi=\dfrac{|\vec{b}\cdot\vec{n}|}{|\vec{b}||\vec{n}|}$.
\item Line $\perp$ plane $\iff \vec{b}\parallel\vec{n}$; line parallel to plane $\iff \vec{b}\cdot\vec{n}=0$.
\item Line \emph{contained} in plane $\iff \vec{b}\cdot\vec{n}=0$ \textbf{and} a point of the line satisfies the plane's equation.
\end{itemize}
\end{aretenir}

\begin{exercice}[Practice]
\begin{enumerate}
\item Find the acute angle between the planes $x+y+z=1$ and $x-2y+z=4$.
\item Find the angle between the line $\dfrac{x-1}{2}=\dfrac{y+1}{-1}=\dfrac{z}{2}$ and the plane $2x-2y+z=6$. What do you notice?
\item Show that the line $\vec{r}=\begin{pmatrix}1\\0\\3\end{pmatrix}+t\begin{pmatrix}1\\1\\0\end{pmatrix}$ lies entirely in the plane $x-y+z=4$.
\item The plane $\mathcal{P}$ contains the point $A(2,1,0)$ and is perpendicular to both planes $x+y-z=3$ and $2x-y+z=1$. Explain how the normal of $\mathcal{P}$ can be found, and determine the angle between $\mathcal{P}$ and the line through $A$ with direction $(1,2,2)$.
\end{enumerate}
\end{exercice}

\begin{corrige}[Exercise 1]
\textbf{1.} $\vec{n}_1=\begin{pmatrix}1\\1\\1\end{pmatrix}$, $\vec{n}_2=\begin{pmatrix}1\\-2\\1\end{pmatrix}$. $\vec{n}_1\cdot\vec{n}_2=1-2+1=0$, so the planes are perpendicular: $\theta=90^{\circ}$.

\textbf{2.} Line direction $\vec{b}=\begin{pmatrix}2\\-1\\2\end{pmatrix}$, plane normal $\vec{n}=\begin{pmatrix}2\\-2\\1\end{pmatrix}$. $\vec{b}\cdot\vec{n}=4+2+2=8$, $|\vec{b}|=3$, $|\vec{n}|=3$, so $\sin\varphi=\tfrac{8}{9}$, $\varphi\approx 62.7^{\circ}$. Notice the numbers match part 3 of the worked exercise — different geometry, same computation.

\textbf{3.} $\vec{b}=\begin{pmatrix}1\\1\\0\end{pmatrix}$, $\vec{n}=\begin{pmatrix}1\\-1\\1\end{pmatrix}$. $\vec{b}\cdot\vec{n}=1-1+0=0$, so the line is parallel to the plane. Point $(1,0,3)$ on the line: $1-0+3=4$, satisfies the plane equation. Hence the line lies entirely in the plane.

\textbf{4.} The normal $\vec{n}$ of $\mathcal{P}$ is perpendicular to both given normals $\vec{n}_1=\begin{pmatrix}1\\1\\-1\end{pmatrix}$ and $\vec{n}_2=\begin{pmatrix}2\\-1\\1\end{pmatrix}$, so take $\vec{n}=\vec{n}_1\times\vec{n}_2=\begin{vmatrix}\mathbf{i}&\mathbf{j}&\mathbf{k}\\1&1&-1\\2&-1&1\end{vmatrix}=\mathbf{i}(1-1)-\mathbf{j}(1+2)+\mathbf{k}(-1-2)=\begin{pmatrix}0\\-3\\-3\end{pmatrix}$, i.e. direction $\begin{pmatrix}0\\1\\1\end{pmatrix}$. Then for $\vec{b}=\begin{pmatrix}1\\2\\2\end{pmatrix}$: $\vec{b}\cdot\vec{n}=0+2+2=4$, $|\vec{b}|=3$, $|\vec{n}|=\sqrt{2}$, so $\sin\varphi=\dfrac{4}{3\sqrt{2}}=\dfrac{2\sqrt{2}}{3}$, giving $\varphi\approx 70.5^{\circ}$.
\end{corrige}

\begin{savaistu}[Dihedral angles around us]
Dihedral angles are everywhere in engineering and nature. The wings of many aircraft are tilted upward at a small dihedral angle (typically a few degrees) to improve stability in flight. Closer to home, the pitch of a corrugated-iron roof on a house in Bamenda or Buea is chosen so that the dihedral angle between the two slopes lets heavy tropical rain run off quickly while keeping the structure cheap to build. The mathematics you have just used — normals and scalar products — is exactly what architects and engineers use to compute such angles before a single sheet of zinc is cut.
\end{savaistu}

\coursec{Intersection of a Line with a Plane and of Two Planes}

In the previous section we learned how to measure the \emph{angle} between two planes and between a line and a plane. Those questions already told us something important: a line and a plane are rarely indifferent to each other. Either they meet, or they run alongside each other forever. In this section we go one step further and answer the question \emph{where} exactly they meet, and we do the same for two planes. We then study the family of all planes passing through the line of intersection of two given planes — a classic GCE Advanced Level topic. These problems reduce to solving one linear equation (for a line and a plane) or a small system of two linear equations (for two planes), followed by a clear geometric interpretation.

\courssub{Point of Intersection of a Line and a Plane}

\textbf{The guiding idea.}\par
Imagine a javelin thrown towards the ground at the Buea stadium. The javelin travels along a straight line; the ground is (approximately) a plane. To find where the javelin lands, we simply ask: \emph{at what instant does a point of the trajectory satisfy the equation of the ground?} Algebraically, a point lies on the line for some value of the parameter $\lambda$, and it lies on the plane if its position vector satisfies the plane equation. The meeting point is the point that does \emph{both}.

\begin{methode}[Line--plane intersection: substitute and solve]
To find where the line $\vec{r} = \vec{a} + \lambda \vec{b}$ meets the plane $\vec{r}\cdot\vec{n} = d$:
\begin{enumerate}
\item Substitute $\vec{r} = \vec{a} + \lambda\vec{b}$ into the plane equation:
\[ (\vec{a} + \lambda\vec{b})\cdot\vec{n} = d. \]
\item Expand the scalar product: $\vec{a}\cdot\vec{n} + \lambda\,(\vec{b}\cdot\vec{n}) = d$.
\item Solve for $\lambda$: $\displaystyle \lambda = \frac{d - \vec{a}\cdot\vec{n}}{\vec{b}\cdot\vec{n}}$, provided $\vec{b}\cdot\vec{n} \neq 0$.
\item Substitute this value of $\lambda$ back into $\vec{r} = \vec{a} + \lambda\vec{b}$ to obtain the position vector of the point of intersection.
\end{enumerate}
\end{methode}

\begin{propriete}[The three possible cases]
Let the line have equation $\vec{r} = \vec{a} + \lambda\vec{b}$ and the plane $\vec{r}\cdot\vec{n} = d$. Substitution gives
\[ \lambda\,(\vec{b}\cdot\vec{n}) = d - \vec{a}\cdot\vec{n}. \]
\begin{itemize}
\item \textbf{Unique point of intersection:} if $\vec{b}\cdot\vec{n} \neq 0$ (the line is not parallel to the plane), there is exactly one value of $\lambda$, hence exactly one meeting point.
\item \textbf{Line parallel to the plane, no intersection:} if $\vec{b}\cdot\vec{n} = 0$ and $\vec{a}\cdot\vec{n} \neq d$ (the point $\vec{a}$ of the line is not on the plane), the equation becomes $0 = \text{(non-zero)}$, impossible: the line never meets the plane.
\item \textbf{Line contained in the plane:} if $\vec{b}\cdot\vec{n} = 0$ and $\vec{a}\cdot\vec{n} = d$, the equation becomes $0 = 0$, true for \emph{every} $\lambda$: every point of the line lies in the plane.
\end{itemize}
This property is examinable and its proof is simply the discussion of the linear equation above.
\end{propriete}

\begin{popfigure}
\begin{center}
\begin{tikzpicture}[scale=1.0]
% plane as parallelogram
\fill[popblueL, opacity=0.6] (-3,-1.2) -- (1.2,-1.2) -- (3.4,0.6) -- (-0.8,0.6) -- cycle;
\draw[popblue] (-3,-1.2) -- (1.2,-1.2) -- (3.4,0.6) -- (-0.8,0.6) -- cycle;
\node[popblue] at (2.6,0.1) {$\vec{r}\cdot\vec{n}=d$};
% line
\draw[very thick, poporange] (-1.6,2.6) -- (0.9,-1.05);
\fill[popink] (0,-0.35) circle (2.2pt);
\node[popink, right] at (0.05,-0.35) {$P$};
% points along line with lambda
\fill[poporange] (-1.2,1.95) circle (1.6pt);
\node[poporange, left] at (-1.25,1.95) {$\lambda=0\;(\vec{a})$};
\fill[poporange] (-0.6,1.3) circle (1.6pt);
\node[poporange, left] at (-0.65,1.3) {$\lambda=1$};
\fill[poporange] (-0.3,0.475) circle (1.6pt);
\node[poporange, left] at (-0.42,0.5) {$\lambda=2$};
\node[popink, right, align=left, text width=3.4cm] at (1.1,1.4) {unique $\lambda$ gives the meeting point $P$};
% normal
\draw[->, thick, poppurple] (0,-0.35) -- (0.7,0.55);
\node[poppurple, right] at (0.7,0.45) {$\vec{n}$};
\end{tikzpicture}
\end{center}
\legende{A line $\vec{r}=\vec{a}+\lambda\vec{b}$ piercing a plane at a single point $P$; each point of the line corresponds to one value of $\lambda$.}
\end{popfigure}

\begin{exoresolu}[Finding where a line meets a plane]
Find the point of intersection of the line
\[ \vec{r} = \begin{pmatrix} 1 \\ 0 \\ 2 \end{pmatrix} + \lambda \begin{pmatrix} 2 \\ 1 \\ -1 \end{pmatrix} \quad\text{with the plane}\quad \vec{r}\cdot\begin{pmatrix} 1 \\ 2 \\ 1 \end{pmatrix} = 7. \]
\tcblower
\textbf{Solution.} Substitute $\vec{r} = \vec{a} + \lambda\vec{b}$ into the plane equation:
\begin{align*}
\left[\begin{pmatrix} 1 \\ 0 \\ 2 \end{pmatrix} + \lambda \begin{pmatrix} 2 \\ 1 \\ -1 \end{pmatrix}\right]\cdot\begin{pmatrix} 1 \\ 2 \\ 1 \end{pmatrix} &= 7\\
(1 + 0 + 2) + \lambda\,(2 + 2 - 1) &= 7\\
3 + 3\lambda &= 7 \quad\Longrightarrow\quad \lambda = \tfrac{4}{3}.
\end{align*}
Substituting back:
\[ \vec{r} = \begin{pmatrix} 1 \\ 0 \\ 2 \end{pmatrix} + \tfrac{4}{3}\begin{pmatrix} 2 \\ 1 \\ -1 \end{pmatrix} = \begin{pmatrix} 11/3 \\ 4/3 \\ 2/3 \end{pmatrix}. \]
The point of intersection is $\left(\tfrac{11}{3}, \tfrac{4}{3}, \tfrac{2}{3}\right)$. \textbf{Check:} $\tfrac{11}{3} + 2\cdot\tfrac{4}{3} + \tfrac{2}{3} = \tfrac{11+8+2}{3} = 7$. \checkmark
\end{exoresolu}

\begin{exemplebox}[The parallel and contained cases]
Consider the plane $\vec{r}\cdot\begin{pmatrix} 1 \\ 1 \\ 1 \end{pmatrix} = 6$.
\begin{itemize}
\item Line $\vec{r} = \begin{pmatrix} 0 \\ 0 \\ 1 \end{pmatrix} + \lambda\begin{pmatrix} 1 \\ -1 \\ 0 \end{pmatrix}$: here $\vec{b}\cdot\vec{n} = 1 - 1 + 0 = 0$ and $\vec{a}\cdot\vec{n} = 1 \neq 6$. The line is parallel to the plane and never meets it.
\item Line $\vec{r} = \begin{pmatrix} 1 \\ 2 \\ 3 \end{pmatrix} + \lambda\begin{pmatrix} 1 \\ -1 \\ 0 \end{pmatrix}$: again $\vec{b}\cdot\vec{n} = 0$, but now $\vec{a}\cdot\vec{n} = 1+2+3 = 6$. The whole line lies \emph{in} the plane.
\end{itemize}
\end{exemplebox}

\begin{attention}[Do not conclude ``no intersection'' too quickly]
When you find $\vec{b}\cdot\vec{n} = 0$, the work is \emph{not} finished. You must still compute $\vec{a}\cdot\vec{n}$ and compare it with $d$. If $\vec{a}\cdot\vec{n} = d$, the line is contained in the plane (infinitely many common points); only if $\vec{a}\cdot\vec{n} \neq d$ may you write ``the line is parallel to the plane and does not intersect it''. Examiners regularly award a mark specifically for this check.
\end{attention}

\courssub{Intersection of Two Planes}

\textbf{The guiding idea.}\par
Open a book: the two pages meet along the spine. Two non-parallel planes in space behave exactly like the two pages — they meet along a whole straight line. Finding this \cle{line of intersection} means finding all points $(x,y,z)$ that satisfy \emph{both} Cartesian equations at once.

\begin{propriete}[Relative positions of two planes]
Let the planes be $\Pi_1: \vec{r}\cdot\vec{n}_1 = d_1$ and $\Pi_2: \vec{r}\cdot\vec{n}_2 = d_2$.
\begin{itemize}
\item \textbf{Intersecting in a line:} $\vec{n}_1$ and $\vec{n}_2$ are not parallel ($\vec{n}_1 \times \vec{n}_2 \neq \vec{0}$). The intersection is a straight line.
\item \textbf{Parallel and distinct:} $\vec{n}_2 = k\,\vec{n}_1$ for some scalar $k$, but the full equations are not proportional (i.e. $d_2 \neq k\,d_1$). No common point.
\item \textbf{Coincident:} the two complete equations are proportional: $\vec{n}_2 = k\,\vec{n}_1$ \emph{and} $d_2 = k\,d_1$. The two planes are one and the same plane.
\end{itemize}
In Cartesian form, with $a_1x+b_1y+c_1z=d_1$ and $a_2x+b_2y+c_2z=d_2$: the planes are parallel iff $\dfrac{a_1}{a_2}=\dfrac{b_1}{b_2}=\dfrac{c_1}{c_2}$, and coincident iff this common ratio also equals $\dfrac{d_1}{d_2}$.
\end{propriete}

\begin{propriete}[Direction of the line of intersection]
If two planes with normals $\vec{n}_1$ and $\vec{n}_2$ intersect, their line of intersection has direction vector
\[ \vec{b} = \vec{n}_1 \times \vec{n}_2. \]
\textbf{Justification (examinable).} The line lies in $\Pi_1$, so its direction $\vec{b}$ is perpendicular to $\vec{n}_1$. The line also lies in $\Pi_2$, so $\vec{b}$ is perpendicular to $\vec{n}_2$. The cross product $\vec{n}_1 \times \vec{n}_2$ is precisely a vector perpendicular to both $\vec{n}_1$ and $\vec{n}_2$, so we may take $\vec{b} = \vec{n}_1 \times \vec{n}_2$.
\end{propriete}

\begin{attention}[Cross product: direction, not normal]
A very common error is to say that $\vec{n}_1 \times \vec{n}_2$ is ``perpendicular to the line of intersection''. It is the opposite: $\vec{n}_1 \times \vec{n}_2$ \emph{is} the direction vector of the line of intersection — it points \emph{along} the line, not across it. Remember: each normal is perpendicular to the line, and the cross product of the two normals is parallel to the line.
\end{attention}

\begin{methode}[Equation of the line of intersection of two planes]
Given $\Pi_1: \vec{r}\cdot\vec{n}_1 = d_1$ and $\Pi_2: \vec{r}\cdot\vec{n}_2 = d_2$ with non-parallel normals:
\begin{enumerate}
\item Compute the direction $\vec{b} = \vec{n}_1 \times \vec{n}_2$.
\item Find \emph{one} point lying in both planes: set one coordinate to a convenient value (usually $z = 0$; if the resulting pair of equations is inconsistent, try $y = 0$ or $x = 0$) and solve the remaining $2\times 2$ system.
\item Write the line in vector form $\vec{r} = \vec{a} + \lambda\vec{b}$, where $\vec{a}$ is the position vector of the point found.
\end{enumerate}
\end{methode}

\begin{popfigure}
\begin{center}
\begin{tikzpicture}[scale=1.0]
% plane 1
\fill[popblueL, opacity=0.55] (-3,-1.1) -- (1.0,-1.1) -- (2.6,0.9) -- (-1.4,0.9) -- cycle;
\draw[popblue] (-3,-1.1) -- (1.0,-1.1) -- (2.6,0.9) -- (-1.4,0.9) -- cycle;
\node[popblue] at (-2.3,0.5) {$\Pi_1$};
% plane 2 (tilted)
\fill[popgreenL, opacity=0.55] (-1.2,-1.6) -- (2.2,-1.6) -- (0.6,1.6) -- (-2.8,1.6) -- cycle;
\draw[popgreen] (-1.2,-1.6) -- (2.2,-1.6) -- (0.6,1.6) -- (-2.8,1.6) -- cycle;
\node[popgreen] at (-2.2,1.25) {$\Pi_2$};
% line of intersection
\draw[very thick, poporange] (-2.6,-0.32) -- (2.2,0.62);
\node[poporange, right] at (2.2,0.62) {$L:\ \vec{r}=\vec{a}+\lambda\vec{b}$};
% direction vector b on the line
\draw[->, very thick, poporange] (0.4,0.08) -- (1.6,0.39);
\node[poporange, above] at (1.0,0.28) {$\vec{b}=\vec{n}_1\times\vec{n}_2$};
% normals
\draw[->, thick, popblue] (-0.8,-0.13) -- (-0.8,1.0);
\node[popblue, left] at (-0.85,0.9) {$\vec{n}_1$};
\draw[->, thick, popgreen] (-0.2,-0.05) -- (0.85,-0.75);
\node[popgreen, right] at (0.8,-0.72) {$\vec{n}_2$};
\end{tikzpicture}
\end{center}
\legende{Two planes $\Pi_1$ and $\Pi_2$ intersecting along the line $L$, whose direction is $\vec{b}=\vec{n}_1\times\vec{n}_2$, perpendicular to both normals.}
\end{popfigure}

\begin{exoresolu}[Line of intersection of two planes]
Find, in vector form, the equation of the line of intersection of the planes
\[ \Pi_1:\ x + 2y - z = 3 \qquad\text{and}\qquad \Pi_2:\ 2x - y + z = 4. \]
\tcblower
\textbf{Solution.} The normals are $\vec{n}_1 = \begin{pmatrix} 1 \\ 2 \\ -1 \end{pmatrix}$ and $\vec{n}_2 = \begin{pmatrix} 2 \\ -1 \\ 1 \end{pmatrix}$, which are not proportional, so the planes intersect in a line.

\textbf{Step 1: direction.}
\[ \vec{b} = \vec{n}_1 \times \vec{n}_2 = \begin{pmatrix} (2)(1) - (-1)(-1) \\ (-1)(2) - (1)(1) \\ (1)(-1) - (2)(2) \end{pmatrix} = \begin{pmatrix} 2 - 1 \\ -2 - 1 \\ -1 - 4 \end{pmatrix} = \begin{pmatrix} 1 \\ -3 \\ -5 \end{pmatrix}. \]

\textbf{Step 2: one common point.} Set $z = 0$ in both equations:
\[ \begin{cases} x + 2y = 3 \\ 2x - y = 4 \end{cases} \]
From the second equation, $y = 2x - 4$; substituting: $x + 2(2x-4) = 3 \Rightarrow 5x = 11 \Rightarrow x = \tfrac{11}{5}$, $y = \tfrac{22}{5} - 4 = \tfrac{2}{5}$. So $\vec{a} = \begin{pmatrix} 11/5 \\ 2/5 \\ 0 \end{pmatrix}$.

\textbf{Step 3: vector equation.}
\[ \vec{r} = \begin{pmatrix} 11/5 \\ 2/5 \\ 0 \end{pmatrix} + \lambda \begin{pmatrix} 1 \\ -3 \\ -5 \end{pmatrix}. \]
\textbf{Check:} $\tfrac{11}{5} + 2\cdot\tfrac{2}{5} - 0 = \tfrac{15}{5} = 3$ \checkmark\ and $2\cdot\tfrac{11}{5} - \tfrac{2}{5} + 0 = \tfrac{20}{5} = 4$ \checkmark.
\end{exoresolu}

\begin{exemplebox}[Recognising parallel and coincident planes]
\begin{itemize}
\item $x - 2y + 3z = 5$ and $2x - 4y + 6z = 7$: the coefficients of $x,y,z$ are in ratio $1:2$, but $5 \times 2 = 10 \neq 7$. The planes are \textbf{parallel and distinct}.
\item $x - 2y + 3z = 5$ and $2x - 4y + 6z = 10$: every coefficient, including the constant, is doubled. The planes are \textbf{coincident} — the second equation carries no new information.
\end{itemize}
\end{exemplebox}

\begin{exoresolu}[A complete discussion]
Show that the line $\vec{r} = \begin{pmatrix} 2 \\ 1 \\ 0 \end{pmatrix} + \lambda\begin{pmatrix} 1 \\ -3 \\ -5 \end{pmatrix}$ lies entirely in the plane $x + 2y - z = 4$, and state its relation with the plane $2x - y + z = 1$.
\tcblower
\textbf{Solution.} Compute $\vec{b}\cdot\vec{n} = (1)(1) + (-3)(2) + (-5)(-1) = 1 - 6 + 5 = 0$, so the line is parallel to the plane. Then $\vec{a}\cdot\vec{n} = 2 + 2(1) - 0 = 4 = d$: the point $(2,1,0)$ lies in the plane. Hence the line is \textbf{contained} in the plane $x + 2y - z = 4$.

For the second plane, $\vec{b}\cdot\vec{n}_2 = 2 + 3 - 5 = 0$ and $\vec{a}\cdot\vec{n}_2 = 4 - 1 + 0 = 3 \neq 1$: the line is \textbf{parallel to the plane but does not meet it}.
\end{exoresolu}

\courssub{Equation of a Plane Through the Line of Intersection of Two Planes (Pencil of Planes)}

\textbf{The guiding idea.}\par
When two planes $\Pi_1$ and $\Pi_2$ intersect, they form a line $L$. There are infinitely many planes that contain this line $L$; they form a \cle{pencil of planes} (or \cle{sheaf of planes}). Any plane in this pencil can be described by a linear combination of the equations of $\Pi_1$ and $\Pi_2$. This is a powerful tool for solving problems like ``Find the plane through the line of intersection of $\Pi_1$ and $\Pi_2$ that also passes through a given point $P$'' or ``that is perpendicular to a given plane $\Pi_3$''.

\begin{propriete}[Equation of the pencil of planes]
Let $\Pi_1: \vec{r}\cdot\vec{n}_1 = d_1$ and $\Pi_2: \vec{r}\cdot\vec{n}_2 = d_2$ be two intersecting planes. The equation of any plane $\Pi$ containing their line of intersection $L$ is
\[ \vec{r}\cdot\vec{n}_1 - d_1 + \lambda\,(\vec{r}\cdot\vec{n}_2 - d_2) = 0 \quad\text{or equivalently}\quad \vec{r}\cdot(\vec{n}_1 + \lambda\vec{n}_2) = d_1 + \lambda d_2, \]
where $\lambda \in \mathbb{R}$ is a parameter.
In Cartesian form, if $\Pi_1: a_1x+b_1y+c_1z=d_1$ and $\Pi_2: a_2x+b_2y+c_2z=d_2$, the pencil is
\[ (a_1x+b_1y+c_1z-d_1) + \lambda\,(a_2x+b_2y+c_2z-d_2) = 0. \]
\textbf{Note:} The plane $\Pi_2$ itself corresponds to $\lambda \to \infty$ (or we can write the symmetric form $\mu(\Pi_1) + \lambda(\Pi_2) = 0$). The single parameter $\lambda$ covers all planes except $\Pi_2$; to include $\Pi_2$, use $\mu\Pi_1 + \lambda\Pi_2 = 0$ with $(\mu,\lambda) \neq (0,0)$.
\end{propriete}

\begin{methode}[Finding a particular plane in the pencil]
Given two intersecting planes $\Pi_1$ and $\Pi_2$ and an extra condition (point, perpendicularity, parallelism to a line, etc.):
\begin{enumerate}
\item Write the general equation of the pencil: $\Pi_1 + \lambda\Pi_2 = 0$ (in Cartesian or vector form).
\item Impose the extra condition to obtain an equation for $\lambda$.
\item Solve for $\lambda$ and substitute back to get the required plane.
\end{enumerate}
\end{methode}

\begin{exoresolu}[Plane through the line of intersection and a given point]
Find the equation of the plane that passes through the line of intersection of
\[ \Pi_1:\ x + y + z = 6 \quad\text{and}\quad \Pi_2:\ 2x - y + z = 3 \]
and through the point $P(1,2,3)$.
\tcblower
\textbf{Solution.}
The pencil of planes through the intersection line is:
\[ (x + y + z - 6) + \lambda\,(2x - y + z - 3) = 0. \]
Substitute $P(1,2,3)$:
\[ (1+2+3-6) + \lambda\,(2-2+3-3) = 0 \;\Longrightarrow\; 0 + \lambda\cdot 0 = 0. \]
This is satisfied for \emph{any} $\lambda$, which means $P$ lies on the intersection line itself (check: $1+2+3=6$ and $2-2+3=3$). The problem is degenerate: infinitely many planes through the line contain $P$. 

\textbf{Modified example:} Find the plane through the same line and the point $Q(0,0,0)$.
Substitute $Q$:
\[ (0+0+0-6) + \lambda\,(0-0+0-3) = 0 \;\Longrightarrow\; -6 - 3\lambda = 0 \;\Longrightarrow\; \lambda = -2. \]
The required plane is:
\[ (x+y+z-6) - 2(2x-y+z-3) = 0 \;\Longrightarrow\; x+y+z-6 -4x+2y-2z+6 = 0 \;\Longrightarrow\; -3x + 3y - z = 0, \]
or $\mathbf{3x - 3y + z = 0}$.
\end{exoresolu}

\begin{exoresolu}[Plane through the line of intersection, perpendicular to a third plane]
Find the equation of the plane through the line of intersection of
\[ \Pi_1:\ x + 2y - z = 4 \quad\text{and}\quad \Pi_2:\ 2x - y + 3z = 5 \]
that is perpendicular to the plane $\Pi_3:\ 3x + y - 2z = 1$.
\tcblower
\textbf{Solution.}
The pencil is: $(x + 2y - z - 4) + \lambda\,(2x - y + 3z - 5) = 0$.
Collect terms:
\[ (1+2\lambda)x + (2-\lambda)y + (-1+3\lambda)z + (-4-5\lambda) = 0. \]
The normal of this plane is $\vec{n} = \begin{pmatrix} 1+2\lambda \\ 2-\lambda \\ -1+3\lambda \end{pmatrix}$.
Perpendicular to $\Pi_3$ means $\vec{n} \cdot \vec{n}_3 = 0$, where $\vec{n}_3 = \begin{pmatrix} 3 \\ 1 \\ -2 \end{pmatrix}$.
\[ 3(1+2\lambda) + 1(2-\lambda) - 2(-1+3\lambda) = 0 \]
\[ 3 + 6\lambda + 2 - \lambda + 2 - 6\lambda = 0 \;\Longrightarrow\; 7 - \lambda = 0 \;\Longrightarrow\; \lambda = 7. \]
Substitute $\lambda = 7$:
\[ (1+14)x + (2-7)y + (-1+21)z + (-4-35) = 0 \]
\[ 15x - 5y + 20z - 39 = 0. \]
The required plane is $\mathbf{15x - 5y + 20z = 39}$.
\end{exoresolu}

\begin{attention}[The parameter $\lambda$ and the second plane]
The form $\Pi_1 + \lambda\Pi_2 = 0$ generates all planes through the intersection line \emph{except} $\Pi_2$ itself (which would require $\lambda \to \infty$). If the answer turns out to be $\Pi_2$, you will not find a finite $\lambda$. Always check at the end whether $\Pi_2$ itself satisfies the extra condition. Alternatively, use the symmetric form $\mu\Pi_1 + \lambda\Pi_2 = 0$ and treat the ratio $\mu:\lambda$ as the unknown.
\end{attention}

\courssub{Synthesis Problems}

Many GCE questions combine the concepts of this chapter: finding a line of intersection, then a plane through that line, then the intersection of that new plane with another line, etc. The key is to proceed step by step, clearly labelling each object.

\begin{exoresolu}[Synthesis: line of intersection, then plane through it, then intersection with a line]
The planes $\Pi_1: x - y + z = 2$ and $\Pi_2: 2x + y - z = 1$ intersect in a line $L$.
\begin{enumerate}
\item Find a vector equation of $L$.
\item Find the equation of the plane $\Pi$ that contains $L$ and is parallel to the vector $\vec{v} = \begin{pmatrix} 1 \\ 1 \\ 1 \end{pmatrix}$.
\item Find the point where the line $\vec{r} = \begin{pmatrix} 0 \\ 0 \\ 0 \end{pmatrix} + \mu\begin{pmatrix} 1 \\ 2 \\ 3 \end{pmatrix}$ meets $\Pi$.
\end{enumerate}
\tcblower
\textbf{Solution.}
\textbf{(i) Line $L$:} Normals $\vec{n}_1 = (1,-1,1)$, $\vec{n}_2 = (2,1,-1)$. Direction $\vec{b} = \vec{n}_1 \times \vec{n}_2 = \begin{pmatrix} 0 \\ 3 \\ 3 \end{pmatrix} \parallel \begin{pmatrix} 0 \\ 1 \\ 1 \end{pmatrix}$.
Point on $L$: set $z=0 \Rightarrow x-y=2,\; 2x+y=1 \Rightarrow 3x=3 \Rightarrow x=1,\; y=-1$. So $\vec{a} = (1,-1,0)$.
$L:\ \vec{r} = \begin{pmatrix} 1 \\ -1 \\ 0 \end{pmatrix} + \lambda\begin{pmatrix} 0 \\ 1 \\ 1 \end{pmatrix}$.

\textbf{(ii) Plane $\Pi$ through $L$ parallel to $\vec{v}$:} Since $\Pi$ contains $L$, its normal $\vec{n}$ is perpendicular to $\vec{b}$. Since $\Pi$ is parallel to $\vec{v}$, $\vec{n}$ is also perpendicular to $\vec{v}$. Thus $\vec{n} = \vec{b} \times \vec{v} = \begin{pmatrix} 0 \\ 1 \\ 1 \end{pmatrix} \times \begin{pmatrix} 1 \\ 1 \\ 1 \end{pmatrix} = \begin{pmatrix} 0 \\ 1 \\ -1 \end{pmatrix}$.
Using point $(1,-1,0)$ on $L$: $\vec{r}\cdot\begin{pmatrix} 0 \\ 1 \\ -1 \end{pmatrix} = (1,-1,0)\cdot(0,1,-1) = -1$.
Equation: $y - z = -1$ or $\mathbf{y - z + 1 = 0}$.
(Alternative using pencil: $\Pi_1 + \lambda\Pi_2 = 0 \Rightarrow (1+2\lambda)x + (-1+\lambda)y + (1-\lambda)z = 2+\lambda$. Normal $\vec{n} = (1+2\lambda, -1+\lambda, 1-\lambda)$. Condition $\vec{n}\cdot\vec{v}=0 \Rightarrow (1+2\lambda)+(-1+\lambda)+(1-\lambda)=0 \Rightarrow 1+2\lambda=0 \Rightarrow \lambda=-1/2$. Substitute: $0x -3/2 y + 3/2 z = 3/2 \Rightarrow -y+z=1 \Rightarrow y-z=-1$. Same result.)

\textbf{(iii) Intersection with line:} Line: $\vec{r} = \mu(1,2,3)$. Substitute into $y-z+1=0$: $2\mu - 3\mu + 1 = 0 \Rightarrow -\mu + 1 = 0 \Rightarrow \mu = 1$.
Point: $(1,2,3)$.
\end{exoresolu}

\begin{aretenir}[Key points of this chapter]
\begin{itemize}
\item Line $\vec{r}=\vec{a}+\lambda\vec{b}$ meets plane $\vec{r}\cdot\vec{n}=d$: solve $(\vec{a}+\lambda\vec{b})\cdot\vec{n}=d$, giving $\lambda = \dfrac{d-\vec{a}\cdot\vec{n}}{\vec{b}\cdot\vec{n}}$.
\item If $\vec{b}\cdot\vec{n}=0$: check $\vec{a}\cdot\vec{n}$. Equal to $d$ $\Rightarrow$ line in the plane; different from $d$ $\Rightarrow$ parallel, no intersection.
\item Two planes intersect in a line iff their normals are not parallel; the line has direction $\vec{b}=\vec{n}_1\times\vec{n}_2$.
\item To find a point on that line, set one coordinate (e.g. $z=0$) and solve the resulting $2\times 2$ system.
\item Planes are parallel iff their normals are proportional; coincident iff the \emph{whole} equations (constants included) are proportional.
\item \cle{Pencil of planes}: All planes through the intersection of $\Pi_1$ and $\Pi_2$ have equation $\Pi_1 + \lambda\Pi_2 = 0$ (or $\mu\Pi_1 + \lambda\Pi_2 = 0$). Use the extra condition to find $\lambda$.
\end{itemize}
\end{aretenir}

\begin{exercice}[Practice]
\begin{enumerate}
\item Find the point where the line $\vec{r} = \begin{pmatrix} 3 \\ -1 \\ 2 \end{pmatrix} + \lambda\begin{pmatrix} 1 \\ 2 \\ -2 \end{pmatrix}$ meets the plane $\vec{r}\cdot\begin{pmatrix} 2 \\ -1 \\ 3 \end{pmatrix} = 11$.
\item Determine whether the line $\vec{r} = \begin{pmatrix} 1 \\ 1 \\ 1 \end{pmatrix} + \lambda\begin{pmatrix} 2 \\ 1 \\ -1 \end{pmatrix}$ intersects the plane $x - y + 3z = 3$, is parallel to it, or lies in it.
\item Find the vector equation of the line of intersection of $x + y + z = 6$ and $2x - y + z = 3$.
\item For what value of $k$ are the planes $x + 2y - z = 4$ and $3x + 6y - 3z = k$ coincident? What happens for other values of $k$?
\item Find the equation of the plane that passes through the line of intersection of $x - y + 2z = 3$ and $2x + y - z = 4$ and through the point $(1,0,-1)$.
\item Find the equation of the plane through the line of intersection of $x + 2y - z = 5$ and $3x - y + 2z = 1$ that is perpendicular to the plane $x - y + z = 0$.
\end{enumerate}
\end{exercice}

\begin{corrige}[Exercises]
\textbf{1.} Line: $\vec{a}=(3,-1,2)$, $\vec{b}=(1,2,-2)$. Plane: $\vec{n}=(2,-1,3)$, $d=11$.
$\vec{a}\cdot\vec{n} = 6 + 1 + 6 = 13$. $\vec{b}\cdot\vec{n} = 2 - 2 - 6 = -6$.
Equation: $13 - 6\lambda = 11 \Rightarrow -6\lambda = -2 \Rightarrow \lambda = \tfrac{1}{3}$.
Point: $\begin{pmatrix} 3 \\ -1 \\ 2 \end{pmatrix} + \tfrac{1}{3}\begin{pmatrix} 1 \\ 2 \\ -2 \end{pmatrix} = \begin{pmatrix} 10/3 \\ -1/3 \\ 4/3 \end{pmatrix}$.
\textbf{Answer:} $\left(\tfrac{10}{3}, -\tfrac{1}{3}, \tfrac{4}{3}\right)$.

\textbf{2.} Line: $\vec{a}=(1,1,1)$, $\vec{b}=(2,1,-1)$. Plane: $x-y+3z=3 \Rightarrow \vec{n}=(1,-1,3)$, $d=3$.
$\vec{b}\cdot\vec{n} = 2 - 1 - 3 = -2 \neq 0$ $\Rightarrow$ unique intersection.
$\vec{a}\cdot\vec{n} = 1 - 1 + 3 = 3 = d$. So $\lambda = 0$.
The line meets the plane at the point $\mathbf{(1,1,1)}$ (which is the point $\vec{a}$ on the line).

\textbf{3.} $\Pi_1: x+y+z=6$, $\Pi_2: 2x-y+z=3$. Normals $\vec{n}_1=(1,1,1)$, $\vec{n}_2=(2,-1,1)$.
Direction $\vec{b} = \vec{n}_1 \times \vec{n}_2 = \begin{pmatrix} 1\cdot1 - 1\cdot(-1) \\ 1\cdot2 - 1\cdot1 \\ 1\cdot(-1) - 1\cdot2 \end{pmatrix} = \begin{pmatrix} 2 \\ 1 \\ -3 \end{pmatrix}$.
Point: set $z=0 \Rightarrow x+y=6,\; 2x-y=3 \Rightarrow 3x=9 \Rightarrow x=3,\; y=3$.
Line: $\vec{r} = \begin{pmatrix} 3 \\ 3 \\ 0 \end{pmatrix} + \lambda\begin{pmatrix} 2 \\ 1 \\ -3 \end{pmatrix}$.

\textbf{4.} Planes: $x+2y-z=4$ and $3x+6y-3z=k$. Normals are proportional ($1:3$).
Coincident iff constants are in same ratio: $k = 3 \times 4 = \mathbf{12}$.
For $k \neq 12$, the planes are \textbf{parallel and distinct} (no intersection).

\textbf{5.} Pencil: $(x-y+2z-3) + \lambda(2x+y-z-4) = 0$.
Point $(1,0,-1)$: $(1-0-2-3) + \lambda(2+0+1-4) = -4 + \lambda(-1) = 0 \Rightarrow \lambda = -4$.
Plane: $(x-y+2z-3) - 4(2x+y-z-4) = 0 \Rightarrow x-y+2z-3 -8x-4y+4z+16 = 0 \Rightarrow -7x -5y +6z +13 = 0$.
\textbf{Answer:} $\mathbf{7x + 5y - 6z = 13}$.

\textbf{6.} Pencil: $(x+2y-z-5) + \lambda(3x-y+2z-1) = 0$.
Normal $\vec{n} = (1+3\lambda,\; 2-\lambda,\; -1+2\lambda)$.
Perpendicular to $x-y+z=0 \Rightarrow \vec{n}\cdot(1,-1,1)=0$:
$(1+3\lambda) - (2-\lambda) + (-1+2\lambda) = 0 \Rightarrow 1+3\lambda -2 +\lambda -1 +2\lambda = 0 \Rightarrow 6\lambda -2 = 0 \Rightarrow \lambda = \tfrac{1}{3}$.
Plane: $(x+2y-z-5) + \tfrac{1}{3}(3x-y+2z-1) = 0$. Multiply by 3:
$3x+6y-3z-15 + 3x-y+2z-1 = 0 \Rightarrow 6x + 5y - z - 16 = 0$.
\textbf{Answer:} $\mathbf{6x + 5y - z = 16}$.
\end{corrige}

\begin{savaistu}[Why this matters beyond the examination room]
When engineers design the junction of two roof panels of a building, or when a GPS system computes where a satellite signal path crosses a layer of the atmosphere, the mathematics is exactly this: intersecting lines with planes and planes with planes. Computer graphics software renders every object you see on a screen by computing millions of line--plane intersections per second — the very formula $\lambda = \dfrac{d - \vec{a}\cdot\vec{n}}{\vec{b}\cdot\vec{n}}$ that you have just mastered. The pencil of planes is used in computer vision to reconstruct 3D scenes from multiple camera views (the \emph{epipolar geometry}).
\end{savaistu}

\coursec{Pencils of Planes and Synthesis Problems}

In the previous section we learnt how to find the line of intersection of two planes $\Pi_1$ and $\Pi_2$. A natural question now arises: \emph{how many planes can pass through that same line?} Imagine a door swinging on its hinges: the hinge line stays fixed while the door occupies infinitely many positions. Each position of the door is a plane through the hinge line. This family of planes is called a \cle{pencil of planes}, and it gives us one of the most powerful time-saving tools of the whole chapter. We then finish the chapter with synthesis problems in the style of the GCE Advanced Level Paper, and a revision grid to help you choose the right tool for each question.

\courssub{The Pencil of Planes Through a Line}

\begin{experience}[The door-hinge model]
Take two non-parallel planes $\Pi_1$ and $\Pi_2$. They meet in a line $\Delta$. Now rotate $\Pi_1$ about $\Delta$: at every instant you obtain a new plane, and every one of them contains $\Delta$. Conversely, any plane containing $\Delta$ is one of these positions. The set of all planes through $\Delta$ is the pencil of planes with \emph{axis} $\Delta$.
\end{experience}

\begin{propriete}[Pencil of planes through the line of intersection of two planes]
Let $\Pi_1 : a_1x+b_1y+c_1z+d_1=0$ and $\Pi_2 : a_2x+b_2y+c_2z+d_2=0$ be two non-parallel planes, intersecting in a line $\Delta$. Then every plane through $\Delta$, \textbf{except $\Pi_2$ itself}, has an equation of the form
\[
\Pi_1+\lambda\Pi_2=0,\qquad\text{i.e.}\qquad (a_1x+b_1y+c_1z+d_1)+\lambda(a_2x+b_2y+c_2z+d_2)=0,
\]
for some real number $\lambda$. This proof \textbf{is examinable} in outline.
\end{propriete}

\textbf{Justification.} Write $P_1(x,y,z)=a_1x+b_1y+c_1z+d_1$ and $P_2(x,y,z)=a_2x+b_2y+c_2z+d_2$.
\begin{itemize}
\item \emph{The combination is a plane containing $\Delta$.} The equation $P_1+\lambda P_2=0$ is linear in $x,y,z$, so it represents a plane (its normal is $\vec n_1+\lambda\vec n_2\neq\vec 0$ since $\vec n_1,\vec n_2$ are not parallel). If $M\in\Delta$, then $P_1(M)=0$ and $P_2(M)=0$, so $P_1(M)+\lambda P_2(M)=0$: the plane contains the whole of $\Delta$.
\item \emph{Every plane through $\Delta$ (other than $\Pi_2$) is obtained.} Let $\Pi$ be a plane through $\Delta$, $\Pi\neq\Pi_2$. Pick a point $A\in\Pi$ with $A\notin\Delta$; then $P_2(A)\neq 0$ (otherwise $A\in\Pi_2$, and $\Pi=\Pi_2$ since both contain $\Delta$ and $A$). Setting $\lambda=-\dfrac{P_1(A)}{P_2(A)}$ makes $P_1(A)+\lambda P_2(A)=0$, so the plane $P_1+\lambda P_2=0$ contains $\Delta$ and $A$: it is exactly $\Pi$.
\end{itemize}

\begin{attention}[The plane $\Pi_2$ is missing]
The family $\Pi_1+\lambda\Pi_2=0$ contains \textbf{every} plane through $\Delta$ except $\Pi_2$ itself (no value of $\lambda$ gives $\Pi_2$). If a problem asks for \emph{all} planes through $\Delta$ satisfying a condition, always check separately whether $\Pi_2$ also works. Alternatively, use the symmetric form $\lambda\Pi_1+\mu\Pi_2=0$ with $(\lambda,\mu)\neq(0,0)$, which covers both.
\end{attention}

\begin{popfigure}
\begin{tikzpicture}[scale=1.05]
% common line (axis)
\draw[line width=1.6pt,popdark] (-0.4,-2.6) -- (0.4,2.6) node[above right,font=\small]{$\Delta$};
% plane 1
\fill[popblueL,opacity=0.55] (-3.2,-1.2) -- (0.4,2.6) -- (2.2,1.4) -- (-1.4,-2.4) -- cycle;
\draw[popblue] (-3.2,-1.2) -- (0.4,2.6) -- (2.2,1.4) -- (-1.4,-2.4) -- cycle;
\node[popblue,font=\small] at (-2.6,-1.9) {$\Pi_1$};
% plane 2
\fill[poporange,opacity=0.35] (3.2,-1.2) -- (0.4,2.6) -- (-2.2,1.4) -- (1.4,-2.4) -- cycle;
\draw[poporange] (3.2,-1.2) -- (0.4,2.6) -- (-2.2,1.4) -- (1.4,-2.4) -- cycle;
\node[poporange,font=\small] at (2.7,-1.9) {$\Pi_2$};
% plane 3 (member of pencil)
\fill[popgreenL,opacity=0.5] (-2.6,1.2) -- (0.4,2.6) -- (2.6,-1.8) -- (-0.6,-2.6) -- cycle;
\draw[popgreen] (-2.6,1.2) -- (0.4,2.6) -- (2.6,-1.8) -- (-0.6,-2.6) -- cycle;
\node[popgreen,font=\small] at (2.3,-2.3) {$\Pi_1+\lambda\Pi_2=0$};
\draw[line width=1.6pt,popdark] (-0.4,-2.6) -- (0.4,2.6);
\end{tikzpicture}
\legende{A pencil of planes: all members fan about their common line $\Delta=\Pi_1\cap\Pi_2$.}
\end{popfigure}

\courssub{Determining $\lambda$ From an Extra Condition}

\begin{methode}[Finding the member of a pencil satisfying one extra condition]
\begin{enumerate}
\item Write the pencil $(a_1x+b_1y+c_1z+d_1)+\lambda(a_2x+b_2y+c_2z+d_2)=0$.
\item Translate the extra condition into an equation in $\lambda$:
\begin{itemize}
\item \emph{Passes through $A(x_0,y_0,z_0)$:} substitute the coordinates, giving $P_1(A)+\lambda P_2(A)=0$.
\item \emph{Perpendicular to a plane with normal $\vec n$:} use the normals, $\vec n\cdot(\vec n_1+\lambda\vec n_2)=0$.
\item \emph{At distance $d$ from a point $A$:} apply the distance formula to the combined plane and solve for $\lambda$ (often two values).
\end{itemize}
\item Solve for $\lambda$ and substitute back to get the Cartesian equation. Check whether $\Pi_2$ itself should also be retained.
\end{enumerate}
\end{methode}

\begin{attention}[Perpendicularity uses normals only]
When the condition is ``perpendicular to a given plane'', work with the \textbf{normal vectors}: $\vec n\cdot(\vec n_1+\lambda\vec n_2)=0$. A very common GCE error is to substitute the point or to mix full plane equations. Two planes are perpendicular if and only if their normals are perpendicular.
\end{attention}

\begin{exoresolu}[Plane through a line of intersection and a point]
Find the equation of the plane containing the line of intersection of $\Pi_1: x+y+z-3=0$ and $\Pi_2: 2x-y+z-1=0$ and passing through $A(1,2,1)$.
\tcblower
The pencil is $(x+y+z-3)+\lambda(2x-y+z-1)=0$. Substituting $A(1,2,1)$:
\[
(1+2+1-3)+\lambda(2-2+1-1)=0 \;\Longrightarrow\; 1+0\cdot\lambda=0,
\]
which is impossible. So no member of the pencil works: the required plane must be $\Pi_2$ itself. Check: $2(1)-2+1-1=0$ \checkmark, so $A\in\Pi_2$. The answer is $2x-y+z-1=0$. This example shows exactly why the warning about $\Pi_2$ matters.
\end{exoresolu}

\begin{exoresolu}[Plane through a line of intersection, perpendicular to a given plane]
Find the plane through the line of intersection of $\Pi_1: x+2y-z+1=0$ and $\Pi_2: x-y+2z-3=0$ which is perpendicular to the plane $\Sigma: x+y+z=0$.
\tcblower
The pencil is $(x+2y-z+1)+\lambda(x-y+2z-3)=0$, with normal
\[
\vec n=\vec n_1+\lambda\vec n_2=(1+\lambda,\;2-\lambda,\;-1+2\lambda).
\]
Perpendicularity to $\Sigma$ (normal $\vec m=(1,1,1)$) requires $\vec n\cdot\vec m=0$:
\[
(1+\lambda)+(2-\lambda)+(-1+2\lambda)=0 \;\Longrightarrow\; 2+2\lambda=0 \;\Longrightarrow\; \lambda=-1.
\]
Substituting back: $(x+2y-z+1)-(x-y+2z-3)=0$, i.e.
\[
3y-3z+4=0.
\]
Check $\Pi_2$ separately: $\vec n_2\cdot\vec m=1-1+2=2\neq0$, so $\Pi_2$ is not a solution. The unique answer is $3y-3z+4=0$.
\end{exoresolu}

\courssub{Planes Through the Point of Intersection of Three Planes}

Three planes $\Pi_1,\Pi_2,\Pi_3$ generally meet at a single point $I$, found by solving the $3\times3$ system of their Cartesian equations. To build a new plane through $I$ satisfying further conditions (e.g.\ containing a given line, or perpendicular to a given plane), first compute $I$, then construct the plane using $I$ and the extra data.

\begin{methode}[Plane through the intersection point of three planes]
\begin{enumerate}
\item Solve the system $\Pi_1=0,\ \Pi_2=0,\ \Pi_3=0$ simultaneously (elimination or substitution) to find $I$.
\item Use $I$ together with the extra condition: e.g.\ if the plane must contain a line with direction $\vec u$ and a point $B$, take normal $\vec n=\overrightarrow{IB}\wedge\vec u$.
\item Write the Cartesian equation $\vec n\cdot\overrightarrow{IM}=0$ and expand.
\end{enumerate}
\end{methode}

\begin{exemplebox}[Building a plane through a triple intersection point]
Let $\Pi_1: x+y+z=6$, $\Pi_2: x-y+z=2$, $\Pi_3: 2x+y-z=3$. Subtracting $\Pi_2$ from $\Pi_1$: $2y=4$, so $y=2$. Adding $\Pi_1$ and $\Pi_3$: $3x+2y=9$, so $3x=5$, $x=\tfrac53$. Then $z=6-\tfrac53-2=\tfrac73$. Hence $I\left(\tfrac53,2,\tfrac73\right)$. Any plane required to pass through the common point of these three planes must pass through $I$; the remaining conditions then fix it completely.
\end{exemplebox}

\courssub{Bisector Planes of Two Intersecting Planes (GCE Extension)}

The \cle{bisector planes} of $\Pi_1$ and $\Pi_2$ are the two planes through $\Delta=\Pi_1\cap\Pi_2$ whose points are \emph{equidistant} from $\Pi_1$ and $\Pi_2$ — they bisect the dihedral angles. Using the distance formula, a point $M(x,y,z)$ lies on a bisector iff
\[
\frac{|a_1x+b_1y+c_1z+d_1|}{\sqrt{a_1^2+b_1^2+c_1^2}}=\frac{|a_2x+b_2y+c_2z+d_2|}{\sqrt{a_2^2+b_2^2+c_2^2}},
\]
giving the \textbf{two} bisector planes by taking the $+$ and $-$ signs. The two bisectors are perpendicular to each other. This is occasionally examined as a starred GCE question: flag it as an extension of the distance formula.

\courssub{Synthesis: GCE Paper-Style Problem}

\begin{exoresolu}[Full synthesis problem]
The planes $\Pi_1: x-2y+2z-3=0$ and $\Pi_2: 2x+y-2z+1=0$ are given, together with the point $A(1,0,2)$.
\begin{enumerate}
\item Show that $\Pi_1$ and $\Pi_2$ are perpendicular.
\item Find the perpendicular distance from $A$ to $\Pi_1$.
\item Find the equation of the plane through the line of intersection of $\Pi_1$ and $\Pi_2$ which passes through $A$.
\item Find the angle between $\Pi_1$ and the plane $\Sigma: z=0$.
\end{enumerate}
\tcblower
\textbf{1.} Normals: $\vec n_1=(1,-2,2)$, $\vec n_2=(2,1,-2)$. Then $\vec n_1\cdot\vec n_2=2-2-4=-4\neq0$? Recompute: $1(2)+(-2)(1)+2(-2)=2-2-4=-4$. Since this is not zero, the planes are \emph{not} perpendicular — the statement to verify fails, and in an examination you would say so honestly. (The angle between them satisfies $\cos\theta=\dfrac{|-4|}{3\cdot3}=\dfrac49$.)

\textbf{2.} Distance from $A$ to $\Pi_1$:
\[
d=\frac{|1-2(0)+2(2)-3|}{\sqrt{1+4+4}}=\frac{|1+4-3|}{3}=\frac{2}{3}.
\]

\textbf{3.} Pencil: $(x-2y+2z-3)+\lambda(2x+y-2z+1)=0$. Substituting $A(1,0,2)$:
\[
(1-0+4-3)+\lambda(2+0-4+1)=0 \;\Longrightarrow\; 2-\lambda=0 \;\Longrightarrow\; \lambda=2.
\]
Hence $(x-2y+2z-3)+2(2x+y-2z+1)=0$, i.e.\ $5x-2z-1=0$.

\textbf{4.} $\Sigma$ has normal $\vec k=(0,0,1)$. The angle $\theta$ between the planes equals the angle between the normals:
\[
\cos\theta=\frac{|\vec n_1\cdot\vec k|}{\|\vec n_1\|\,\|\vec k\|}=\frac{|2|}{3\cdot1}=\frac23,\qquad \theta\approx 48.2^\circ.
\]
\end{exoresolu}

\courssub{Revision Grid: Choosing the Right Tool}

\begin{center}
\begin{tabularx}{\linewidth}{|l|X|X|}
\hline
\rowcolor{popblueL} \textbf{Given} & \textbf{Tool} & \textbf{Result} \\
\hline
Point $A$ + normal $\vec n$ & $\vec n\cdot\overrightarrow{AM}=0$ & Cartesian equation of plane \\
\hline
Three non-collinear points & $\vec n=\overrightarrow{AB}\wedge\overrightarrow{AC}$ & Equation of the plane \\
\hline
Distance point--plane & $d=\dfrac{|ax_0+by_0+cz_0+d|}{\sqrt{a^2+b^2+c^2}}$ & Numerical distance \\
\hline
Angle between two planes & $\cos\theta=\dfrac{|\vec n_1\cdot\vec n_2|}{\|\vec n_1\|\|\vec n_2\|}$ & Acute angle \\
\hline
Angle line--plane & $\sin\theta=\dfrac{|\vec u\cdot\vec n|}{\|\vec u\|\|\vec n\|}$ & Acute angle \\
\hline
Line meets plane & Substitute parametric line into plane & Intersection point \\
\hline
Two planes meet & Solve system / $\vec u=\vec n_1\wedge\vec n_2$ & Line of intersection \\
\hline
Plane through a line of intersection & Pencil $\Pi_1+\lambda\Pi_2=0$ & Family; fix $\lambda$ by extra condition \\
\hline
\end{tabularx}
\end{center}

\begin{popfigure}
\begin{tikzpicture}[font=\small, node distance=0.45cm,
box/.style={draw=popblue, fill=popblueL, rounded corners=2pt, text width=3.1cm, align=center, minimum height=0.85cm},
tool/.style={draw=poporange, fill=white, rounded corners=2pt, text width=3.4cm, align=center, minimum height=0.85cm},
res/.style={draw=popgreen, fill=popgreenL, rounded corners=2pt, text width=2.9cm, align=center, minimum height=0.85cm}]
\node[box] (g1) {Point + normal, or 3 points};
\node[tool, right=0.7cm of g1] (t1) {Normal $\vec n$; cross product $\overrightarrow{AB}\wedge\overrightarrow{AC}$};
\node[res, right=0.7cm of t1] (r1) {Plane equation};
\node[box, below=of g1] (g2) {Point + plane};
\node[tool, right=0.7cm of g2] (t2) {Distance formula};
\node[res, right=0.7cm of t2] (r2) {Perpendicular distance};
\node[box, below=of g2] (g3) {Two planes / line + plane};
\node[tool, right=0.7cm of g3] (t3) {Dot product of normals (or $\vec u,\vec n$)};
\node[res, right=0.7cm of t3] (r3) {Angle; intersection};
\node[box, below=of g3] (g4) {Line of intersection + 1 condition};
\node[tool, right=0.7cm of g4] (t4) {Pencil $\Pi_1+\lambda\Pi_2=0$};
\node[res, right=0.7cm of t4] (r4) {Required plane};
\foreach \a/\b in {g1/t1,t1/r1,g2/t2,t2/r2,g3/t3,t3/r3,g4/t4,t4/r4}
  \draw[-latex, popmuted, thick] (\a) -- (\b);
\end{tikzpicture}
\legende{Summary flowchart: Given $\to$ Tool $\to$ Result for the main problem types of the chapter.}
\end{popfigure}

\begin{aretenir}[Key points of the section]
\begin{itemize}
\item The pencil of planes through $\Delta=\Pi_1\cap\Pi_2$ is $\Pi_1+\lambda\Pi_2=0$; it contains every plane through $\Delta$ \textbf{except $\Pi_2$}.
\item Fix $\lambda$ by the extra condition: substitute a point, use $\vec n\cdot(\vec n_1+\lambda\vec n_2)=0$ for perpendicularity, or the distance formula for a distance condition.
\item Three planes generally meet at one point: solve the $3\times3$ system, then build the required plane.
\item Bisector planes come from equal perpendicular distances (two solutions, $+$ and $-$).
\item In synthesis problems, read each part separately and pick the tool from the revision grid.
\end{itemize}
\end{aretenir}

\begin{exercice}[Practice]
\begin{enumerate}
\item Find the plane through the line of intersection of $x+y-z=2$ and $2x-y+z=1$ passing through the origin.
\item Find the planes of the pencil of Question 1 which are at distance $\sqrt{6}$ from the point $(2,2,2)$.
\item Find the bisector planes of $3x+4y+1=0$ and $4x-3y+2=0$ (in 3D, with $z$ absent).
\item Show that the planes $x+y+z=6$, $x-y+z=2$ and $2x+y-z=3$ meet at a single point, and find the plane through this point perpendicular to the line $\dfrac{x}{1}=\dfrac{y}{-1}=\dfrac{z}{2}$.
\end{enumerate}
\end{exercice}

\begin{corrige}[Exercises 1--4]
\textbf{1.} Pencil: $(x+y-z-2)+\lambda(2x-y+z-1)=0$. Origin: $-2+\lambda(-1)=0$, so $\lambda=-2$. Substituting: $(x+y-z-2)-2(2x-y+z-1)=0$, i.e.\ $-3x+3y-3z=0$, or $x-y+z=0$.

\textbf{2.} Combined plane: $(1+2\lambda)x+(1-\lambda)y+(-1+\lambda)z-(2+\lambda)=0$. Distance from $(2,2,2)$:
\[
\frac{|2(1+2\lambda)+2(1-\lambda)+2(-1+\lambda)-(2+\lambda)|}{\sqrt{(1+2\lambda)^2+(1-\lambda)^2+(-1+\lambda)^2}}=\frac{|3\lambda|}{\sqrt{6\lambda^2+3}}=\sqrt6.
\]
Squaring: $9\lambda^2=6(6\lambda^2+3)$, i.e.\ $9\lambda^2=36\lambda^2+18$, which gives $-27\lambda^2=18$: no real solution. Hence \emph{no} plane of the pencil is at distance $\sqrt6$ from that point — a useful reminder that a condition may have no solution.

\textbf{3.} Equal distances: $\dfrac{|3x+4y+1|}{5}=\dfrac{|4x-3y+2|}{5}$, so $3x+4y+1=\pm(4x-3y+2)$. With $+$: $x-7y+1=0$. With $-$: $7x+y+3=0$. These are the two bisector planes (both parallel to the $z$-axis).

\textbf{4.} The three planes are exactly those of the \texttt{exemplebox} above. We already found their unique intersection point $I\left(\tfrac53,2,\tfrac73\right)$. The line $\frac{x}{1}=\frac{y}{-1}=\frac{z}{2}$ has direction vector $\vec u=(1,-1,2)$. The required plane has normal $\vec n=\vec u=(1,-1,2)$ and passes through $I$. Its equation is
\[
1\left(x-\tfrac53\right)-1(y-2)+2\left(z-\tfrac73\right)=0
\;\Longrightarrow\; x-\tfrac53-y+2+2z-\tfrac{14}{3}=0
\;\Longrightarrow\; x-y+2z-\tfrac{13}{3}=0,
\]
or, clearing denominators, $3x-3y+6z-13=0$.
\end{corrige}

\begin{savaistu}[Pencils beyond the GCE]
Pencils are a first taste of \emph{linear families} of geometric objects, an idea that runs through university mathematics: pencils of lines, of circles, of conics. The single parameter $\lambda$ organises an infinity of objects into one simple algebraic family — a beautiful example of algebra taming geometry.
\end{savaistu}

\newpage

\coursec{Integration activity}

\coursec{Activity of Integration: Designing a School Roof in Bamenda}

\courssub{Aims of the Activity}
This integration activity closes the chapter on vectors and planes by mobilising, in a single authentic task, all the competencies developed in Lessons 42 to 51. By the end of the activity, each group of learners should be able to:
\begin{itemize}
\item determine the \cle{Cartesian equation of a plane} through three given non-collinear points, using a normal vector obtained from a cross product;
\item verify that a line lies in a plane and obtain the \cle{vector equation of the line of intersection} of two planes;
\item compute the \cle{angle between two planes} from their normal vectors, and interpret the result (acute angle versus interior dihedral angle) in a real design specification;
\item apply the formula for the \cle{perpendicular distance from a point to a plane} in a concrete measurement problem;
\item find the \cle{point of intersection of a line and a plane} by substituting a parametric line into a Cartesian equation;
\item communicate a complete mathematical argument on a poster, with a clear three-dimensional sketch, as required in the problem-solving and presentation component of the GCE Advanced Level examination.
\end{itemize}

\begin{aretenir}[Skills Mobilised]
Cross product and normal vectors; Cartesian and vector equations of planes; line of intersection of two planes; angle between two planes; distance from a point to a plane; intersection of a line with a plane; critical interpretation of a numerical result against a specification.
\end{aretenir}

\courssub{Situation}
\textbf{Designing a School Roof in Bamenda.}\par
A contractor in Bamenda has been awarded the construction of the roof of a new classroom block for a secondary school. The roof consists of two rectangular metal panels joined along a horizontal \cle{ridge line}. The architect has supplied a three-dimensional coordinate model (units in metres), with the $x$-axis along the ridge, the $y$-axis across the building and the $z$-axis vertical.

\textbf{Panel 1} (the lower, gentler slope) passes through the points
\[ A(0,0,4),\qquad B(6,0,4),\qquad C(6,4,2). \]

\textbf{Panel 2} (the steeper slope) passes through the points
\[ B(6,0,4),\qquad D(12,0,4),\qquad E(12,4,6). \]

The two panels must meet along the segment $BD$, which is the ridge line of the roof.

The architect's specification booklet imposes two checks before the panels can be ordered from the workshop:
\begin{enumerate}
\item the \emph{interior} angle between the two panels (the angle measured inside the building, below the ridge) must be $120^\circ$;
\item a concrete pillar whose top is at the point $P(3,2,0)$ must just reach Panel 1, so its required height must be computed exactly.
\end{enumerate}

Finally, an electric cable is to be run in a straight line along
\[ \vec{r} = (\vec{i} + 6\vec{k}) + \lambda(\vec{i} + \vec{j} - \vec{k}), \]
and the electrician needs the exact point where this cable pierces Panel 1 in order to fix a junction box there.

Each group of learners plays the role of the contractor's technical team: carry out all the verifications, decide whether the roof design satisfies the architect's specification, and present the full solution on a poster including a three-dimensional sketch of the roof.

\begin{attention}[Units and Interpretation]
All coordinates are in metres. Angles computed from normal vectors using the modulus formula are always acute (between $0^\circ$ and $90^\circ$): deciding whether the \emph{interior} angle of the roof equals this acute angle or its supplement is part of the modelling work.
\end{attention}

\courssub{Tasks}

\textbf{Task 1: Equation of each panel.}
\begin{enumerate}
\item For Panel 1, compute the vectors $\overrightarrow{AB}$ and $\overrightarrow{AC}$, then the cross product $\overrightarrow{AB} \times \overrightarrow{AC}$. Deduce a normal vector $\vec{n}_1$ and the Cartesian equation of Panel 1.
\item In the same way, using $\overrightarrow{BD}$ and $\overrightarrow{BE}$, find a normal vector $\vec{n}_2$ and the Cartesian equation of Panel 2.
\end{enumerate}

\textbf{Task 2: The ridge line.}
\begin{enumerate}
\item[3.] Show that every point of the segment $BD$ satisfies the equations of \emph{both} panels. What does this prove geometrically?
\item[4.] Give a vector equation of the line of intersection of the two panels (the ridge line).
\end{enumerate}

\textbf{Task 3: The angle between the panels.}
\begin{enumerate}
\item[5.] Compute the acute angle $\theta$ between the two panels using $\cos\theta = \dfrac{|\vec{n}_1 \cdot \vec{n}_2|}{|\vec{n}_1|\,|\vec{n}_2|}$.
\item[6.] Deduce the interior angle of the roof below the ridge. Does the design satisfy the architect's specification of $120^\circ$? Write one sentence of advice to the contractor.
\end{enumerate}

\textbf{Task 4: The pillar.}
\begin{enumerate}
\item[7.] Compute the perpendicular distance from $P(3,2,0)$ to Panel 1. Hence state the exact height the pillar must have.
\end{enumerate}

\textbf{Task 5: The cable.}
\begin{enumerate}
\item[8.] Find the value of $\lambda$ for which the cable line meets Panel 1, and give the coordinates of the piercing point. Verify that this point lies within the rectangular panel (check the ranges of $x$ and $y$).
\end{enumerate}

\textbf{Task 6: Poster.} Present all results on a poster with a labelled 3D sketch showing the two panels, the ridge line, the pillar and the cable.

\courssub{Model Answer}

\textbf{Task 1. Equations of the Panels.}
\paragraph{Panel 1.}
\[
\overrightarrow{AB} = \begin{pmatrix} 6 \\ 0 \\ 0 \end{pmatrix},\qquad
\overrightarrow{AC} = \begin{pmatrix} 6 \\ 4 \\ -2 \end{pmatrix}.
\]
\[
\vec{n}_1 = \overrightarrow{AB} \times \overrightarrow{AC}
= \begin{pmatrix}
0\cdot(-2) - 0\cdot4 \\
0\cdot6 - 6\cdot(-2) \\
6\cdot4 - 0\cdot6
\end{pmatrix}
= \begin{pmatrix} 0 \\ 12 \\ 24 \end{pmatrix}
\parallel \begin{pmatrix} 0 \\ 1 \\ 2 \end{pmatrix}.
\]
A Cartesian equation is $0\cdot x + 1\cdot y + 2\cdot z = d$. Substituting $A(0,0,4)$ gives $d = 8$. Hence
\[
\[\boxed{\text{Panel 1: } y + 2z = 8.}\]
\]

\paragraph{Panel 2.}
\[
\overrightarrow{BD} = \begin{pmatrix} 6 \\ 0 \\ 0 \end{pmatrix},\qquad
\overrightarrow{BE} = \begin{pmatrix} 6 \\ 4 \\ 2 \end{pmatrix},\qquad
\vec{n}_2 = \overrightarrow{BD} \times \overrightarrow{BE}
= \begin{pmatrix}
0\cdot2 - 0\cdot4 \\
0\cdot6 - 6\cdot2 \\
6\cdot4 - 0\cdot6
\end{pmatrix}
= \begin{pmatrix} 0 \\ -12 \\ 24 \end{pmatrix}
\parallel \begin{pmatrix} 0 \\ -1 \\ 2 \end{pmatrix}.
\]
A Cartesian equation is $0\cdot x - 1\cdot y + 2\cdot z = d$. Substituting $B(6,0,4)$ gives $d = 8$. Hence
\[
\[\boxed{\text{Panel 2: } -y + 2z = 8.}\]
\]

\textbf{Task 2. The Ridge Line.}
\paragraph{3.} A general point of the segment $BD$ has coordinates $(x,0,4)$ with $6 \le x \le 12$.
\begin{itemize}
\item Panel 1: $0 + 2(4) = 8$ \checkmark.
\item Panel 2: $-0 + 2(4) = 8$ \checkmark.
\end{itemize}
Every point of $BD$ satisfies both equations; therefore the whole segment $BD$ lies in both planes. Geometrically, this proves that the two panels intersect exactly along the ridge line $BD$, as required by the design.

\paragraph{4.} Using point $B(6,0,4)$ and direction vector $\vec{i} = (1,0,0)$, a vector equation of the ridge line is
\[
\vec{r} = \begin{pmatrix} 6 \\ 0 \\ 4 \end{pmatrix} + t \begin{pmatrix} 1 \\ 0 \\ 0 \end{pmatrix},\qquad t \in \mathbb{R}.
\]

\textbf{Task 3. Angle Between the Panels.}
\paragraph{5.} With simplified normal vectors $\vec{n}_1 = (0,1,2)$ and $\vec{n}_2 = (0,-1,2)$:
\[
\vec{n}_1 \cdot \vec{n}_2 = 0\cdot0 + 1\cdot(-1) + 2\cdot2 = 3,\qquad
|\vec{n}_1| = |\vec{n}_2| = \sqrt{0^2+1^2+2^2} = \sqrt{5}.
\]
\[
\cos\theta = \frac{|\vec{n}_1 \cdot \vec{n}_2|}{|\vec{n}_1|\,|\vec{n}_2|}
= \frac{|3|}{\sqrt{5}\cdot\sqrt{5}} = \frac{3}{5}
\quad\Longrightarrow\quad
\theta = \arccos\!\left(\frac{3}{5}\right) \approx 53.13^\circ.
\]

\paragraph{6.} The acute angle between the planes is $\theta \approx 53.13^\circ$. The normals $\vec{n}_1$ and $\vec{n}_2$ both have positive $z$-components (pointing generally upward), so they point to the same side of the ridge line (the exterior of the roof). Consequently, the interior angle of the roof (measured inside the building, below the ridge) is the supplement of $\theta$:
\[
\text{Interior angle} = 180^\circ - \theta \approx 180^\circ - 53.13^\circ = 126.87^\circ.
\]
\textbf{Conclusion:} The design does \emph{not} satisfy the architect's specification of $120^\circ$; the interior angle is about $6.9^\circ$ too large (the roof is too open). The contractor should ask the architect to adjust the slopes (for example by lowering point $C$ or raising point $E$) before ordering the panels.

\textbf{Task 4. The Pillar Height.}
\paragraph{7.} Panel 1 has equation $y + 2z - 8 = 0$. The perpendicular distance from $P(3,2,0)$ to this plane is
\[
d = \frac{|\,2 + 2(0) - 8\,|}{\sqrt{0^2 + 1^2 + 2^2}}
= \frac{6}{\sqrt{5}}
= \frac{6\sqrt{5}}{5}\ \mathrm{m}
\approx 2.68\ \mathrm{m}.
\]
The pillar must be exactly $\dfrac{6\sqrt{5}}{5}\ \mathrm{m}$ tall.

\textbf{Task 5. The Cable Piercing Point.}
\paragraph{8.} The cable line is given in parametric form:
\[
x = 1+\lambda,\quad y = \lambda,\quad z = 6-\lambda.
\]
Substitute into the equation of Panel 1 ($y + 2z = 8$):
\[
\lambda + 2(6-\lambda) = 8
\;\Longrightarrow\;
12 - \lambda = 8
\;\Longrightarrow\;
\lambda = 4.
\]
The piercing point is
\[
(1+4,\; 4,\; 6-4) = (5,\,4,\,2).
\]
Panel 1 is the rectangle with vertices $A(0,0,4)$, $B(6,0,4)$, $C(6,4,2)$ and the fourth vertex $(0,4,2)$. Its projection on the $xy$-plane covers $0 \le x \le 6$, $0 \le y \le 4$. The point $(5,4,2)$ satisfies $0 \le 5 \le 6$ and $y=4$ (on the boundary). It lies on the lower edge of Panel 1 (the edge through $C$ parallel to the ridge). The junction box will be fixed at $(5,4,2)$, on that lower edge.

\begin{popfigure}
\begin{tikzpicture}[scale=0.7, line join=round, line cap=round]
% Define 3D to 2D projection: x right, y up-left, z up
% Basis vectors in 2D:
% x -> (1.0, 0.0)
% y -> (-0.3, -0.3)
% z -> (0.0, 0.6)
\def\px{1.0} \def\py{0.0}
\def\qx{-0.3} \def\qy{-0.3}
\def\rx{0.0} \def\ry{0.6}
% Function to project (x,y,z) -> (X,Y)
% X = x*px + y*qx + z*rx
% Y = x*py + y*qy + z*ry
% Points:
% A(0,0,4) -> (0, 2.4)
% B(6,0,4) -> (6, 2.4)
% C(6,4,2) -> (6 -1.2, -1.2 +1.2) = (4.8, 0)
% D(12,0,4) -> (12, 2.4)
% E(12,4,6) -> (12 -1.2, -1.2 +3.6) = (10.8, 2.4)
% Fourth vertex of Panel 1: A + (C-B) = (0,4,2) -> (-1.2, 0)
% Fourth vertex of Panel 2: B + (E-D) = (6,4,6) -> (6 -1.2, -1.2 +3.6) = (4.8, 2.4)
% P(3,2,0) -> (3 -0.6, -0.6) = (2.4, -0.6)
% Cable point (5,4,2) -> (5 -1.2, -1.2 +1.2) = (3.8, 0)

\coordinate (A) at (0, 2.4);
\coordinate (B) at (6, 2.4);
\coordinate (C) at (4.8, 0);
\coordinate (A4) at (-1.2, 0); % (0,4,2)
\coordinate (D) at (12, 2.4);
\coordinate (E) at (10.8, 2.4);
\coordinate (B4) at (4.8, 2.4); % (6,4,6)
\coordinate (P) at (2.4, -0.6);
\coordinate (Cable) at (3.8, 0);

% Axes (origin at (0,0,0) -> (0,0))
\coordinate (O) at (0,0);
\draw[popmuted, ->] (O) -- (7,0) node[right] {$x$ (ridge)};
\draw[popmuted, ->] (O) -- (-2.5,-2.5) node[below left] {$y$ (across)};
\draw[popmuted, ->] (O) -- (0,4) node[above] {$z$ (vertical)};

% Panel 1 (lower slope) - popblue
\fill[popblueL, opacity=0.6] (A) -- (B) -- (C) -- (A4) -- cycle;
\draw[popblue, thick] (A) -- (B) -- (C) -- (A4) -- cycle;
\draw[popblue, dashed] (A) -- (C); % diagonal for depth

% Panel 2 (steeper slope) - popgreen
\fill[popgreenL, opacity=0.6] (B) -- (D) -- (E) -- (B4) -- cycle;
\draw[popgreen!70!black, thick] (B) -- (D) -- (E) -- (B4) -- cycle;
\draw[popgreen!70!black, dashed] (B) -- (E);

% Ridge line
\draw[poporange, very thick] (B) -- (D);
\node[poporange, above] at (9, 2.5) {Ridge $BD$};

% Labels
\node[popblue, left] at (A) {$A(0,0,4)$};
\node[popblue, above] at (B) {$B(6,0,4)$};
\node[popblue, right] at (C) {$C(6,4,2)$};
\node[popblue, left] at (A4) {$(0,4,2)$};
\node[popgreen!70!black, above] at (D) {$D(12,0,4)$};
\node[popgreen!70!black, right] at (E) {$E(12,4,6)$};
\node[popgreen!70!black, above] at (B4) {$(6,4,6)$};

% Pillar
\draw[popdark, dashed, thick] (P) -- (3.8, 0); % P to panel 1 along vertical? No, perpendicular.
% Perpendicular from P to plane y+2z=8. Normal (0,1,2). In projection, direction is (qy, ry) scaled? 
% Actually, just draw a dashed line from P to the piercing point on panel 1 vertically below? 
% The problem says pillar top at P(3,2,0) reaches Panel 1. The distance is perpendicular.
% The foot of perpendicular from P to plane: P + d*n/|n|. n=(0,1,2). d=6/√5.
% Foot = (3,2,0) + (6/5)*(0,1,2) = (3, 3.2, 2.4). Projected: (3 -0.96, -0.96 +1.44) = (2.04, 0.48).
\coordinate (Foot) at (2.04, 0.48);
\draw[popdark, dashed, thick] (P) -- (Foot);
\node[popdark, below] at (P) {$P(3,2,0)$};
\node[popdark, left] at (Foot) {Foot};
\node[popdark, right] at (3.1, -0.2) {$\frac{6\sqrt5}{5}$ m};

% Cable
\draw[poppink, thick, ->] (1, 3.6) -- (Cable); % from (1,0,6) projected: (1, 3.6)
\fill[poppink] (Cable) circle (2pt);
\node[poppink, below right] at (Cable) {Cable $(5,4,2)$};
\node[poppink, above left] at (1, 3.6) {$(1,0,6)$};

% Panel labels
\node[popblue, opacity=0.8] at (2, 1.2) {Panel 1};
\node[popgreen!70!black, opacity=0.8] at (8.5, 1.8) {Panel 2};
\end{tikzpicture}
\legende{Sketch of the two roof panels meeting along the ridge line $BD$. The pillar from $P$ meets Panel 1 perpendicularly; the cable pierces Panel 1 at $(5,4,2)$ on its lower edge. Axes show the coordinate orientation.}
\end{popfigure}

\begin{attention}[What the Examiners Expect]
In such synthesis tasks, marks are awarded for:
\begin{itemize}
\item correct computation of cross products and normal vectors;
\item correct Cartesian equations of planes (simplified integer coefficients preferred);
\item correct use of the modulus in the angle formula $\cos\theta = \frac{|\vec{n}_1\cdot\vec{n}_2|}{|\vec{n}_1||\vec{n}_2|}$;
\item a clear justification that the interior angle is the supplement of the acute angle between the normals (by considering the orientation of the normals relative to the interior of the roof);
\item exact values left in surd form (e.g. $\frac{6\sqrt5}{5}$) rather than decimal approximations;
\item verification that the intersection point lies within the physical boundaries of the panel;
\item a concluding sentence that answers the practical question directly (here: the $120^\circ$ specification is \emph{not} met).
\end{itemize}
\end{attention}

\newpage

\coursec{Methods and worked examples}

\coursec{Equations of a Plane: Vector, Parametric and Cartesian Forms}

\begin{definition}[Plane in space]
A \cle{plane} $\Pi$ is a flat, two-dimensional surface extending infinitely in all directions. It is uniquely determined by:
\begin{itemize}
\item a point $A$ and a non-zero normal vector $\vec{n}$ (perpendicular to the plane), or
\item three non-collinear points $A,B,C$, or
\item a point $A$ and two non-parallel direction vectors $\vec{u},\vec{v}$ lying in the plane.
\end{itemize}
\end{definition}

\begin{propriete}[Vector equation of a plane (normal form)]
Let $\Pi$ be a plane with normal vector $\vec{n}\neq\vec{0}$ passing through point $A$ with position vector $\vec{a}$. A point $P$ with position vector $\vec{r}$ lies on $\Pi$ if and only if $\overrightarrow{AP}$ is perpendicular to $\vec{n}$, i.e.
\[
(\vec{r}-\vec{a})\cdot\vec{n}=0 \quad\Longleftrightarrow\quad \vec{r}\cdot\vec{n}=\vec{a}\cdot\vec{n}.
\]
If $\vec{n}=a\vec{i}+b\vec{j}+c\vec{k}$ and $\vec{a}=x_0\vec{i}+y_0\vec{j}+z_0\vec{k}$, this becomes
\[
\vec{r}\cdot(a\vec{i}+b\vec{j}+c\vec{k})=ax_0+by_0+cz_0.
\]
\end{propriete}

\begin{propriete}[Cartesian equation of a plane]
Writing $\vec{r}=x\vec{i}+y\vec{j}+z\vec{k}$ in the vector equation gives the \cle{Cartesian equation}
\[
ax+by+cz=d,\qquad\text{where } d=ax_0+by_0+cz_0.
\]
The coefficients $(a,b,c)$ are the components of a normal vector $\vec{n}$. Conversely, any equation of the form $ax+by+cz=d$ with $(a,b,c)\neq(0,0,0)$ represents a plane with normal $\vec{n}=a\vec{i}+b\vec{j}+c\vec{k}$.
\end{propriete}

\begin{propriete}[Parametric vector equation of a plane]
If $\Pi$ passes through $A(\vec{a})$ and contains two non-parallel direction vectors $\vec{u}$ and $\vec{v}$, every point $P$ on $\Pi$ can be written as
\[
\vec{r}=\vec{a}+\lambda\vec{u}+\mu\vec{v},\qquad \lambda,\mu\in\mathbb{R}.
\]
In Cartesian parametric form:
\[
\begin{cases}
x=x_0+\lambda u_1+\mu v_1,\\
y=y_0+\lambda u_2+\mu v_2,\\
z=z_0+\lambda u_3+\mu v_3.
\end{cases}
\]
\end{propriete}

\begin{methode}[From a point and a normal vector to the Cartesian equation]
To find the equation of a plane $\Pi$ passing through $A(x_0,y_0,z_0)$ with normal vector $\vec{n}=a\vec{i}+b\vec{j}+c\vec{k}$:
\begin{enumerate}
\item Write the vector equation $\vec{r}\cdot\vec{n}=\vec{a}\cdot\vec{n}$.
\item Compute the constant $d=\vec{a}\cdot\vec{n}=ax_0+by_0+cz_0$.
\item Replace $\vec{r}$ by $x\vec{i}+y\vec{j}+z\vec{k}$ to obtain $ax+by+cz=d$.
\item \textbf{Reflex:} the coefficients of $x,y,z$ are \emph{always} the components of a normal vector.
\item \textbf{Check:} substitute the coordinates of $A$; the equality must hold.
\end{enumerate}
\end{methode}

\begin{exoresolu}[Equation of a plane given a point and a normal]
Find the vector and Cartesian equations of the plane passing through $A(2,-1,3)$ and perpendicular to $\vec{n}=3\vec{i}+2\vec{j}-\vec{k}$.
\tcblower
\textbf{Step 1: vector equation.}
\[
\vec{r}\cdot(3\vec{i}+2\vec{j}-\vec{k})=\vec{a}\cdot\vec{n}.
\]
\textbf{Step 2: compute the constant.}
\[
\vec{a}\cdot\vec{n}=(2)(3)+(-1)(2)+(3)(-1)=6-2-3=1.
\]
Vector equation: $\vec{r}\cdot(3\vec{i}+2\vec{j}-\vec{k})=1$.

\textbf{Step 3: Cartesian form.}
\[
\boxed{3x+2y-z=1}.
\]
\textbf{Check:} at $A(2,-1,3)$: $3(2)+2(-1)-3=6-2-3=1$. \checkmark
\end{exoresolu}

\begin{exoresolu}[Converting between forms]
A plane has parametric equations $x=1+2\lambda-\mu$, $y=2-\lambda+3\mu$, $z=-1+\lambda+2\mu$. Find its Cartesian equation.
\tcblower
\textbf{Step 1: identify $\vec{a},\vec{u},\vec{v}$.}
\[
\vec{a}=\begin{pmatrix}1\\2\\-1\end{pmatrix},\quad
\vec{u}=\begin{pmatrix}2\\-1\\1\end{pmatrix},\quad
\vec{v}=\begin{pmatrix}-1\\3\\2\end{pmatrix}.
\]
\textbf{Step 2: normal vector $\vec{n}=\vec{u}\times\vec{v}$.}
\[
\vec{n}=\begin{vmatrix}
\vec{i} & \vec{j} & \vec{k}\\
2 & -1 & 1\\
-1 & 3 & 2
\end{vmatrix}
=\vec{i}((-1)(2)-(1)(3))-\vec{j}((2)(2)-(1)(-1))+\vec{k}((2)(3)-(-1)(-1))
=\begin{pmatrix}-5\\-5\\5\end{pmatrix}.
\]
Simplify to $\vec{n}=(1,1,-1)$.

\textbf{Step 3: constant $d=\vec{n}\cdot\vec{a}=1+2+1=4$.}
Cartesian equation: $\boxed{x+y-z=4}$.
\end{exoresolu}

\begin{popfigure}
\begin{tikzpicture}[scale=0.8]
% Plane visualization
\draw[popblue, thick, fill=popblue!10] (0,0,0) -- (3,0,1) -- (3,2,1) -- (0,2,0) -- cycle;
\draw[->, thick, poporange] (1,1,0.5) -- (1.5,1.5,1.5) node[above right] {$\vec{n}$};
\draw[->, thick, popgreen] (1,1,0.5) -- (2,0.5,0.8) node[below right] {$\vec{u}$};
\draw[->, thick, popgreen] (1,1,0.5) -- (1.2,2,0.7) node[above] {$\vec{v}$};
\draw[fill=popdark] (1,1,0.5) circle (2pt) node[below left] {$A$};
\draw[dashed, gray] (0,0,0) -- (0,0,2) node[above] {$z$};
\draw[dashed, gray] (0,0,0) -- (4,0,0) node[right] {$x$};
\draw[dashed, gray] (0,0,0) -- (0,3,0) node[left] {$y$};
\end{tikzpicture}
\legende{Plane with normal $\vec{n}$ and direction vectors $\vec{u},\vec{v}$}
\end{popfigure}

\coursec{Plane Through Three Points and Special Planes}

\begin{methode}[Plane through three non-collinear points]
Given $A(x_1,y_1,z_1)$, $B(x_2,y_2,z_2)$, $C(x_3,y_3,z_3)$:
\begin{enumerate}
\item Form two direction vectors $\overrightarrow{AB}$ and $\overrightarrow{AC}$.
\item \textbf{Check collinearity:} if $\overrightarrow{AB}\times\overrightarrow{AC}=\vec{0}$, the points are collinear — infinitely many planes contain them; no unique plane exists.
\item \textbf{Cartesian route:} compute normal $\vec{n}=\overrightarrow{AB}\times\overrightarrow{AC}$, then $d=\vec{n}\cdot\vec{a}$ using any of the three points.
\item \textbf{Parametric route:} $\vec{r}=\vec{a}+\lambda\overrightarrow{AB}+\mu\overrightarrow{AC}$.
\item \textbf{Check:} verify all three points satisfy the final Cartesian equation.
\end{enumerate}
\end{methode}

\begin{exoresolu}[Plane through three points]
Find (a) the Cartesian equation and (b) the parametric vector equation of the plane through $A(1,0,2)$, $B(2,1,1)$ and $C(0,3,1)$.
\tcblower
\textbf{Step 1: direction vectors.}
\[
\overrightarrow{AB}=\begin{pmatrix}1\\1\\-1\end{pmatrix},\qquad
\overrightarrow{AC}=\begin{pmatrix}-1\\3\\-1\end{pmatrix}.
\]
\textbf{Step 2: normal vector.}
\[
\vec{n}=\overrightarrow{AB}\times\overrightarrow{AC}
=\begin{vmatrix}
\vec{i} & \vec{j} & \vec{k}\\
1 & 1 & -1\\
-1 & 3 & -1
\end{vmatrix}
=\vec{i}((1)(-1)-(-1)(3))-\vec{j}((1)(-1)-(-1)(-1))+\vec{k}((1)(3)-(1)(-1))
=\begin{pmatrix}2\\2\\4\end{pmatrix}.
\]
Simplify to $\vec{n}=(1,1,2)$.

\textbf{Step 3: constant term.} Using $A(1,0,2)$: $d=1(1)+1(0)+2(2)=5$.

\textbf{Cartesian equation:} $\boxed{x+y+2z=5}$.

\textbf{Check:} $B$: $2+1+2=5$ \checkmark; $C$: $0+3+2=5$ \checkmark.

\textbf{Parametric form:}
\[
\vec{r}=\begin{pmatrix}1\\0\\2\end{pmatrix}+\lambda\begin{pmatrix}1\\1\\-1\end{pmatrix}+\mu\begin{pmatrix}-1\\3\\-1\end{pmatrix},\quad \lambda,\mu\in\mathbb{R},
\]
i.e. $x=1+\lambda-\mu$, $y=\lambda+3\mu$, $z=2-\lambda-\mu$.
\end{exoresolu}

\begin{attention}[Collinear points]
If $\overrightarrow{AB}\times\overrightarrow{AC}=\vec{0}$, the points are collinear and do \emph{not} define a unique plane. Always test before computing. Also, if you rescale $\vec{n}$, recompute $d$ using the same point.
\end{attention}

\begin{aretenir}[Special planes]
\begin{itemize}
\item Plane through origin: $d=0$, equation $ax+by+cz=0$.
\item Plane parallel to $xy$-plane: $z=k$ (normal $\vec{k}$).
\item Plane parallel to $yz$-plane: $x=k$ (normal $\vec{i}$).
\item Plane parallel to $xz$-plane: $y=k$ (normal $\vec{j}$).
\item Plane parallel to a given plane $ax+by+cz=d_1$: same normal, equation $ax+by+cz=d_2$.
\end{itemize}
\end{aretenir}

\begin{exoresolu}[Special planes]
Find the equation of the plane through $P(2,-1,3)$ parallel to the plane $2x-3y+z=7$.
\tcblower
Parallel planes have the same normal $\vec{n}=(2,-3,1)$. The required plane has equation $2x-3y+z=d$. Substitute $P$: $d=2(2)-3(-1)+3=4+3+3=10$. Equation: $\boxed{2x-3y+z=10}$.
\end{exoresolu}

\coursec{Perpendicular Distance from a Point to a Plane}

\begin{propriete}[Perpendicular distance formula]
The perpendicular distance $\delta$ from point $P(x_1,y_1,z_1)$ to the plane $\Pi:\ ax+by+cz=d$ is
\[
\delta=\frac{|ax_1+by_1+cz_1-d|}{\sqrt{a^2+b^2+c^2}}.
\]
The numerator is the absolute value of substituting $P$ into the left-hand side minus $d$; the denominator is $|\vec{n}|$.
\end{propriete}

\begin{methode}[Distance and foot of the perpendicular]
To find the perpendicular distance from $P(x_1,y_1,z_1)$ to $\Pi:\ ax+by+cz=d$:
\begin{enumerate}
\item Ensure the equation is in the form $ax+by+cz-d=0$.
\item Apply $\delta=\dfrac{|ax_1+by_1+cz_1-d|}{\sqrt{a^2+b^2+c^2}}$.
\item \textbf{Foot of the perpendicular $F$:} parametrise the line through $P$ along $\vec{n}$:
\[
x=x_1+at,\quad y=y_1+bt,\quad z=z_1+ct.
\]
Substitute into the plane equation, solve for $t$, then find $F$.
\item \textbf{Reflex:} if $P$ lies on $\Pi$, the numerator is zero — quick consistency check.
\end{enumerate}
\end{methode}

\begin{exoresolu}[Distance and foot of the perpendicular]
Find the perpendicular distance from $P(1,2,3)$ to $\Pi:\ 2x-y+2z=5$. Hence find the foot $F$.
\tcblower
\textbf{Distance.} $a=2$, $b=-1$, $c=2$, $d=5$.
\[
\delta=\frac{|2(1)-(2)+2(3)-5|}{\sqrt{2^2+(-1)^2+2^2}}=\frac{|2-2+6-5|}{\sqrt{9}}=\frac{1}{3}.
\]
Distance: $\boxed{\tfrac{1}{3}\text{ unit}}$.

\textbf{Foot of the perpendicular.} Normal $\vec{n}=(2,-1,2)$. Line through $P$:
\[
x=1+2t,\quad y=2-t,\quad z=3+2t.
\]
Substitute into $2x-y+2z=5$:
\[
2(1+2t)-(2-t)+2(3+2t)=5 \;\Longrightarrow\; 2+4t-2+t+6+4t=5 \;\Longrightarrow\; 9t=-1 \;\Longrightarrow\; t=-\tfrac{1}{9}.
\]
\[
F=\left(1-\tfrac{2}{9},\;2+\tfrac{1}{9},\;3-\tfrac{2}{9}\right)=\boxed{\left(\tfrac{7}{9},\tfrac{19}{9},\tfrac{25}{9}\right)}.
\]
\textbf{Check:} $|PF|=|t|\,|\vec{n}|=\tfrac{1}{9}\times 3=\tfrac13$ \checkmark.
\end{exoresolu}

\begin{exoresolu}[Distance between parallel planes]
Find the distance between the parallel planes $\Pi_1:\ 2x-y+2z=5$ and $\Pi_2:\ 2x-y+2z=11$.
\tcblower
Pick any point on $\Pi_1$, e.g. $A(0,0,2.5)$ (since $2(0)-0+2(2.5)=5$). Distance from $A$ to $\Pi_2$:
\[
\delta=\frac{|2(0)-0+2(2.5)-11|}{\sqrt{9}}=\frac{|5-11|}{3}=\frac{6}{3}=2.
\]
Alternatively, for parallel planes $ax+by+cz=d_1$ and $ax+by+cz=d_2$:
\[
\delta=\frac{|d_2-d_1|}{\sqrt{a^2+b^2+c^2}}=\frac{|11-5|}{3}=2.
\]
Distance: $\boxed{2\text{ units}}$.
\end{exoresolu}

\coursec{Angles: Between Two Planes and Between a Line and a Plane}

\begin{propriete}[Angle between two planes]
The angle $\theta$ between two planes is the acute angle between their normal vectors $\vec{n}_1$ and $\vec{n}_2$:
\[
\cos\theta=\frac{|\vec{n}_1\cdot\vec{n}_2|}{|\vec{n}_1|\,|\vec{n}_2|},\qquad 0^\circ\le\theta\le 90^\circ.
\]
\begin{itemize}
\item Planes parallel $\Leftrightarrow$ $\vec{n}_1\parallel\vec{n}_2$ $\Leftrightarrow$ $\theta=0^\circ$.
\item Planes perpendicular $\Leftrightarrow$ $\vec{n}_1\cdot\vec{n}_2=0$ $\Leftrightarrow$ $\theta=90^\circ$.
\end{itemize}
\end{propriete}

\begin{propriete}[Angle between a line and a plane]
Let a line have direction vector $\vec{d}$ and a plane have normal $\vec{n}$. The angle $\theta$ between the line and the plane is the complement of the angle $\varphi$ between $\vec{d}$ and $\vec{n}$:
\[
\sin\theta=\frac{|\vec{d}\cdot\vec{n}|}{|\vec{d}|\,|\vec{n}|},\qquad 0^\circ\le\theta\le 90^\circ.
\]
\begin{itemize}
\item Line parallel to plane (or lies in it) $\Leftrightarrow$ $\vec{d}\cdot\vec{n}=0$ $\Leftrightarrow$ $\theta=0^\circ$.
\item Line perpendicular to plane $\Leftrightarrow$ $\vec{d}\parallel\vec{n}$ $\Leftrightarrow$ $\theta=90^\circ$.
\end{itemize}
\end{propriete}

\begin{methode}[Angles involving planes]
\begin{enumerate}
\item \textbf{Two planes:} identify normals $\vec{n}_1,\vec{n}_2$ from Cartesian equations; compute $\cos\theta=\dfrac{|\vec{n}_1\cdot\vec{n}_2|}{|\vec{n}_1||\vec{n}_2|}$.
\item \textbf{Line and plane:} identify direction $\vec{d}$ of line and normal $\vec{n}$ of plane; compute $\sin\theta=\dfrac{|\vec{d}\cdot\vec{n}|}{|\vec{d}||\vec{n}|}$.
\item \textbf{Reflexes:} parallel $\Leftrightarrow$ dot product zero (line-plane) or normals parallel (plane-plane); perpendicular $\Leftrightarrow$ dot product zero (plane-plane) or direction parallel to normal (line-plane).
\item \textbf{Trap:} line–plane angle uses \cle{sine}, not cosine.
\end{enumerate}
\end{methode}

\begin{exoresolu}[Two angle computations]
(a) Find the acute angle between $\Pi_1:\ x+2y+2z=3$ and $\Pi_2:\ 2x-2y+z=6$.

(b) Find the angle between the line $\vec{r}=(\vec{i}+2\vec{k})+t(2\vec{i}+\vec{j}-2\vec{k})$ and $\Pi_1$.
\tcblower
\textbf{(a)} Normals: $\vec{n}_1=(1,2,2)$, $\vec{n}_2=(2,-2,1)$.
\[
\vec{n}_1\cdot\vec{n}_2=2-4+2=0 \;\Longrightarrow\; \theta=90^\circ.
\]
Planes are perpendicular: $\boxed{\theta=90^\circ}$.

\textbf{(b)} Direction $\vec{d}=(2,1,-2)$; normal $\vec{n}=(1,2,2)$.
\[
\vec{d}\cdot\vec{n}=2+2-4=0 \;\Longrightarrow\; \sin\theta=0 \;\Longrightarrow\; \theta=0^\circ.
\]
The line is parallel to the plane. Check if it lies in the plane: point $(1,0,2)$ gives $1+0+4=5\neq3$, so line is parallel and distinct. Angle: $\boxed{0^\circ}$.

\medskip
\textbf{Non-degenerate practice:} line with $\vec{d}=(1,1,1)$ and $\Pi_1$.
\[
\sin\theta=\frac{|1+2+2|}{\sqrt{3}\,\sqrt{9}}=\frac{5}{3\sqrt{3}}=\frac{5\sqrt{3}}{9}\approx 0.9623 \;\Longrightarrow\; \theta\approx 74.2^\circ.
\]
\end{exoresolu}

\begin{popfigure}
\begin{tikzpicture}[scale=0.9]
% Angle between line and plane
\draw[popblue, thick, fill=popblue!10] (0,0) -- (4,0) -- (4.5,1.5) -- (0.5,1.5) -- cycle;
\draw[->, thick, poporange] (2,0.75) -- (2,2) node[above] {$\vec{n}$};
\draw[->, thick, popgreen] (2,0.75) -- (3.5,1.5) node[right] {$\vec{d}$};
\draw[dashed, popmuted] (2,0.75) -- (3.5,0.75) node[right] {proj};
\draw[poppurple, <->] (2.3,0.9) arc (0:45:0.5) node[midway, right] {$\varphi$};
\draw[poppurple, <->] (3.2,1.1) arc (180:135:0.5) node[midway, above] {$\theta$};
\node at (2,-0.5) {Plane $\Pi$};
\end{tikzpicture}
\legende{Angle $\theta$ between line (direction $\vec{d}$) and plane (normal $\vec{n}$): $\theta=90^\circ-\varphi$, $\sin\theta=|\vec{d}\cdot\vec{n}|/(|\vec{d}||\vec{n}|)$}
\end{popfigure}

\coursec{Intersection of a Line with a Plane and of Two Planes}

\begin{methode}[Intersection of a line and a plane]
Given line $\vec{r}=\vec{a}+t\vec{d}$ and plane $ax+by+cz=d$:
\begin{enumerate}
\item Write parametric coordinates $x(t),y(t),z(t)$.
\item Substitute into plane equation $\rightarrow$ linear equation in $t$.
\item \textbf{Unique solution $t_0$:} line pierces plane at $\vec{a}+t_0\vec{d}$.
\item \textbf{No solution (e.g. $0=5$):} line parallel to plane, no intersection.
\item \textbf{Identity (e.g. $0=0$):} line lies entirely in the plane.
\item \textbf{Quick check:} compute $\vec{d}\cdot\vec{n}$. If zero, line $\parallel$ plane — then test if $\vec{a}$ satisfies plane equation to distinguish cases 4 and 5.
\end{enumerate}
\end{methode}

\begin{exoresolu}[Line meeting a plane]
Find the intersection of $\vec{r}=\begin{pmatrix}1\\-1\\2\end{pmatrix}+t\begin{pmatrix}2\\1\\-1\end{pmatrix}$ with $x-2y+3z=9$.
\tcblower
\textbf{Parametric:} $x=1+2t$, $y=-1+t$, $z=2-t$.

\textbf{Substitute:} $(1+2t)-2(-1+t)+3(2-t)=9 \;\Longrightarrow\; 1+2t+2-2t+6-3t=9 \;\Longrightarrow\; 9-3t=9 \;\Longrightarrow\; t=0$.

\textbf{Point:} $\boxed{(1,-1,2)}$.

\textbf{Check:} $\vec{d}\cdot\vec{n}=(2)(1)+(1)(-2)+(-1)(3)=-3\neq0$ $\rightarrow$ unique intersection expected. \checkmark
\end{exoresolu}

\begin{exoresolu}[Line parallel to plane]
Determine whether the line $\vec{r}=\begin{pmatrix}1\\2\\-1\end{pmatrix}+t\begin{pmatrix}1\\-2\\1\end{pmatrix}$ meets the plane $x-2y+z=4$.
\tcblower
$\vec{d}=(1,-2,1)$, $\vec{n}=(1,-2,1)$ $\rightarrow$ $\vec{d}\cdot\vec{n}=1+4+1=6\neq0$? Wait: $\vec{d}\cdot\vec{n}=1(1)+(-2)(-2)+1(1)=1+4+1=6\neq0$. So line is \emph{not} parallel to plane; it intersects.

Let's do a parallel case: line $\vec{r}=\begin{pmatrix}1\\2\\-1\end{pmatrix}+t\begin{pmatrix}1\\1\\1\end{pmatrix}$ with plane $x+y+z=3$.
$\vec{d}\cdot\vec{n}=1+1+1=3\neq0$ — still intersects.

True parallel: line $\vec{r}=\begin{pmatrix}1\\2\\-1\end{pmatrix}+t\begin{pmatrix}1\\-1\\0\end{pmatrix}$ with plane $x+y+z=3$.
$\vec{d}\cdot\vec{n}=1-1+0=0$ $\rightarrow$ parallel. Point $(1,2,-1)$: $1+2-1=2\neq3$ $\rightarrow$ no intersection.
\end{exoresolu}

\begin{methode}[Line of intersection of two planes]
Given $\Pi_1:\ a_1x+b_1y+c_1z=d_1$ and $\Pi_2:\ a_2x+b_2y+c_2z=d_2$:
\begin{enumerate}
\item Compare normals $\vec{n}_1,\vec{n}_2$: if parallel, planes are parallel (no line, or coincident).
\item Direction of intersection line: $\vec{d}=\vec{n}_1\times\vec{n}_2$ (perpendicular to both normals).
\item Find one point on the line: set one coordinate to a convenient value (often $z=0$) and solve the $2\times2$ system. If that fails, try $x=0$ or $y=0$.
\item Write line as $\vec{r}=\vec{p}+t\vec{d}$.
\item \textbf{Alternative:} solve the system with one variable as parameter.
\end{enumerate}
\end{methode}

\begin{exoresolu}[Line of intersection]
Find the vector equation of the line of intersection of
\[
\Pi_1:\ x+y+z=6,\qquad \Pi_2:\ 2x-y+z=3.
\]
\tcblower
\textbf{Normals:} $\vec{n}_1=(1,1,1)$, $\vec{n}_2=(2,-1,1)$ — not parallel.

\textbf{Direction:}
\[
\vec{d}=\vec{n}_1\times\vec{n}_2=\begin{vmatrix}
\vec{i} & \vec{j} & \vec{k}\\
1 & 1 & 1\\
2 & -1 & 1
\end{vmatrix}
=\vec{i}(1+1)-\vec{j}(1-2)+\vec{k}(-1-2)=\begin{pmatrix}2\\1\\-3\end{pmatrix}.
\]

\textbf{Point on line:} set $z=0$:
\[
\begin{cases}
x+y=6\\
2x-y=3
\end{cases}
\;\Longrightarrow\; 3x=9 \;\Longrightarrow\; x=3,\; y=3.
\]
Point $P(3,3,0)$. Check: $3+3+0=6$ \checkmark; $6-3+0=3$ \checkmark.

\textbf{Line equation:}
\[
\boxed{\vec{r}=\begin{pmatrix}3\\3\\0\end{pmatrix}+t\begin{pmatrix}2\\1\\-3\end{pmatrix},\quad t\in\mathbb{R}}.
\]
\end{exoresolu}

\begin{exoresolu}[Intersection via parameter]
Find the line of intersection of $x+2y-z=4$ and $2x-y+3z=5$ by using $z$ as parameter.
\tcblower
Let $z=t$. System:
\[
\begin{cases}
x+2y=4+t\\
2x-y=5-3t
\end{cases}
\]
Solve: from second, $y=2x-5+3t$. Substitute into first:
\[
x+2(2x-5+3t)=4+t \;\Longrightarrow\; 5x-10+6t=4+t \;\Longrightarrow\; 5x=14-5t \;\Longrightarrow\; x=\tfrac{14}{5}-t.
\]
Then $y=2(\tfrac{14}{5}-t)-5+3t=\tfrac{28}{5}-2t-5+3t=\tfrac{3}{5}+t$.
Parametric: $x=\tfrac{14}{5}-t$, $y=\tfrac{3}{5}+t$, $z=t$.
Vector form: $\vec{r}=\begin{pmatrix}14/5\\3/5\\0\end{pmatrix}+t\begin{pmatrix}-1\\1\\1\end{pmatrix}$.
\end{exoresolu}

\coursec{Pencils of Planes and Synthesis Problems}

\begin{propriete}[Pencil of planes through the line of intersection]
All planes passing through the line of intersection of
\[
\Pi_1:\ a_1x+b_1y+c_1z=d_1,\qquad \Pi_2:\ a_2x+b_2y+c_2z=d_2
\]
are given by the one-parameter family
\[
(a_1x+b_1y+c_1z-d_1)+\lambda(a_2x+b_2y+c_2z-d_2)=0,\qquad \lambda\in\mathbb{R}.
\]
This represents every plane of the pencil \emph{except} $\Pi_2$ itself (which corresponds to $\lambda\to\infty$). To include $\Pi_2$, use the symmetric form $\mu\Pi_1+\lambda\Pi_2=0$ with $(\mu,\lambda)\neq(0,0)$.
\end{propriete}

\begin{methode}[Plane of a pencil satisfying an extra condition]
\begin{enumerate}
\item Write the pencil: $\Pi_1+\lambda\Pi_2=0$ (i.e. $(a_1x+b_1y+c_1z-d_1)+\lambda(a_2x+b_2y+c_2z-d_2)=0$).
\item Use the extra condition to find $\lambda$:
\begin{itemize}
\item \textbf{Passes through point $Q$:} substitute $Q$ into the pencil equation.
\item \textbf{Perpendicular to plane $\Pi_3$:} combined normal $\vec{n}=\vec{n}_1+\lambda\vec{n}_2$ must satisfy $\vec{n}\cdot\vec{n}_3=0$.
\item \textbf{Parallel to a line with direction $\vec{d}$:} $(\vec{n}_1+\lambda\vec{n}_2)\cdot\vec{d}=0$.
\item \textbf{Parallel to a plane $\Pi_3$:} $\vec{n}_1+\lambda\vec{n}_2$ parallel to $\vec{n}_3$.
\end{itemize}
\item Substitute $\lambda$ back and simplify.
\end{enumerate}
\end{methode}

\begin{exoresolu}[Plane of a pencil through a given point]
Find the plane through the line of intersection of $\Pi_1:\ x+2y-z=4$ and $\Pi_2:\ 2x-y+z=1$ that also passes through $Q(1,1,4)$.
\tcblower
\textbf{Pencil:} $(x+2y-z-4)+\lambda(2x-y+z-1)=0$.

\textbf{Condition $Q\in\Pi$:} $(1+2-4-4)+\lambda(2-1+4-1)=0 \;\Longrightarrow\; -5+4\lambda=0 \;\Longrightarrow\; \lambda=\tfrac54$.

\textbf{Substitute:} $(x+2y-z-4)+\tfrac54(2x-y+z-1)=0$. Multiply by $4$:
\[
4x+8y-4z-16+10x-5y+5z-5=0 \;\Longrightarrow\; \boxed{14x+3y+z=21}.
\]
\textbf{Check:} $Q$: $14+3+4=21$ \checkmark. Any point on the intersection line satisfies both $\Pi_1$ and $\Pi_2$, hence the combination.
\end{exoresolu}

\begin{exoresolu}[Plane of a pencil perpendicular to another plane]
Find the plane through the intersection of $\Pi_1:\ x+y+z=1$ and $\Pi_2:\ 2x-y+3z=4$ that is perpendicular to $\Pi_3:\ x-2y+z=5$.
\tcblower
\textbf{Pencil:} $(x+y+z-1)+\lambda(2x-y+3z-4)=0$.
Combined normal: $\vec{n}=(1+2\lambda,\;1-\lambda,\;1+3\lambda)$.
$\Pi_3$ normal: $\vec{n}_3=(1,-2,1)$.
Perpendicular $\Rightarrow$ $\vec{n}\cdot\vec{n}_3=0$:
\[
(1+2\lambda)(1)+(1-\lambda)(-2)+(1+3\lambda)(1)=0
\]
\[
1+2\lambda-2+2\lambda+1+3\lambda=0 \;\Longrightarrow\; 7\lambda=0 \;\Longrightarrow\; \lambda=0.
\]
Thus the required plane is $\Pi_1$ itself: $\boxed{x+y+z=1}$.
\end{exoresolu}

\begin{propriete}[Equation of planes through the point of intersection of two planes]
If two planes $\Pi_1$ and $\Pi_2$ intersect in a \emph{point} (which occurs when we have three planes meeting at a point, or when considering the intersection of two planes with a third condition), the family of planes through that point is not a simple pencil. However, if two planes intersect in a line $L$, and we want planes through a \emph{specific point} $P$ on $L$, we use the pencil method with the condition $P\in\Pi$.

If the problem asks for planes through the \emph{point of intersection of three planes} (i.e. the unique point where $\Pi_1,\Pi_2,\Pi_3$ meet), then any plane through that point has equation
\[
\alpha(a_1x+b_1y+c_1z-d_1)+\beta(a_2x+b_2y+c_2z-d_2)+\gamma(a_3x+b_3y+c_3z-d_3)=0
\]
with $(\alpha,\beta,\gamma)\neq(0,0,0)$. This is a two-parameter family.
\end{propriete}

\begin{exoresolu}[Planes through intersection point of three planes]
The planes $\Pi_1:\ x+y+z=6$, $\Pi_2:\ 2x-y+z=3$, $\Pi_3:\ x+2y-z=2$ intersect at a unique point $P$. Find the equation of the plane through $P$ parallel to the vector $\vec{v}=(1,1,1)$.
\tcblower
\textbf{Step 1: find $P$ by solving the system.}
\[
\begin{cases}
x+y+z=6 & (1)\\
2x-y+z=3 & (2)\\
x+2y-z=2 & (3)
\end{cases}
\]
$(1)+(3)$: $2x+3y=8$ $\Rightarrow$ $x=4-\tfrac32y$.
$(2)-(1)$: $x-2y=-3$ $\Rightarrow$ $x=2y-3$.
Equate: $4-\tfrac32y=2y-3$ $\Rightarrow$ $7=\tfrac72y$ $\Rightarrow$ $y=2$.
Then $x=1$, and from (1): $1+2+z=6$ $\Rightarrow$ $z=3$.
Point $P(1,2,3)$.

\textbf{Step 2: plane through $P$ parallel to $\vec{v}=(1,1,1)$.}
A plane parallel to a vector means the vector lies in the plane (is a direction vector). We need another direction or a normal.
Since only one direction is given, there are infinitely many such planes — they form a pencil around the line through $P$ parallel to $\vec{v}$.
General plane through $P$: $a(x-1)+b(y-2)+c(z-3)=0$ with normal $\vec{n}=(a,b,c)$.
Condition $\vec{n}\cdot\vec{v}=0$: $a+b+c=0$ $\Rightarrow$ $c=-a-b$.
Equation: $a(x-1)+b(y-2)-(a+b)(z-3)=0$.
This is a one-parameter family (e.g. set $a=1$, $b=\lambda$).
\end{exoresolu}

\begin{methode}[Plane defined by a point and two directions (synthesis)]
To find the plane through $A$ parallel to lines with directions $\vec{d}_1$ and $\vec{d}_2$ (non-parallel):
\begin{enumerate}
\item Compute normal $\vec{n}=\vec{d}_1\times\vec{d}_2$.
\item Form $ax+by+cz=d$ with $d=\vec{n}\cdot\vec{a}$.
\item If the plane must also \emph{contain} a given line, verify: (i) line's direction $\perp\vec{n}$, (ii) a point of the line satisfies the plane equation.
\end{enumerate}
\end{methode}

\begin{exoresolu}[GCE-style synthesis problem]
Lines $L_1$ and $L_2$:
\[
L_1:\ \vec{r}=\begin{pmatrix}1\\0\\1\end{pmatrix}+s\begin{pmatrix}1\\2\\1\end{pmatrix},\qquad
L_2:\ \vec{r}=\begin{pmatrix}2\\1\\0\end{pmatrix}+t\begin{pmatrix}2\\-1\\1\end{pmatrix}.
\]
(a) Find the Cartesian equation of plane $\Pi$ containing $L_1$ and parallel to $L_2$.

(b) Find the perpendicular distance from $B(2,1,0)$ on $L_2$ to $\Pi$, and interpret.
\tcblower
\textbf{(a)} $\Pi$ contains $\vec{d}_1=(1,2,1)$ and is parallel to $\vec{d}_2=(2,-1,1)$.
Normal: $\vec{n}=\vec{d}_1\times\vec{d}_2=\begin{vmatrix}
\vec{i} & \vec{j} & \vec{k}\\
1 & 2 & 1\\
2 & -1 & 1
\end{vmatrix}=\vec{i}(2+1)-\vec{j}(1-2)+\vec{k}(-1-4)=\begin{pmatrix}3\\1\\-5\end{pmatrix}$.
$\Pi$ contains $A(1,0,1)$: $d=3(1)+1(0)-5(1)=-2$.
Equation: $\boxed{3x+y-5z=-2}$.

\textbf{Check:} $\vec{d}_1\cdot\vec{n}=3+2-5=0$ \checkmark; $\vec{d}_2\cdot\vec{n}=6-1-5=0$ \checkmark; $A$: $3+0-5=-2$ \checkmark.

\textbf{(b)} Distance from $B(2,1,0)$:
\[
\delta=\frac{|3(2)+1-5(0)+2|}{\sqrt{9+1+25}}=\frac{|6+1+2|}{\sqrt{35}}=\frac{9}{\sqrt{35}}=\frac{9\sqrt{35}}{35}\approx 1.52.
\]
\textbf{Interpretation:} $L_2$ is parallel to $\Pi$ and $L_1\subset\Pi$. Every point of $L_2$ is at the same distance from $\Pi$; this common distance $\dfrac{9}{\sqrt{35}}$ is the \cle{shortest distance between the two skew lines} $L_1$ and $L_2$.
\end{exoresolu}

\begin{exoresolu}[Shortest distance between two skew lines]
Find the shortest distance between the skew lines
\[
L_1:\ \frac{x-1}{2}=\frac{y+1}{1}=\frac{z-2}{-1},\qquad
L_2:\ \frac{x-3}{1}=\frac{y}{2}=\frac{z+1}{1}.
\]
\tcblower
\textbf{Method:} distance = distance from any point on $L_2$ to the plane containing $L_1$ parallel to $L_2$.

$L_1$ point $A(1,-1,2)$, direction $\vec{d}_1=(2,1,-1)$.
$L_2$ point $B(3,0,-1)$, direction $\vec{d}_2=(1,2,1)$.
Normal to plane: $\vec{n}=\vec{d}_1\times\vec{d}_2=\begin{vmatrix}
\vec{i} & \vec{j} & \vec{k}\\
2 & 1 & -1\\
1 & 2 & 1
\end{vmatrix}=\vec{i}(1+2)-\vec{j}(2+1)+\vec{k}(4-1)=\begin{pmatrix}3\\-3\\3\end{pmatrix}\sim(1,-1,1)$.
Plane through $A$: $1(x-1)-1(y+1)+1(z-2)=0$ $\Rightarrow$ $x-y+z=4$.
Distance from $B(3,0,-1)$ to this plane:
\[
\delta=\frac{|3-0-1-4|}{\sqrt{3}}=\frac{2}{\sqrt{3}}=\frac{2\sqrt{3}}{3}.
\]
Shortest distance: $\boxed{\frac{2\sqrt{3}}{3}}$.
\end{exoresolu}

\coursec{Activities of Integration, Assessment and Remediation}

\begin{exercice}[Integration Exercise 1]
The planes $\Pi_1$ and $\Pi_2$ have equations $2x-y+3z=5$ and $x+2y-z=3$ respectively.
\begin{enumerate}
\item Find the acute angle between $\Pi_1$ and $\Pi_2$.
\item Find a vector equation of the line of intersection of $\Pi_1$ and $\Pi_2$.
\item Find the equation of the plane through the line of intersection of $\Pi_1$ and $\Pi_2$ which is perpendicular to the plane $x-y+2z=4$.
\end{enumerate}
\end{exercice}

\begin{exercice}[Integration Exercise 2]
A line $L$ has equation $\vec{r}=\begin{pmatrix}2\\-1\\3\end{pmatrix}+t\begin{pmatrix}1\\2\\-2\end{pmatrix}$. A plane $\Pi$ has equation $3x-y+2z=7$.
\begin{enumerate}
\item Show that $L$ is not parallel to $\Pi$ and find the coordinates of the intersection point.
\item Find the acute angle between $L$ and $\Pi$.
\item Find the coordinates of the foot of the perpendicular from the point $A(4,1,1)$ to $\Pi$.
\end{enumerate}
\end{exercice}

\begin{exercice}[Integration Exercise 3]
Three points are given: $A(1,2,-1)$, $B(3,0,1)$, $C(-1,1,2)$.
\begin{enumerate}
\item Show that $A,B,C$ are non-collinear.
\item Find the Cartesian equation of the plane $\Pi$ through $A,B,C$.
\item Find the perpendicular distance from the origin to $\Pi$.
\item Find the equation of the plane through the line of intersection of $\Pi$ and the plane $x+y+z=0$ that passes through the point $D(2,2,2)$.
\end{enumerate}
\end{exercice}

\begin{exercice}[Integration Exercise 4 — GCE Style]
The line $L_1$ passes through $A(1,2,3)$ and is parallel to $\vec{d}_1=2\vec{i}-\vec{j}+\vec{k}$.
The line $L_2$ passes through $B(4,-1,2)$ and is parallel to $\vec{d}_2=\vec{i}+2\vec{j}-\vec{k}$.
\begin{enumerate}
\item Find the Cartesian equation of the plane $\Pi$ that contains $L_1$ and is parallel to $L_2$.
\item Find the shortest distance between $L_1$ and $L_2$.
\item Find the coordinates of the points $P$ on $L_1$ and $Q$ on $L_2$ such that $\overrightarrow{PQ}$ is perpendicular to both lines.
\end{enumerate}
\end{exercice}

\begin{corrige}[Integration Exercise 1]
\begin{enumerate}
\item Normals: $\vec{n}_1=(2,-1,3)$, $\vec{n}_2=(1,2,-1)$.
$\vec{n}_1\cdot\vec{n}_2=2-2-3=-3$, $|\vec{n}_1|=\sqrt{14}$, $|\vec{n}_2|=\sqrt{6}$.
$\cos\theta=\frac{3}{\sqrt{84}}=\frac{3}{2\sqrt{21}}=\frac{\sqrt{21}}{14}$, $\theta\approx 70.5^\circ$.
\item Direction $\vec{d}=\vec{n}_1\times\vec{n}_2=\begin{vmatrix}\vec{i}&\vec{j}&\vec{k}\\2&-1&3\\1&2&-1\end{vmatrix}=(-5,5,5)\sim(-1,1,1)$.
Point on line: set $z=0$: $2x-y=5$, $x+2y=3$ $\Rightarrow$ $5x=13$, $x=13/5$, $y=1/5$.
Line: $\vec{r}=\begin{pmatrix}13/5\\1/5\\0\end{pmatrix}+t\begin{pmatrix}-1\\1\\1\end{pmatrix}$.
\item Pencil: $(2x-y+3z-5)+\lambda(x+2y-z-3)=0$.
Normal: $(2+\lambda,\,-1+2\lambda,\,3-\lambda)$.
Perpendicular to $x-y+2z=4$ $\Rightarrow$ $(2+\lambda)(1)+(-1+2\lambda)(-1)+(3-\lambda)(2)=0$
$\Rightarrow 2+\lambda+1-2\lambda+6-2\lambda=0 \Rightarrow 9-3\lambda=0 \Rightarrow \lambda=3$.
Plane: $(2x-y+3z-5)+3(x+2y-z-3)=0 \Rightarrow 5x+5y-14=0$ or $\boxed{x+y=\frac{14}{5}}$.
\end{enumerate}
\end{corrige}

\begin{corrige}[Integration Exercise 2]
\begin{enumerate}
\item $\vec{d}=(1,2,-2)$, $\vec{n}=(3,-1,2)$. $\vec{d}\cdot\vec{n}=3-2-4=-3\neq0$ $\rightarrow$ not parallel.
Parametric: $x=2+t$, $y=-1+2t$, $z=3-2t$.
Substitute: $3(2+t)-(-1+2t)+2(3-2t)=7 \Rightarrow 6+3t+1-2t+6-4t=7 \Rightarrow 13-3t=7 \Rightarrow t=2$.
Intersection: $(4,3,-1)$.
\item $\sin\theta=\frac{|\vec{d}\cdot\vec{n}|}{|\vec{d}||\vec{n}|}=\frac{3}{3\sqrt{14}}=\frac{1}{\sqrt{14}}$, $\theta=\arcsin(1/\sqrt{14})\approx 15.5^\circ$.
\item Foot from $A(4,1,1)$ to $3x-y+2z=7$:
Line: $x=4+3t$, $y=1-t$, $z=1+2t$.
$3(4+3t)-(1-t)+2(1+2t)=7 \Rightarrow 12+9t-1+t+2+4t=7 \Rightarrow 13+14t=7 \Rightarrow t=-3/7$.
Foot: $\left(4-\frac{9}{7},\,1+\frac{3}{7},\,1-\frac{6}{7}\right)=\boxed{\left(\frac{19}{7},\frac{10}{7},\frac{1}{7}\right)}$.
\end{enumerate}
\end{corrige}

\begin{corrige}[Integration Exercise 3]
\begin{enumerate}
\item $\overrightarrow{AB}=(2,-2,2)$, $\overrightarrow{AC}=(-2,-1,3)$. Cross product $\neq\vec{0}$ $\rightarrow$ non-collinear.
\item $\vec{n}=\overrightarrow{AB}\times\overrightarrow{AC}=\begin{vmatrix}\vec{i}&\vec{j}&\vec{k}\\2&-2&2\\-2&-1&3\end{vmatrix}=(-4,-10,-6)\sim(2,5,3)$.
Using $A$: $2(1)+5(2)+3(-1)=2+10-3=9$. Plane: $\boxed{2x+5y+3z=9}$.
\item Distance from origin: $\delta=\frac{|0-9|}{\sqrt{4+25+9}}=\frac{9}{\sqrt{38}}=\frac{9\sqrt{38}}{38}$.
\item Pencil: $(2x+5y+3z-9)+\lambda(x+y+z)=0$.
Through $D(2,2,2)$: $(4+10+6-9)+\lambda(2+2+2)=0 \Rightarrow 11+6\lambda=0 \Rightarrow \lambda=-11/6$.
Plane: $6(2x+5y+3z-9)-11(x+y+z)=0 \Rightarrow 12x+30y+18z-54-11x-11y-11z=0 \Rightarrow \boxed{x+19y+7z=54}$.
\end{enumerate}
\end{corrige}

\begin{corrige}[Integration Exercise 4]
\begin{enumerate}
\item $\vec{n}=\vec{d}_1\times\vec{d}_2=\begin{vmatrix}\vec{i}&\vec{j}&\vec{k}\\2&-1&1\\1&2&-1\end{vmatrix}=(-1,3,5)$.
Plane through $A(1,2,3)$: $-1(x-1)+3(y-2)+5(z-3)=0 \Rightarrow -x+1+3y-6+5z-15=0 \Rightarrow \boxed{-x+3y+5z=20}$ or $x-3y-5z=-20$.
\item Distance = distance from $B(4,-1,2)$ to plane: $\delta=\frac{|4-3(-1)-5(2)+20|}{\sqrt{1+9+25}}=\frac{|4+3-10+20|}{\sqrt{35}}=\frac{17}{\sqrt{35}}=\frac{17\sqrt{35}}{35}$.
\item $\overrightarrow{PQ}$ must be parallel to $\vec{n}=(-1,3,5)$.
$P=(1,2,3)+s(2,-1,1)$, $Q=(4,-1,2)+t(1,2,-1)$.
$\overrightarrow{PQ}=(3+t-2s,\,-3+2t+s,\,-1-t-s)=k(-1,3,5)$.
Solve: $3+t-2s=-k$, $-3+2t+s=3k$, $-1-t-s=5k$.
From first: $k=2s-t-3$. Substitute into second: $-3+2t+s=6s-3t-9 \Rightarrow 5t-5s=-6 \Rightarrow t-s=-6/5$.
Third: $-1-t-s=10s-5t-15 \Rightarrow 4t-11s=-14$.
Solve: $t=s-6/5$, $4(s-6/5)-11s=-14 \Rightarrow 4s-24/5-11s=-14 \Rightarrow -7s=-46/5 \Rightarrow s=46/35$, $t=46/35-42/35=4/35$.
$P=\left(1+\frac{92}{35},\,2-\frac{46}{35},\,3+\frac{46}{35}\right)=\left(\frac{127}{35},\frac{24}{35},\frac{151}{35}\right)$.
$Q=\left(4+\frac{4}{35},\,-1+\frac{8}{35},\,2-\frac{4}{35}\right)=\left(\frac{144}{35},-\frac{27}{35},\frac{66}{35}\right)$.
Check $\overrightarrow{PQ}=\frac{1}{35}(17,-51,-85)=\frac{17}{35}(1,-3,-5)$ parallel to $(-1,3,5)$ \checkmark.
\end{enumerate}
\end{corrige}

\begin{savaistu}[Historical note]
The vector approach to planes was systematised by Hermann Grassmann (1809–1877) and later refined by Josiah Willard Gibbs (1839–1903) and Oliver Heaviside (1850–1925). The "pencil of planes" concept comes from projective geometry (Poncelet, 1822). In Cameroon, vector methods are essential for engineering entrance exams (ENSP, ENSET) and for modelling terrain in surveying (e.g. mapping the slopes of Mount Cameroon).
\end{savaistu}

\begin{attention}[Common GCE traps]
\begin{itemize}
\item Forgetting the absolute value in the distance formula.
\item Using cosine instead of sine for line–plane angle.
\item Not checking collinearity of three points before finding a plane.
\item Confusing "plane through the line of intersection" (pencil, 1 parameter) with "plane through the point of intersection of three planes" (2 parameters).
\item Forgetting that the pencil $\Pi_1+\lambda\Pi_2=0$ misses $\Pi_2$; use symmetric form if $\Pi_2$ might be the answer.
\item In skew line distance, the vector connecting the closest points is parallel to $\vec{d}_1\times\vec{d}_2$.
\end{itemize}
\end{attention}

\begin{aretenir}[Key formulas summary]
\begin{center}
\begin{tabularx}{\linewidth}{|l|X|}\hline
\textbf{Situation} & \textbf{Formula / Condition} \\ \hline
Plane through $A(\vec{a})$, normal $\vec{n}$ & $\vec{r}\cdot\vec{n}=\vec{a}\cdot\vec{n}$ or $ax+by+cz=d$ \\ \hline
Plane through $A$, directions $\vec{u},\vec{v}$ & $\vec{r}=\vec{a}+\lambda\vec{u}+\mu\vec{v}$ \\ \hline
Plane through $A,B,C$ & $\vec{n}=\overrightarrow{AB}\times\overrightarrow{AC}$, then $d=\vec{n}\cdot\vec{a}$ \\ \hline
Distance point $P$ to plane $ax+by+cz=d$ & $\delta=\dfrac{|ax_1+by_1+cz_1-d|}{\sqrt{a^2+b^2+c^2}}$ \\ \hline
Angle between planes $\vec{n}_1,\vec{n}_2$ & $\cos\theta=\dfrac{|\vec{n}_1\cdot\vec{n}_2|}{|\vec{n}_1||\vec{n}_2|}$ \\ \hline
Angle line $\vec{d}$ and plane $\vec{n}$ & $\sin\theta=\dfrac{|\vec{d}\cdot\vec{n}|}{|\vec{d}||\vec{n}|}$ \\ \hline
Line–plane intersection & Substitute parametrics, solve for $t$ \\ \hline
Line of intersection of two planes & $\vec{d}=\vec{n}_1\times\vec{n}_2$, find one point \\ \hline
Pencil through intersection line & $\Pi_1+\lambda\Pi_2=0$ \\ \hline
Shortest distance between skew lines & Distance from point on $L_2$ to plane containing $L_1$ parallel to $L_2$ \\ \hline
\end{tabularx}
\end{center}
\end{aretenir}

\newpage

\coursec{Exercises}

\courssub{I check my knowledge}

\begin{exercice}[Multiple choice questions]
For each question, choose the single correct answer, justifying your choice briefly.
\begin{enumerate}
\item The plane $3x-2y+z=7$ has normal vector:
\begin{enumerate}
\item $(3,-2,1)$ \quad (b) $(3,2,1)$ \quad (c) $(7,-2,1)$ \quad (d) $(3,-2,7)$
\end{enumerate}
\item The perpendicular distance from the origin to the plane $2x+6y-3z=14$ is:
\begin{enumerate}
\item $1$ \quad (b) $2$ \quad (c) $7$ \quad (d) $14$
\end{enumerate}
\item The planes $x+y+z=1$ and $2x+2y+2z=5$ are:
\begin{enumerate}
\item identical \quad (b) parallel and distinct \quad (c) intersecting in a line \quad (d) perpendicular
\end{enumerate}
\item If a line has direction vector $\vec{d}$ and a plane has normal vector $\vec{n}$, then the line is parallel to the plane when:
\begin{enumerate}
\item $\vec{d}=\vec{n}$ \quad (b) $\vec{d}\cdot\vec{n}=0$ \quad (c) $\vec{d}\times\vec{n}=\vec{0}$ \quad (d) $\vec{d}\cdot\vec{n}=1$
\end{enumerate}
\end{enumerate}
\end{exercice}

\begin{exercice}[True or false]
State whether each assertion is true or false, and justify your answer.
\begin{enumerate}
\item Three collinear points determine a unique plane.
\item The angle between two planes is equal to the angle between their normal vectors (or its supplement).
\item If a line is perpendicular to a plane, then its direction vector is parallel to the normal vector of the plane.
\item Every plane of the pencil through the line of intersection of $\Pi_1$ and $\Pi_2$ has an equation of the form $\Pi_1+\lambda\,\Pi_2=0$.
\end{enumerate}
\end{exercice}

\begin{exercice}[Definitions and vocabulary]
\begin{enumerate}
\item Define the \cle{normal vector} of a plane and explain its geometric meaning.
\item Give the \cle{vector equation} and the \cle{parametric equations} of a plane passing through a point $A$ and parallel to two non-parallel vectors $\vec{u}$ and $\vec{v}$.
\item State the formula for the perpendicular distance from a point $P(x_0,y_0,z_0)$ to the plane $ax+by+cz+d=0$.
\item What is a \cle{pencil of planes}? Illustrate your answer with a sketch.
\end{enumerate}
\end{exercice}

\begin{exercice}[Direct calculations]
\begin{enumerate}
\item Write down a normal vector and a point of the plane $5x-y+4z=12$.
\item Find the Cartesian equation of the plane through $A(1,0,-2)$ with normal vector $\vec{n}=2\vec{i}-3\vec{j}+\vec{k}$.
\item Compute $\vec{u}\times\vec{v}$ where $\vec{u}=(1,2,-1)$ and $\vec{v}=(0,1,3)$, and interpret the result geometrically.
\item Determine whether the point $M(2,1,-1)$ lies on the plane $3x-2y+z=3$.
\end{enumerate}
\end{exercice}

\courssub{I practise}

\begin{exercice}[Plane through three points]
Find the Cartesian equation of the plane through $A(1,2,-1)$, $B(2,0,3)$ and $C(0,1,1)$. Verify first that the three points are not collinear.
\end{exercice}

\begin{exercice}[From parametric to Cartesian form]
A plane has parametric (vector) equation
\[
\vec{r}=(\vec{i}+2\vec{j}-\vec{k})+\lambda(\vec{i}-\vec{j})+\mu(\vec{j}+\vec{k}).
\]
Find its Cartesian equation and the coordinates of its intercepts with the three coordinate axes.
\end{exercice}

\begin{exercice}[Distance and foot of the perpendicular]
Calculate the perpendicular distance from $P(3,-1,2)$ to the plane $2x-y+2z=5$, and find the coordinates of the foot of the perpendicular from $P$ to the plane.
\end{exercice}

\begin{exercice}[Angles involving planes and lines]
\begin{enumerate}
\item Find the acute angle between the planes $x+2y-2z=3$ and $2x-y+z=7$.
\item Find the angle between the line $\vec{r}=(2\vec{i}+\vec{k})+\lambda(\vec{i}-2\vec{j}+2\vec{k})$ and the plane $3x+4z=1$.
\end{enumerate}
\end{exercice}

\courssub{I prepare for the GCE}

\begin{exercice}[Line, plane and pencil \hfill (12 marks)]
\begin{enumerate}
\item Determine the point of intersection of the line $\vec{r}=(\vec{i}-\vec{j}+2\vec{k})+\lambda(2\vec{i}+\vec{j}-\vec{k})$ with the plane $x+y+z=6$. \hfill (3 marks)
\item The line is now replaced by one through the same point but with direction $(\vec{i}+\vec{j}-2\vec{k})$. Show that this new line does not meet the plane, and state the geometrical relationship between the line and the plane. \hfill (3 marks)
\item Show that the planes $2x-y+z=4$ and $x+y-z=1$ intersect in a line $L$, and find a vector equation of $L$. \hfill (4 marks)
\item Find the acute angle between the two planes of part (c). \hfill (2 marks)
\end{enumerate}
\end{exercice}

\begin{exercice}[Pencils of planes \hfill (10 marks)]
\begin{enumerate}
\item Find the equation of the plane passing through the line of intersection of the planes $x+y+z=3$ and $2x-y+z=1$ and through the point $(1,2,3)$. \hfill (5 marks)
\item Find the equation of the plane through the line of intersection of the planes $x+2y-z=0$ and $3x-y+z=5$ which is perpendicular to the plane $x+y+z=2$. \hfill (5 marks)
\end{enumerate}
\end{exercice}

\begin{exercice}[GCE past-paper style: study of two planes \hfill (16 marks)]
The planes $\Pi_1:\ 2x-y+2z=9$ and $\Pi_2:\ x+2y-z=4$ are given.
\begin{enumerate}
\item Find the acute angle between $\Pi_1$ and $\Pi_2$, giving your answer to the nearest degree. \hfill (4 marks)
\item Find a vector equation of the line of intersection of $\Pi_1$ and $\Pi_2$. \hfill (5 marks)
\item Find the perpendicular distance from the origin to $\Pi_1$, and the coordinates of the foot of the perpendicular from the origin to $\Pi_1$. \hfill (4 marks)
\item Find the Cartesian equation of the plane through the line of intersection of $\Pi_1$ and $\Pi_2$ that passes through the point $(0,1,-1)$. \hfill (3 marks)
\end{enumerate}
\end{exercice}

\newpage

\coursec{Solutions to the exercises}

\courssub{I check my knowledge}

\begin{corrige}[Exercise 1 --- Multiple choice questions]
\begin{enumerate}
\item \textbf{Answer (a).} The Cartesian equation $ax+by+cz=d$ has normal vector $(a,b,c)$: the normal vector is read directly from the coefficients of $x$, $y$ and $z$. Here these coefficients are $3$, $-2$, $1$, so $\vec{n}=(3,-2,1)$. The constant term $7$ never appears in the normal vector, which eliminates (c) and (d); the sign of the $y$-coefficient eliminates (b).

\item \textbf{Answer (b).} Using the perpendicular distance formula with $P=(0,0,0)$ and the plane $2x+6y-3z-14=0$:
\[
d=\dfrac{|2(0)+6(0)-3(0)-14|}{\sqrt{2^{2}+6^{2}+(-3)^{2}}}=\dfrac{14}{\sqrt{4+36+9}}=\dfrac{14}{\sqrt{49}}=\dfrac{14}{7}=2.
\]
\textbf{Trap avoided:} do not forget the absolute value in the numerator, and remember that the denominator is the modulus of the normal vector, $\|\vec{n}\|=\sqrt{a^{2}+b^{2}+c^{2}}$, not $a+b+c$.

\item \textbf{Answer (b).} The normal vectors are $\vec{n}_1=(1,1,1)$ and $\vec{n}_2=(2,2,2)=2\vec{n}_1$: the normals are parallel, so the planes are parallel. They are not identical, since dividing the second equation by $2$ gives $x+y+z=\tfrac{5}{2}\neq 1$. Hence the planes are \cle{parallel and distinct} (they have no common point).

\item \textbf{Answer (b).} A line is parallel to a plane when its direction vector is perpendicular to the plane's normal vector, i.e. $\vec{d}\cdot\vec{n}=0$. (Option (c), $\vec{d}\times\vec{n}=\vec{0}$, would mean $\vec{d}\parallel\vec{n}$, i.e. the line \emph{perpendicular} to the plane.) Note that $\vec{d}\cdot\vec{n}=0$ covers two cases: the line is strictly parallel to the plane, or it lies entirely in the plane; to distinguish them, test whether one point of the line satisfies the equation of the plane.
\end{enumerate}
\end{corrige}

\begin{corrige}[Exercise 2 --- True or false]
\begin{enumerate}
\item \textbf{False.} If $A$, $B$, $C$ are collinear, then $\overrightarrow{AB}$ and $\overrightarrow{AC}$ are parallel, so $\overrightarrow{AB}\times\overrightarrow{AC}=\vec{0}$ and no unique normal vector exists: \emph{infinitely many} planes contain the line $(ABC)$ (they form the pencil of planes about that line). Three points determine a unique plane only when they are \textbf{non-collinear}.

\item \textbf{True.} The angle between two planes is defined as the acute angle between their normal vectors. Since $\vec{n}_1\cdot\vec{n}_2$ may be negative (giving an obtuse angle between the chosen normals), we take the absolute value:
\[
\cos\theta=\dfrac{|\vec{n}_1\cdot\vec{n}_2|}{\|\vec{n}_1\|\,\|\vec{n}_2\|},\qquad 0\leq\theta\leq 90^{\circ},
\]
which corresponds to the angle between the normals or its supplement, whichever is acute.

\item \textbf{True.} A line perpendicular to a plane is parallel to every normal direction of that plane; hence $\vec{d}=\lambda\vec{n}$ for some $\lambda\neq 0$ (equivalently $\vec{d}\times\vec{n}=\vec{0}$). This is the standard way to write the equation of the perpendicular from a point to a plane.

\item \textbf{True, with one caveat.} Every plane through the line of intersection of $\Pi_1=0$ and $\Pi_2=0$ can be written $\Pi_1+\lambda\Pi_2=0$ for some $\lambda\in\mathbb{R}$, \emph{except} the plane $\Pi_2$ itself (which would formally require ``$\lambda=\infty$''). In examinations, check separately whether $\Pi_2$ is a solution before concluding; if the required plane turns out to be $\Pi_2$, the equation in $\lambda$ will be inconsistent, which is the signal.
\end{enumerate}
\end{corrige}

\begin{corrige}[Exercise 3 --- Definitions and vocabulary]
\begin{enumerate}
\item A \cle{normal vector} of a plane $\Pi$ is any non-zero vector $\vec{n}$ perpendicular to every vector lying in $\Pi$ (equivalently, perpendicular to two non-parallel direction vectors of $\Pi$). Geometrically, $\vec{n}$ gives the direction of the line perpendicular to $\Pi$; it is unique up to a non-zero scalar multiple, so $(a,b,c)$ and $(2a,2b,2c)$ describe the same plane orientation.

\item Let $A$ have position vector $\vec{a}$, and let $\vec{u}$, $\vec{v}$ be two non-parallel direction vectors of the plane. The \cle{vector equation} is
\[
\vec{r}=\vec{a}+\lambda\vec{u}+\mu\vec{v},\qquad \lambda,\mu\in\mathbb{R}.
\]
If $A=(x_0,y_0,z_0)$, $\vec{u}=(u_1,u_2,u_3)$, $\vec{v}=(v_1,v_2,v_3)$, the \cle{parametric equations} are
\[
\begin{cases}
x=x_0+\lambda u_1+\mu v_1,\\
y=y_0+\lambda u_2+\mu v_2,\\
z=z_0+\lambda u_3+\mu v_3.
\end{cases}
\]
Each choice of the pair $(\lambda,\mu)$ gives exactly one point of the plane.

\item The perpendicular distance from $P(x_0,y_0,z_0)$ to the plane $ax+by+cz+d=0$ is
\[
d=\dfrac{|ax_0+by_0+cz_0+d|}{\sqrt{a^{2}+b^{2}+c^{2}}}.
\]
\textbf{Warning:} the equation must first be written with \emph{all} terms on one side; for the plane $ax+by+cz=k$ the numerator is $|ax_0+by_0+cz_0-k|$.

\item A \cle{pencil of planes} is the family of all planes passing through a fixed line $L$ (the \emph{axis} of the pencil). If $L$ is the intersection of $\Pi_1=0$ and $\Pi_2=0$, the pencil is $\Pi_1+\lambda\Pi_2=0$, $\lambda\in\mathbb{R}$ (together with $\Pi_2$ itself).
\begin{popfigure}
\begin{center}
\begin{tikzpicture}[scale=0.9]
\draw[thick,popink] (-3.2,0) -- (3.2,0) node[right] {\small $L$};
\draw[popblue,fill=popblueL,opacity=0.55] (-2.2,-1.4) -- (2.2,-1.4) -- (2.9,0.6) -- (-1.5,0.6) -- cycle;
\draw[poporange,fill=poporange,opacity=0.25] (-0.4,-1.9) -- (2.6,-0.5) -- (0.4,1.9) -- (-2.6,0.5) -- cycle;
\draw[popgreen,fill=popgreenL,opacity=0.55] (-2.6,-0.5) -- (0.4,-1.9) -- (2.6,0.5) -- (-0.4,1.9) -- cycle;
\draw[thick,popink] (-3.2,0) -- (3.2,0);
\node[popblue] at (2.6,-1.3) {\small $\Pi_1$};
\node[poporange] at (2.9,1.1) {\small $\Pi_2$};
\node[popgreen!60!black] at (-2.9,-1.0) {\small $\Pi_3$};
\end{tikzpicture}
\end{center}
\legende{Planes of a pencil all sharing the same axis $L$: each value of $\lambda$ gives one plane through $L$.}
\end{popfigure}
\end{enumerate}
\end{corrige}

\begin{corrige}[Exercise 4 --- Direct calculations]
\begin{enumerate}
\item Normal vector: $\vec{n}=(5,-1,4)$, read from the coefficients of $x$, $y$, $z$. A point: setting $x=y=0$ gives $4z=12$, so $(0,0,3)$ lies on the plane. (Any convenient choice of two coordinates works; choose values that make the remaining equation easy to solve.)

\item The plane through $A(1,0,-2)$ with normal $\vec{n}=(2,-3,1)$ consists of the points $M(x,y,z)$ such that $\vec{n}\cdot\overrightarrow{AM}=0$:
\[
2(x-1)-3(y-0)+1(z+2)=0 \;\Longrightarrow\; 2x-3y+z=0.
\]
\[
\[\boxed{2x-3y+z=0}\]
\]
\textbf{Check:} $A$ satisfies $2(1)-3(0)+(-2)=0$ \checkmark.

\item With $\vec{u}=(1,2,-1)$ and $\vec{v}=(0,1,3)$:
\[
\vec{u}\times\vec{v}=
\begin{pmatrix} (2)(3)-(-1)(1) \\ (-1)(0)-(1)(3) \\ (1)(1)-(2)(0) \end{pmatrix}
=\begin{pmatrix} 7 \\ -3 \\ 1 \end{pmatrix}.
\]
\textbf{Interpretation:} $(7,-3,1)$ is perpendicular to both $\vec{u}$ and $\vec{v}$ (verify: $7-2-3=0$ and $0-3+3=0$ \checkmark); it is a normal vector of any plane parallel to $\vec{u}$ and $\vec{v}$, and its magnitude $\sqrt{7^{2}+3^{2}+1^{2}}=\sqrt{59}$ is the area of the parallelogram spanned by $\vec{u}$ and $\vec{v}$.

\item Substituting $M(2,1,-1)$ into the left-hand side of $3x-2y+z=3$:
\[
3(2)-2(1)+(-1)=6-2-1=3.
\]
The equation is satisfied, so \textbf{$M$ lies on the plane}.
\end{enumerate}
\end{corrige}

\courssub{I practise}

\begin{corrige}[Exercise 5 --- Plane through three points]
\textbf{Step 1: check non-collinearity.}
\[
\overrightarrow{AB}=(1,-2,4),\qquad \overrightarrow{AC}=(-1,-1,2).
\]
These vectors are not proportional (e.g. $\tfrac{1}{-1}\neq\tfrac{-2}{-1}$), so $A$, $B$, $C$ are \textbf{not collinear} and determine a unique plane. This check is essential: with collinear points the cross product in Step 2 would give $\vec{0}$.

\textbf{Step 2: normal vector.}
\[
\vec{n}=\overrightarrow{AB}\times\overrightarrow{AC}
=\begin{pmatrix} (-2)(2)-(4)(-1) \\ (4)(-1)-(1)(2) \\ (1)(-1)-(-2)(-1) \end{pmatrix}
=\begin{pmatrix} -4+4 \\ -4-2 \\ -1-2 \end{pmatrix}
=\begin{pmatrix} 0 \\ -6 \\ -3 \end{pmatrix}
=-3\begin{pmatrix} 0 \\ 2 \\ 1 \end{pmatrix}.
\]
Take the simpler normal $\vec{n}=(0,2,1)$ (any non-zero multiple is acceptable).

\textbf{Step 3: equation.} Using point $A(1,2,-1)$ and $\vec{n}\cdot\overrightarrow{AM}=0$:
\[
0(x-1)+2(y-2)+1(z+1)=0 \;\Longrightarrow\; \boxed{2y+z=3}.
\]
\textbf{Verification:} $B$: $2(0)+3=3$ \checkmark; $C$: $2(1)+1=3$ \checkmark. Note that the equation contains no $x$-term: the plane is parallel to the $x$-axis, which is consistent with the normal vector having a zero first component.
\end{corrige}

\begin{corrige}[Exercise 6 --- From parametric to Cartesian form]
The plane passes through $A(1,2,-1)$ with direction vectors $\vec{u}=(1,-1,0)$ and $\vec{v}=(0,1,1)$.

\textbf{Normal vector:}
\[
\vec{n}=\vec{u}\times\vec{v}
=\begin{pmatrix} (-1)(1)-(0)(1) \\ (0)(0)-(1)(1) \\ (1)(1)-(-1)(0) \end{pmatrix}
=\begin{pmatrix} -1 \\ -1 \\ 1 \end{pmatrix}.
\]

\textbf{Cartesian equation:} for $M(x,y,z)$ on the plane, $\vec{n}\cdot\overrightarrow{AM}=0$:
\[
-1(x-1)-1(y-2)+1(z+1)=0
\;\Longrightarrow\; -x-y+z+4=0
\;\Longrightarrow\; \boxed{x+y-z=4}.
\]
\textbf{Check with the parametric form:} substituting $x=1+\lambda$, $y=2-\lambda+\mu$, $z=-1+\mu$ gives $(1+\lambda)+(2-\lambda+\mu)-(-1+\mu)=4$ for all $\lambda,\mu$ \checkmark.

\textbf{Intercepts with the axes:}
\begin{itemize}
\item $x$-axis ($y=z=0$): $x=4$, giving $(4,0,0)$;
\item $y$-axis ($x=z=0$): $y=4$, giving $(0,4,0)$;
\item $z$-axis ($x=y=0$): $-z=4$, giving $(0,0,-4)$.
\end{itemize}
These three intercepts are a convenient way to sketch the plane: it is the plane through the triangle with vertices $(4,0,0)$, $(0,4,0)$, $(0,0,-4)$.
\end{corrige}

\begin{corrige}[Exercise 7 --- Distance and foot of the perpendicular]
Plane: $2x-y+2z-5=0$, normal $\vec{n}=(2,-1,2)$, $\|\vec{n}\|=\sqrt{4+1+4}=3$.

\textbf{Distance:}
\[
d=\dfrac{|2(3)-(-1)+2(2)-5|}{3}=\dfrac{|6+1+4-5|}{3}=\dfrac{6}{3}=\boxed{2}.
\]

\textbf{Foot of the perpendicular.} The perpendicular through $P(3,-1,2)$ is parallel to $\vec{n}$, so it has equation
\[
\vec{r}=(3,-1,2)+t(2,-1,2),\qquad t\in\mathbb{R}.
\]
Substituting $x=3+2t$, $y=-1-t$, $z=2+2t$ into the plane:
\[
2(3+2t)-(-1-t)+2(2+2t)=5
\;\Longrightarrow\; 6+4t+1+t+4+4t=5
\;\Longrightarrow\; 9t=-6 \;\Longrightarrow\; t=-\tfrac{2}{3}.
\]
Hence the foot is
\[
H=\left(3-\tfrac{4}{3},\;-1+\tfrac{2}{3},\;2-\tfrac{4}{3}\right)=\boxed{\left(\tfrac{5}{3},\;-\tfrac{1}{3},\;\tfrac{2}{3}\right)}.
\]
\textbf{Check:} $2\cdot\tfrac{5}{3}+\tfrac{1}{3}+2\cdot\tfrac{2}{3}=\tfrac{10+1+4}{3}=5$ \checkmark, and $|PH|=|t|\,\|\vec{n}\|=\tfrac{2}{3}\times 3=2$ \checkmark, which agrees with the distance found above. This double check is quick and strongly recommended in the examination.
\end{corrige}

\begin{corrige}[Exercise 8 --- Angles involving planes and lines]
\begin{enumerate}
\item Normals: $\vec{n}_1=(1,2,-2)$, $\|\vec{n}_1\|=\sqrt{1+4+4}=3$; $\vec{n}_2=(2,-1,1)$, $\|\vec{n}_2\|=\sqrt{4+1+1}=\sqrt{6}$.
\[
\vec{n}_1\cdot\vec{n}_2=2-2-2=-2,\qquad
\cos\theta=\dfrac{|-2|}{3\sqrt{6}}=\dfrac{2}{3\sqrt{6}}\approx 0.2722.
\]
\[
\boxed{\theta\approx 74.2^{\circ}}.
\]
The absolute value is essential here: without it, $\cos\theta<0$ would give the obtuse angle $105.8^{\circ}$ between the chosen normals instead of the acute angle between the planes.

\item Line direction $\vec{d}=(1,-2,2)$, $\|\vec{d}\|=3$; plane normal $\vec{n}=(3,0,4)$, $\|\vec{n}\|=5$. The angle $\theta$ between a line and a plane is the complement of the angle between $\vec{d}$ and $\vec{n}$, so it satisfies $\sin\theta=\dfrac{|\vec{d}\cdot\vec{n}|}{\|\vec{d}\|\,\|\vec{n}\|}$:
\[
\vec{d}\cdot\vec{n}=3+0+8=11,\qquad
\sin\theta=\dfrac{11}{15}\approx 0.7333
\;\Longrightarrow\; \boxed{\theta\approx 47.2^{\circ}}.
\]
(Recall: use $\sin$ here because $\theta=90^{\circ}-\varphi$, where $\varphi$ is the angle between $\vec{d}$ and $\vec{n}$, and $\cos\varphi=\sin\theta$.)
\end{enumerate}
\end{corrige}

\courssub{I prepare for the GCE}

\begin{corrige}[Exercise 9 --- Line, plane and pencil (12 marks)]
\begin{enumerate}
\item Parametrize the line: $x=1+2\lambda$, $y=-1+\lambda$, $z=2-\lambda$. Substitute into $x+y+z=6$:
\[
(1+2\lambda)+(-1+\lambda)+(2-\lambda)=6
\;\Longrightarrow\; 2+2\lambda=6
\;\Longrightarrow\; \lambda=2.
\]
Point of intersection: $\boxed{(5,1,0)}$. \hfill \emph{[M1 substitution, A1 value of $\lambda$, A1 point]}

\item New direction $\vec{d}=(1,1,-2)$, plane normal $\vec{n}=(1,1,1)$:
\[
\vec{d}\cdot\vec{n}=1+1-2=0,
\]
so the line is parallel to the plane. Moreover the point $(1,-1,2)$ of the line gives $1-1+2=2\neq 6$, so this point is \emph{not} on the plane. Hence the line \textbf{does not meet the plane}: it is \cle{strictly parallel} to it. \hfill \emph{[M1 dot product, A1 conclusion of parallelism, A1 point not on plane]}

\item Normals $\vec{n}_1=(2,-1,1)$ and $\vec{n}_2=(1,1,-1)$ are not proportional (e.g. $\tfrac{2}{1}\neq\tfrac{-1}{1}$), so the planes \textbf{intersect in a line} $L$. The direction of $L$ is perpendicular to both normals:
\[
\vec{n}_1\times\vec{n}_2=
\begin{pmatrix} (-1)(-1)-(1)(1) \\ (1)(1)-(2)(-1) \\ (2)(1)-(-1)(1) \end{pmatrix}
=\begin{pmatrix} 1-1 \\ 1+2 \\ 2+1 \end{pmatrix}
=\begin{pmatrix} 0 \\ 3 \\ 3 \end{pmatrix}
=3\begin{pmatrix} 0 \\ 1 \\ 1 \end{pmatrix}.
\]
A point of $L$: set $z=0$; then $2x-y=4$ and $x+y=1$. Adding gives $3x=5$, so $x=\tfrac{5}{3}$, $y=1-\tfrac{5}{3}=-\tfrac{2}{3}$. Hence
\[
\[\boxed{\vec{r}=\begin{pmatrix} 5/3 \\ -2/3 \\ 0 \end{pmatrix}+\lambda\begin{pmatrix} 0 \\ 1 \\ 1 \end{pmatrix},\quad \lambda\in\mathbb{R}.}\]
\]
\textbf{Check:} the point satisfies $2(\tfrac{5}{3})-(-\tfrac{2}{3})+0=4$ \checkmark\ and $\tfrac{5}{3}+(-\tfrac{2}{3})-0=1$ \checkmark. \hfill \emph{[B1 normals not parallel, M1 cross product, A1 a common point, A1 equation]}

\item $\vec{n}_1\cdot\vec{n}_2=2-1-1=0$: the normals are perpendicular, so the planes are perpendicular:
\[
\boxed{\theta=90^{\circ}}.
\]
\hfill \emph{[M1 dot product of normals, A1 answer]}
\end{enumerate}
\textbf{Examiner's expectations:} clear substitution in (a); in (b) \emph{both} conditions ($\vec{d}\cdot\vec{n}=0$ and the point not on the plane) are required for full marks --- a line with $\vec{d}\cdot\vec{n}=0$ could also lie \emph{in} the plane; in (c), if setting $z=0$ had led to an inconsistent system, one would try $x=0$ or $y=0$ instead (the line of intersection may be parallel to the plane $z=0$).
\end{corrige}

\begin{corrige}[Exercise 10 --- Pencils of planes (10 marks)]
\begin{enumerate}
\item The pencil through the line of intersection of $x+y+z=3$ and $2x-y+z=1$ is
\[
(x+y+z-3)+\lambda(2x-y+z-1)=0.
\]
Imposing passage through $(1,2,3)$:
\[
(1+2+3-3)+\lambda(2-2+3-1)=0
\;\Longrightarrow\; 3+2\lambda=0
\;\Longrightarrow\; \lambda=-\tfrac{3}{2}.
\]
Substituting and multiplying by $2$ to clear the fraction:
\[
2(x+y+z-3)-3(2x-y+z-1)=0
\;\Longrightarrow\; -4x+5y-z-3=0,
\]
\[
\boxed{4x-5y+z+3=0}.
\]
\textbf{Check:} $4(1)-5(2)+3+3=0$ \checkmark. \hfill \emph{[M1 pencil, M1 substitution, A1 $\lambda$, M1 simplification, A1 equation]}

\item The pencil through the line of intersection of $x+2y-z=0$ and $3x-y+z=5$ is
\[
(x+2y-z)+\lambda(3x-y+z-5)=0,
\]
with normal vector $\vec{n}=(1+3\lambda,\;2-\lambda,\;-1+\lambda)$. Perpendicularity to the plane $x+y+z=2$ (normal $(1,1,1)$) requires $\vec{n}\cdot(1,1,1)=0$:
\[
(1+3\lambda)+(2-\lambda)+(-1+\lambda)=0
\;\Longrightarrow\; 2+3\lambda=0
\;\Longrightarrow\; \lambda=-\tfrac{2}{3}.
\]
Substituting and multiplying by $3$:
\[
3(x+2y-z)-2(3x-y+z-5)=0
\;\Longrightarrow\; -3x+8y-5z+10=0,
\]
\[
\boxed{3x-8y+5z=10}.
\]
\textbf{Check:} normals $(3,-8,5)\cdot(1,1,1)=3-8+5=0$ \checkmark. \hfill \emph{[M1 pencil, M1 normal in terms of $\lambda$, M1 perpendicularity condition, A1 $\lambda$, A1 equation]}
\end{enumerate}
\textbf{Key idea:} in a pencil problem, the extra condition (passage through a point, perpendicularity or parallelism to a given plane or line) is always translated into \emph{one} equation in the single unknown $\lambda$.
\end{corrige}

\begin{corrige}[Exercise 11 --- GCE past-paper style: study of two planes (16 marks)]
Given $\Pi_1: 2x-y+2z=9$ with $\vec{n}_1=(2,-1,2)$, $\|\vec{n}_1\|=\sqrt{4+1+4}=3$, and $\Pi_2: x+2y-z=4$ with $\vec{n}_2=(1,2,-1)$, $\|\vec{n}_2\|=\sqrt{1+4+1}=\sqrt{6}$.

\begin{enumerate}
\item \textbf{Angle between the planes.}
\[
\vec{n}_1\cdot\vec{n}_2=2-2-2=-2,\qquad
\cos\theta=\dfrac{|-2|}{3\sqrt{6}}=\dfrac{2}{3\sqrt{6}}\approx 0.2722.
\]
\[
\boxed{\theta\approx 74^{\circ}\ \text{(nearest degree)}}.
\]
\hfill \emph{[B1 both normals and moduli, M1 correct formula with absolute value, A1 $\cos\theta$, A1 rounded answer]}

\item \textbf{Line of intersection.} The normals are not proportional, so the planes meet in a line $L$ whose direction is perpendicular to both normals:
\[
\vec{n}_1\times\vec{n}_2=
\begin{pmatrix} (-1)(-1)-(2)(2) \\ (2)(1)-(2)(-1) \\ (2)(2)-(-1)(1) \end{pmatrix}
=\begin{pmatrix} 1-4 \\ 2+2 \\ 4+1 \end{pmatrix}
=\begin{pmatrix} -3 \\ 4 \\ 5 \end{pmatrix}.
\]
A common point: set $z=0$; then $2x-y=9$ and $x+2y=4$. From the first equation, $y=2x-9$; substituting into the second: $x+2(2x-9)=4$, so $5x=22$, $x=\tfrac{22}{5}$, $y=\tfrac{44}{5}-9=-\tfrac{1}{5}$.
\[
\[\boxed{\vec{r}=\begin{pmatrix} 22/5 \\ -1/5 \\ 0 \end{pmatrix}+\lambda\begin{pmatrix} -3 \\ 4 \\ 5 \end{pmatrix},\quad \lambda\in\mathbb{R}.}\]
\]
\textbf{Check:} $2(\tfrac{22}{5})-(-\tfrac{1}{5})=9$ \checkmark\ and $\tfrac{22}{5}+2(-\tfrac{1}{5})=4$ \checkmark. \hfill \emph{[M1 cross product, A1 direction, M1 fixing a coordinate, A1 point, A1 full equation]}

\item \textbf{Distance from the origin to $\Pi_1$:}
\[
d=\dfrac{|2(0)-(0)+2(0)-9|}{3}=\dfrac{9}{3}=\boxed{3}.
\]
\textbf{Foot of the perpendicular:} the perpendicular through $O$ has direction $\vec{n}_1$, so $\vec{r}=t(2,-1,2)$. Substituting into $\Pi_1$: $2(2t)-(-t)+2(2t)=9$, so $9t=9$, $t=1$. Hence
\[
\boxed{H=(2,-1,2)}.
\]
\textbf{Check:} $|OH|=\sqrt{4+1+4}=3$ \checkmark, equal to the distance found above. \hfill \emph{[M1 distance formula, A1 value, M1 parametric perpendicular, A1 foot]}

\item \textbf{Plane of the pencil through $(0,1,-1)$.} Pencil:
\[
(2x-y+2z-9)+\lambda(x+2y-z-4)=0.
\]
At $(0,1,-1)$:
\[
(0-1-2-9)+\lambda(0+2+1-4)=0
\;\Longrightarrow\; -12-\lambda=0
\;\Longrightarrow\; \lambda=-12.
\]
Therefore
\[
(2x-y+2z-9)-12(x+2y-z-4)=0
\;\Longrightarrow\; -10x-25y+14z+39=0,
\]
\[
\boxed{10x+25y-14z=39}.
\]
\textbf{Check:} $10(0)+25(1)-14(-1)=25+14=39$ \checkmark. \hfill \emph{[M1 pencil form, A1 $\lambda=-12$, A1 simplified equation]}
\end{enumerate}
\textbf{Examiner's expectations:} method marks are awarded for correct \emph{processes} (cross product, pencil, distance formula) even if an arithmetic slip occurs; always verify the final plane against the given point, give the angle to the accuracy demanded, and state the range $\lambda\in\mathbb{R}$ when writing a vector equation of a line.
\end{corrige}

\newpage

\coursec{Summary sheet}

\begin{popfigure}
\begin{tikzpicture}[
  mindmap, grow cyclic,
  every node/.style={concept, text=white, font=\small},
  root concept/.append style={concept color=popink, font=\small\bfseries},
  level 1/.append style={level distance=3.4cm, sibling angle=60},
  level 1 concept/.append style={font=\scriptsize},
  concept color=popblue
]
\node[root concept]{Vectors: Planes in Space}
  child[concept color=popblue]{ node {Equations of a plane: $\vec{r}\cdot\vec{n}=d$; $ax+by+cz=d$} }
  child[concept color=popgreen]{ node {Plane through 3 points; parametric form $\vec{r}=\vec{a}+\lambda\vec{u}+\mu\vec{v}$} }
  child[concept color=poporange]{ node {Distance point--plane: $\dfrac{|ax_1+by_1+cz_1-d|}{\sqrt{a^2+b^2+c^2}}$} }
  child[concept color=poppurple]{ node {Angles: plane--plane $\cos\theta=\dfrac{|\vec{n}_1\cdot\vec{n}_2|}{|\vec{n}_1||\vec{n}_2|}$; line--plane $\sin\phi=\dfrac{|\vec{u}\cdot\vec{n}|}{|\vec{u}||\vec{n}|}$} }
  child[concept color=poppink]{ node {Intersection line--plane: substitute $\vec{r}=\vec{a}+\lambda\vec{u}$} }
  child[concept color=popgold]{ node {Two planes: line of intersection, direction $\vec{n}_1\times\vec{n}_2$} }
  child[concept color=popmuted]{ node {Pencils: $\pi_1+\lambda\pi_2=0$} };
\end{tikzpicture}
\legende{Mind map of the chapter: Planes in space.}
\end{popfigure}

\begin{bilan}
\textbf{Key definitions}
\begin{itemize}
\item A \cle{plane} is determined by: one point and a \cle{normal vector} $\vec{n}$; or one point and two non-parallel direction vectors; or three non-collinear points.
\item \cle{Vector equation}: $(\vec{r}-\vec{a})\cdot\vec{n}=0$, i.e. $\vec{r}\cdot\vec{n}=\vec{a}\cdot\vec{n}=d$.
\item \cle{Cartesian equation}: $ax+by+cz=d$, where $\vec{n}=\begin{pmatrix}a\\b\\c\end{pmatrix}$ is normal to the plane.
\item \cle{Parametric form}: $\vec{r}=\vec{a}+\lambda\vec{u}+\mu\vec{v}$, $\lambda,\mu\in\mathbb{R}$, with $\vec{u},\vec{v}$ non-parallel directions in the plane; equivalently $x=x_0+\lambda u_1+\mu v_1$, etc.
\item A \cle{pencil of planes}: the family $\pi_1+\lambda\pi_2=0$ of all planes through the line of intersection of $\pi_1$ and $\pi_2$.
\end{itemize}

\textbf{Formulas and theorems to know}
\begin{itemize}
\item Normal from two directions: $\vec{n}=\vec{u}\times\vec{v}$ (cross product).
\item Plane through $A$, $B$, $C$: $\vec{n}=\overrightarrow{AB}\times\overrightarrow{AC}$; check first that $A,B,C$ are \emph{not collinear}.
\item \cle{Perpendicular distance} from $P(x_1,y_1,z_1)$ to $ax+by+cz=d$:
\[ \delta=\frac{|ax_1+by_1+cz_1-d|}{\sqrt{a^2+b^2+c^2}}. \]
\item \cle{Angle between two planes} = angle between their normals:
\[ \cos\theta=\frac{|\vec{n}_1\cdot\vec{n}_2|}{|\vec{n}_1|\,|\vec{n}_2|}. \]
Planes are parallel iff $\vec{n}_1=k\vec{n}_2$; perpendicular iff $\vec{n}_1\cdot\vec{n}_2=0$.
\item \cle{Angle between a line} (direction $\vec{u}$) \cle{and a plane} (normal $\vec{n}$):
\[ \sin\phi=\frac{|\vec{u}\cdot\vec{n}|}{|\vec{u}|\,|\vec{n}|}, \qquad \phi=90^\circ-\theta. \]
\item Line meets plane: substitute $\vec{r}=\vec{a}+\lambda\vec{u}$ into $ax+by+cz=d$, solve for $\lambda$. If $\vec{u}\cdot\vec{n}=0$: line is parallel to the plane (no intersection, or the line lies in the plane).
\item Line of intersection of two planes: direction $\vec{u}=\vec{n}_1\times\vec{n}_2$; find one point by setting e.g. $z=0$ and solving the two equations.
\item Plane through the line of intersection of $\pi_1: a_1x+b_1y+c_1z=d_1$ and $\pi_2: a_2x+b_2y+c_2z=d_2$:
\[ (a_1x+b_1y+c_1z-d_1)+\lambda(a_2x+b_2y+c_2z-d_2)=0, \]
then determine $\lambda$ from an extra condition (a given point, perpendicularity, etc.).
\end{itemize}

\textbf{Methods}
\begin{itemize}
\item To find a plane: identify a point on it and a normal vector (often via a cross product), then expand $\vec{r}\cdot\vec{n}=\vec{a}\cdot\vec{n}$.
\item To convert parametric to Cartesian: eliminate $\lambda$ and $\mu$, or compute $\vec{n}=\vec{u}\times\vec{v}$.
\item For distance problems, write the plane in the form $ax+by+cz-d=0$ \emph{before} applying the formula.
\item For pencil problems, substitute the given point into $\pi_1+\lambda\pi_2=0$ to get a linear equation in $\lambda$.
\end{itemize}

\textbf{Common mistakes}
\begin{itemize}
\item Forgetting the absolute value in the distance formula, or forgetting to move $d$ to the left-hand side.
\item Using $\cos$ instead of $\sin$ for the line--plane angle: the angle uses the \emph{complement} of the angle with the normal.
\item Applying the ``plane through three points'' method when the points are collinear: infinitely many planes then contain them.
\item Sign errors in the cross product: compute each component carefully, e.g. $(u_2v_3-u_3v_2,\ u_3v_1-u_1v_3,\ u_1v_2-u_2v_1)$.
\item Concluding ``no intersection'' when $\vec{u}\cdot\vec{n}=0$ without checking whether the line actually lies in the plane.
\end{itemize}

\textbf{GCE tips}
\begin{itemize}
\item State each formula before substituting: method marks are awarded even if arithmetic slips.
\item Keep vectors in column form and show the dot or cross product explicitly.
\item Give angles in degrees to 1 d.p. unless told otherwise, and distances to 3 s.f.
\item Always verify: does your point satisfy the plane equation? Does your direction vector satisfy $\vec{u}\cdot\vec{n}=0$ for a line claimed to lie in a plane?
\item In pencil questions, remember $\pi_1+\lambda\pi_2=0$ never includes $\pi_2$ itself; check separately if $\pi_2$ is the answer.
\end{itemize}
\end{bilan}

\findecours

\end{document}
